%---------------------------------------------------------------------------------------------------------------------------- 
\graphicspath{{Parti/P03-Meccanica_SR/C04-Vincoli/Figure/}}
\chapter{Prestazioni cinematiche e statiche dei vincoli}\label{cap:vincoli_CR}
%---------------------------------------------------------------------------------------------------------------------------- 


\begin{sommario}
	Si distingue preliminarmente il \emph{vincolo di
		rigidità} --- proprietà interna del corpo che impone uguaglianza di spostamento e rotazione in ogni zona di connessione,
	con le relative nozioni di distorsione e svincolo ---	dai \emph{vincoli di collegamento} che costituiscono	l'oggetto principale del capitolo. Si definisce il	vincolo posizionale, se ne introduce la molteplicità	e si distingue fra vincoli esterni e interni. Per
	ciascuno dei cinque vincoli elementari --- carrello,	bipendolo, cerniera, glifo e incastro --- si ricavano	l'equazione cinematica, il luogo del centro di	rotazione assoluto o relativo e la reazione vincolare,	ottenuta applicando il principio dei lavori virtuali.	Si mostra come i vincoli composti si ottengano per	giustapposizione dei vincoli semplici e come l'incastro	interno sia la forma di collegamento equivalente al	vincolo di rigidità. Si generalizza quindi la	trattazione al caso del cedimento vincolare e si	estende ai vincoli che collegano più di due corpi.	Il capitolo si chiude con una tavola riepilogativa	e una serie di esercizi proposti.
\end{sommario}


%----------------------------------------------------------------------------------------------------------------------------
\section{Il vincolo di rigidità e i vincoli di collegamento}
\label{sec:vincoli_intro}
%----------------------------------------------------------------------------------------------------------------------------

I vincoli che agiscono su un sistema di corpi rigidi
sono di due nature fondamentalmente diverse. Si
introduce dapprima il \emph{vincolo di rigidità},
proprietà interna del corpo, e si passa poi ai
\emph{vincoli di collegamento}, che costituiscono
il catalogo principale del capitolo:

\begin{itemize}[nosep]
	\item il \textbf{vincolo di rigidità}\index{vincolo!di rigidità|textbf}: ogni corpo
	rigido esteso soddisfa per sua natura una
	condizione di continuità --- le sue parti sono
	connesse. La rigidità impone inoltre che questa
	connessione non ammetta deformazioni relative:
	nessuno scorrimento, nessuna rotazione relativa
	tra le parti. Non è un dispositivo aggiunto
	dall'esterno ma una condizione che il corpo
	possiede per sua natura. La sua violazione
	controllata genera le \emph{distorsioni}; il
	suo rilascio parziale genera gli
	\emph{svincoli}. La trattazione quantitativa
	è sviluppata nei capitoli successivi;
	\item i \textbf{vincoli di collegamento}:
	relazioni cinematiche imposte fra due entità
	distinte --- un corpo e il suolo (vincoli
	esterni) oppure due corpi distinti
	$\mathcal{C}_i$ e $\mathcal{C}_j$ (vincoli
	interni), definiti formalmente nel
	\PARref{sec:definizione}.
\end{itemize}



\subsection{Prestazione cinematica}

Il vincolo di rigidità in una zona di connessione impone che le due parti del corpo abbiano, in quella zona, lo stesso spostamento e la stessa rotazione. In un problema piano	ciò si traduce in tre condizioni scalari: le due componenti	dello spostamento devono essere uguali sui due lati della zona di connessione, e altrettanto la rotazione. Non vi è alcun grado	di libertà relativo: la zona di connessione non può aprirsi, scorrere né ruotare.

\subsection{Prestazione statica}

Le reazioni interne trasmesse attraverso la sezione sono
distribuite sulla superficie di connessione. Ridotte a un
polo arbitrario $O$, esse equivalgono a una
\emph{forza risultante} --- due componenti scalari nel
piano --- e un \emph{momento risultante} --- una componente
scalare. In totale tre incognite, duali alle tre condizioni
cinematiche. Il polo $O$ è arbitrario: cambiandolo varia
il momento risultante ma non la forza né il numero delle
incognite.

\begin{oss}
	La dualità fra prestazione cinematica e prestazione
	statica è perfetta: a ciascuna delle tre condizioni
	cinematiche corrisponde una e una sola incognita
	statica. Questa struttura --- tre condizioni, tre
	reazioni --- sarà ritrovata nel catalogo dei vincoli
	di collegamento con il nome di \emph{incastro}.
	Il collegamento fra le due forme è discusso
	nell'osservazione del \PARref{sec:incastro}.
\end{oss}

\subsection{Distorsioni e svincoli}

La rigidità è la condizione di default: in assenza di
indicazioni contrarie ogni zona di connessione è rigida.
È possibile introdurre due tipi di eccezione:

\begin{itemize}[nosep]
	\item le \textbf{distorsioni}\index{distorsione|textbf}\index{distorsione!di Volterra}\index{distorsione|seealso{cedimento}} (o distorsioni di
	Volterra): violazioni \emph{controllate} della
	rigidità, in cui una o più delle tre condizioni
	cinematiche sono sostituite da valori non nulli
	assegnati. Modellano discontinuità cinematiche
	prescritte --- rotazioni relative imposte, variazioni
	termiche differenziali, disassamenti geometrici;
	\item gli \textbf{svincoli}\index{svincolo|textbf}\index{svincolo|seealso{vincolo}}: rilascio di una o più
	delle tre condizioni di rigidità, che vengono
	semplicemente rimosse. Uno svincolo rotazionale
	lascia libera la rotazione relativa fra i due tratti,
	mantenendo le due condizioni di traslazione; uno
	svincolo traslazionale lascia libera una componente
	di spostamento relativo. Ogni svincolo riduce di uno
	il numero di condizioni cinematiche attive e,
	dualmente, di uno il numero di reazioni interne
	trasmesse  attraverso la zona di connessione.
\end{itemize}

\begin{oss}
	I vincoli di collegamento interni possono essere
	interpretati come vincoli di rigidità parzialmente
	rilasciati tramite svincoli. Ad esempio, un vincolo
	che mantiene le due condizioni di traslazione ---
	uguaglianza delle componenti orizzontale e verticale
	dello spostamento relativo --- ma rilascia la
	condizione di rotazione, lasciando libera la rotazione
	relativa fra i due tratti, corrisponde a un vincolo
	di rigidità con svincolo rotazionale. Questa
	prospettiva --- rigidità come punto di partenza,
	svincolo come rilascio --- è del tutto equivalente
	a quella dei vincoli di collegamento sviluppata nelle
	sezioni seguenti, ma offre un quadro concettuale
	più unitario, particolarmente utile quando si
	introduce la deformabilità nei volumi successivi.
	Il lettore riconoscerà in questo esempio la cerniera
	interna, che sarà definita formalmente nel
	\PARref{sec:cerniera}.
\end{oss}

%---------------------------------------------------------------------------------------------------------------------------- 
\section{Il vincolo posizionale: definizione e classificazione}	\label{sec:definizione}
%---------------------------------------------------------------------------------------------------------------------------- 

\begin{definizione}[Vincolo posizionale]\label{def:vincolo}
	Un \emph{vincolo posizionale}\index{vincolo|textbf} è un dispositivo meccanico che, una
	volta applicato a uno o più corpi rigidi, limita le possibilità di
	spostamento ammissibili per i punti del corpo (o dei corpi) cui è
	collegato, imponendo una o più relazioni cinematiche lineari fra le
	componenti di spostamento.
\end{definizione}

La parola \textit{posizionale} indica che le condizioni imposte dal
vincolo riguardano le \emph{coordinate di posizione} dei punti del
corpo (o, in regime infinitesimo, i loro spostamenti), e non le
velocità: il vincolo stabilisce \emph{dove} il corpo può trovarsi,
non \emph{come} si sta muovendo. Vincoli che impongono invece
condizioni sulle velocità sono detti \emph{cinematici}\index{vincolo!cinematico}: alcuni
sono equivalenti a vincoli posizionali --- quando le condizioni
sulle velocità sono integrabili, come accade per il rotolamento
puro di una ruota su una guida rettilinea, dove $\dot{x} = R\,\dot{\varphi}$
si integra in $x = R\varphi + \text{cost.}$ --- altri non lo
sono e si dicono \emph{anolonomi}\index{vincolo!anolonomo}: il caso classico è quello del
pattino, in cui la velocità del centro deve restare allineata con
il pattino senza che ciò restringa le posizioni accessibili.
Vincoli cinematici e anolonomi non rientrano nello scopo di
questo capitolo.

\begin{definizione}[Molteplicità di un vincolo]\label{def:molteplicita}
	Si chiama \emph{molteplicità}\index{vincolo!molteplicita@molteplicità|textbf}\index{molteplicita@molteplicità|see{vincolo}} (o \emph{grado di vincolo}\index{vincolo!grado di|textbf}\index{grado di vincolo|see{vincolo}}) di un
	vincolo posizionale il numero di equazioni cinematiche lineari
	indipendenti che esso impone sulle grandezze cinematiche del corpo
	cui è applicato. Un vincolo di molteplicità $m$ riduce di $m$ il
	numero di gradi di libertà del sistema.
\end{definizione}

Nel seguito ci si restringe al \emph{caso piano}, in cui un
corpo rigido libero possiede tre gradi di libertà: due
componenti di traslazione e una rotazione. I vincoli che
studieremo hanno molteplicità $m \in \{1,2,3\}$ e lasciano
rispettivamente due, uno o nessun grado di libertà. La
trattazione tridimensionale, in cui il corpo rigido libero
possiede sei gradi di libertà, segue gli stessi principi
con le opportune estensioni.


\begin{oss}
	La linearità delle relazioni cinematiche richiesta nella
	\DEFref{def:vincolo} è conseguenza dell'ipotesi di
	spostamenti \emph{infinitesimi} adottata nel
	\CAPref{cap:cinematica_infinitesima}: nell'ambito della
	cinematica infinitesima, le relazioni geometriche si
	linearizzano e le equazioni di vincolo assumono la forma
	lineare utilizzata nel seguito. Per spostamenti finiti
	le relazioni diventerebbero non lineari e la nozione di
	molteplicità andrebbe reinterpretata.
\end{oss}


% ───────────────────────────────────────────────────────────────────────────
\subsection{Vincoli esterni e vincoli interni}
\label{sec:esterno_interno}
% ───────────────────────────────────────────────────────────────────────────

\begin{definizione}[Vincolo esterno]\label{def:esterno}
	Un vincolo si dice \emph{esterno}\index{vincolo!esterno|textbf} quando collega il corpo rigido
	$\mathcal{C}$ a un riferimento fisso (il suolo). Le equazioni che
	esso impone coinvolgono le componenti di \emph{spostamento assoluto}
	dei punti del corpo.
\end{definizione}


\begin{definizione}[Vincolo interno di collegamento]\label{def:interno}
	Un vincolo si dice \emph{interno}\index{vincolo!interno|textbf}\index{collegamento, vincolo di|see{vincolo (interno)}} --- o equivalentemente
	\emph{di collegamento} --- quando connette due corpi rigidi
	distinti $\mathcal{C}_i$ e $\mathcal{C}_j$. Il termine
	\emph{collegamento} sottolinea che si tratta di una relazione
	fra due entità separate, a differenza del vincolo di rigidità
	--- introdotto nel \PARref{sec:vincoli_intro} --- che è una
	proprietà interna del corpo stesso. Le equazioni che il vincolo
	interno impone coinvolgono lo \emph{spostamento relativo} e la
	\emph{rotazione relativa} dei due corpi nella zona di
	collegamento.
\end{definizione}

A seconda della geometria del collegamento si distinguono i
vincoli interni \emph{puntiformi}, in cui i due punti di
applicazione coincidono geometricamente in un unico punto $A$,
e i vincoli interni a \emph{estensione finita}, in cui sono
punti distinti dei due corpi: per i primi le equazioni di
vincolo si scrivono naturalmente attraverso il \emph{salto}
$\llbracket \cdot \rrbracket(A)$ delle grandezze cinematiche
attraverso $A$; per i secondi, attraverso il loro
\emph{incremento} $\Delta$ fra i due punti di applicazione.
La distinzione fra le due notazioni è quella stabilita nel
\PARref{sec:salto_incremento}; nelle sezioni seguenti la si
esemplifica caso per caso.

Il trattamento matematico dei vincoli interni è formalmente
identico a quello dei vincoli esterni: basta sostituire gli
spostamenti assoluti con gli spostamenti relativi e la
rotazione assoluta con la rotazione relativa. In modo analogo,
il \CRa\ individuato da un vincolo esterno diventa, nel caso
interno, il \CRr\ fra i due corpi collegati. Per questa
ragione, nelle sezioni seguenti si svilupperà il trattamento
per il caso esterno e si riporterà poi il risultato per il
caso interno in forma compatta.



%---------------------------------------------------------------------------------------------------------------------------- 
\section{Richiamo: il campo di spostamento e il centro di rotazione}	\label{sec:richiamo}
%---------------------------------------------------------------------------------------------------------------------------- 

I risultati di questo capitolo poggiano sulle formule ricavate
nel \CAPref{cap:cinematica_infinitesima}. Li ricordiamo
brevemente.

Nel caso piano, la formula fondamentale della cinematica infinitesima
del corpo rigido lega lo spostamento di un punto generico $P$ a quello
di un polo $O$ e alla rotazione $\vartheta$:
\begin{equation}\label{eq:FFCI_piano_vett}
	\ub(P) = \ub(O) + \bm{\upvartheta} \times \overrightarrow{OP},
\end{equation}
dove $\bm{\upvartheta} = \vartheta\,\ab_z$ è il vettore di rotazione
diretto lungo l'asse uscente dal piano. Ponendo il polo $O$
nell'origine del riferimento, in componenti cartesiane:
\begin{equation}\label{eq:FFCI_piano}
	\begin{Bmatrix} u(P) \\ v(P) \end{Bmatrix}
	=
	\begin{Bmatrix} u(O) \\ v(O) \end{Bmatrix}
	+
	\begin{bmatrix} 0 & -\vartheta \\ \vartheta & 0 \end{bmatrix}
	\begin{Bmatrix} x(P) \\ y(P) \end{Bmatrix}.
\end{equation}
Lo spostamento rigido piano è caratterizzato da tre parametri
indipendenti: le due componenti della traslazione $u(O)$, $v(O)$ e
la rotazione $\vartheta$.

Per $\vartheta \neq 0$ esiste un unico punto $C$ con spostamento
nullo: il \CRa. Il campo si riscrive come
\begin{equation}\label{eq:centro_campo}
	\ub(P) = \bm{\upvartheta} \times \overrightarrow{CP},
\end{equation}
e le coordinate di $C$, con $O$ nell'origine, si ricavano
imponendo $\ub(C) = \zerovb$:
\begin{equation}\label{eq:coord_C}
	x(C) = -\frac{v(O)}{\vartheta}, \qquad
	y(C) = \frac{u(O)}{\vartheta}.
\end{equation}

Per $\vartheta = 0$ il moto è una traslazione pura e il centro
si sposta all'infinito nella direzione perpendicolare alla
traslazione.

\medskip

Questo strumento geometrico --- la posizione del centro di
rotazione --- è il principale indicatore cinematico\index{centro di rotazione!luogo geometrico} che
utilizzeremo per caratterizzare le prestazioni di ciascun
vincolo: ogni equazione di vincolo può essere letta come
un'informazione sul luogo geometrico in cui il centro di
rotazione può trovarsi. Dualmente, la reazione vincolare\index{reazione vincolare|textbf}
è la grandezza statica che non compie lavoro negli
spostamenti ammessi dal vincolo.


\begin{oss}
	Per il corpo rigido, la cinematica è interamente
	descritta da tre parametri globali: le due componenti
	di traslazione del polo $O$ e la rotazione $\vartheta$.
	Quest'ultima è la stessa per tutti i punti del corpo,
	mentre la traslazione dipende dalla scelta del polo.
	Le equazioni di vincolo coinvolgono quindi grandezze
	globali --- lo spostamento di un punto specifico $A$,
	espresso tramite la formula fondamentale in funzione
	di $\ub(O)$ e $\vartheta$, e la rotazione $\vartheta$
	del corpo --- e il punto di applicazione $A$ individua
	il luogo geometrico in cui il vincolo è fisicamente
	connesso al corpo, senza che ciò limiti la generalità
	della trattazione.
	
	Per i corpi deformabili, come le travi che si 
	incontreranno nelle lezioni successive, la cinematica
	è descritta da un campo di spostamenti $\ub(s)$ che
	varia punto per punto. In quel contesto il punto di
	applicazione $A$ acquista un ruolo ancora più centrale:
	un bipendolo applicato in $A$ impone $\vartheta(A) = 0$
	solo in quel punto, non sull'intera struttura. La
	trattazione sviluppata in questo capitolo per il corpo
	rigido costituisce quindi il caso particolare --- e il
	fondamento concettuale --- della teoria generale dei
	vincoli per corpi deformabili.
\end{oss}

\begin{oss}
	Le reazioni vincolari si dividono in due categorie di
	natura diversa:
	\begin{enumerate}[label=(\roman*)]
		\item le \emph{forze reattive}\index{reazione vincolare!forza reattiva} $\rb(A)$\nomenclature[L]{$\rb(A)$}{forza reattiva (reazione vincolare nel punto $A$)}, che sono
		applicate in un punto specifico $A$ del corpo; il
		loro punto di applicazione è essenziale per il
		bilancio dei momenti e non può essere modificato
		arbitrariamente;
		\item le \emph{coppie reattive}\index{reazione vincolare!coppia reattiva} $\mu$, che sono
		vettori liberi: il loro effetto sul corpo rigido
		è indipendente dal punto in cui vengono applicate.
	\end{enumerate}
	Per uniformità notazionale, anche le coppie reattive
	si scrivono nella forma $\mu(A)$, in cui il punto $A$
	non rappresenta un punto di applicazione --- privo di
	significato per un vettore libero --- ma serve come
	\emph{etichetta} che identifica il vincolo di
	provenienza, in analogia con le componenti scalari
	$X(A), Y(A), R(A)$ delle forze. Questo accorgimento
	è essenziale nella soluzione dei problemi statici, in
	cui un corpo può ricevere coppie reattive da più
	vincoli distinti e ciascuna incognita scalare deve
	essere identificabile in modo univoco. Lo stesso vale
	per le coppie scambiate nei vincoli interni puntiformi,
	dove $A$ identifica il giunto.
	
	Nell'equilibrio del corpo rigido inteso come corpo unico,
	la distinzione (i)/(ii) fra forze applicate in un punto e
	coppie come vettori liberi è automaticamente soddisfatta.
	Diventa invece cruciale quando si considerano le forze
	interne trasmesse attraverso una sezione --- nella trave
	rigida e, a maggior ragione, nei corpi deformabili ---
	dove il punto di applicazione di forze e coppie influenza
	la risposta interna della struttura.
\end{oss}


%---------------------------------------------------------------------------------------------------------------------------- 
\section{Vincoli semplici: molteplicità uno}	\label{sec:semplici}
%---------------------------------------------------------------------------------------------------------------------------- 

I vincoli di molteplicità uno\index{vincolo!semplice} rimuovono un solo grado di libertà,
lasciando il corpo con due gradi di libertà residui. Le due tipologie
fondamentali sono il \emph{carrello} (o biella equivalente) e il
\emph{bipendolo}.

Per ciascun vincolo del catalogo si seguirà la medesima procedura
espositiva in tre passi: si scrive l'\emph{equazione cinematica}
imposta dal vincolo; si determina il \emph{luogo del centro di
	rotazione} compatibile con essa; si ricava la \emph{reazione
	vincolare} applicando il \emph{principio dei lavori virtuali}
(\CAPref{cap:dualismo_CR}), in base al quale la reazione è la
grandezza statica che compie lavoro nullo su tutti gli spostamenti
virtuali compatibili con il vincolo, in dualità con le componenti
cinematiche bloccate.


\subsection{Il carrello}\label{sec:carrello}

\paragraph{Realizzazione meccanica e parametri caratterizzanti.}
Il carrello\index{carrello|textbf}\index{carrello|seealso{vincolo}} è schematizzato dal simbolo di figura~\ref{fig:carrello}:
un dispositivo che impedisce la traslazione del punto di applicazione
$A$ in una direzione assegnata, lasciando libere la traslazione
nella direzione perpendicolare e la rotazione. Il vincolo è
completamente caratterizzato da:
\begin{enumerate}[label=(\roman*)]
	\item il \emph{punto di applicazione} $A$;
	\item il \emph{versore dell'asse} $\ab$, ovvero la direzione
	lungo la quale lo spostamento è impedito.
\end{enumerate}
Una realizzazione meccanica concreta del carrello è la
\emph{biella}\index{biella|textbf}\index{biella|seealso{carrello}}: un'asta rigida incardinata
con una cerniera al corpo nel punto $A$ e con un'altra cerniera
al suolo. Per spostamenti finiti il punto $A$ percorre un arco
di circonferenza centrato nel cardine a terra; nell'ipotesi di
spostamenti infinitesimi, lo spostamento ammesso di $A$ è
perpendicolare all'asse della biella, e quest'ultima fornisce
quindi le stesse prestazioni cinematiche del carrello, a
condizione che il suo asse e il versore $\ab$ del carrello
siano collineari.


%%%%%%%%%%%%%%%%%%%%%%%%%%%%%%%%%%%%%%%%%%%
\begin{figure}[htbp]
	\centering		
	\begin{tikzpicture}[>=Stealth, line width=0.8pt]
		
		% ═══════════════════════════════════════════════════════════
		% PANNELLO (a) — Simbolo del carrello
		% ═══════════════════════════════════════════════════════════
		\begin{scope}[xshift=0cm]
			
			% ── Parametri del simbolo (convenzioni del libro) ──
			\def\xA{0.8}     % ascissa di A
			\def\yA{0.0}     % ordinata di A
			\def\tw{0.30}    % semi-larghezza triangolo
			\def\thR{0.40}   % altezza triangolo carrello
			\def\rn{0.10}    % raggio semicerchio (anellino a metà)
			\def\bw{0.18}    % semi-distanza ruote
			\def\gw{0.45}    % semi-larghezza barre di terra
			
			% ── Asse del carrello: retta verticale grigia continua ──
			\draw[gray!60, line width=0.6pt] 
			(\xA, \yA-\thR-0.5) -- (\xA, 3.85);
			
			% ── Etichetta del luogo del CR ──
			\node[gray, anchor=west, font=\footnotesize] 
			at (\xA+0.08, 2.50) {$C \in$ asse del carrello};
			
			
			% ── Retta-luogo del CR: due tratti, sopra il corpo e sotto il suolo ──
			\draw[dashed, gray!60, line width=0.6pt] 
			(\xA, 2.05) -- (\xA, 2.85);
			\draw[ gray!60, line width=0.6pt] 
			(\xA, \yA-\thR) -- (\xA, \yA-\thR-0.15);
			
			% ── Corpo di fantasia (linea, senza fill) ──
			\draw[line width=1pt] 
			plot[smooth cycle, tension=0.7] coordinates 
			{(0.8, 0.0) (1.6, 0.15) (2.2, 0.7) (2.0, 1.7) 
				(1.0, 2.0) (-0.1, 1.7) (-0.6, 1.0) (-0.3, 0.3)};
			\node at (1.5, 1.1) {$\mathcal{C}$};
			
			% ── Spostamento ammesso: freccia grigia perpendicolare ad a ──
			%\draw[->, gray, line width=0.9pt] 	(\xA+0.20, \yA) -- (\xA+1.60, \yA);
			\node[above=5pt] at (\xA+2.50, \yA+0.5) 	{\footnotesize $\ub(A) \cdot \ab = 0$};
			
			% ── Simbolo del carrello: apice del triangolo IN A ──
			
			% Triangolo (apice in A, base in basso)
			\draw (\xA, \yA) -- ++(-\tw, -\thR) -- ++(2*\tw, 0) -- cycle;
			
			% Prima barra (alla base del triangolo)
			\draw (\xA-\gw, \yA-\thR) -- (\xA+\gw, \yA-\thR);
			
			% Due ruote (cerchietti vuoti)
			\draw (\xA-\bw, \yA-\thR-0.13) circle (3pt);
			\draw (\xA+\bw, \yA-\thR-0.13) circle (3pt);
			
			% Seconda barra (suolo)
			\draw (\xA-\gw, \yA-\thR-0.26) -- (\xA+\gw, \yA-\thR-0.26);
			
			% Tratteggio del suolo
			\foreach \x in {-0.36,-0.18,0,0.18,0.36}
			\draw[thin] (\xA+\x, \yA-\thR-0.26) -- ++(-0.10, -0.13);
			
			% Semicerchio (anellino a metà) in A: rotazione consentita
			\draw[fill=white] 
			(\xA-\rn, \yA) arc(180:360:\rn) -- cycle;
			
			% Punto A (ridisegnato dopo il semicerchio così resta visibile)
			%\fill (\xA, \yA) circle (1.5pt);
			\node[above left=-1pt] at (\xA-0.02, \yA+0.05) {$A$};
			
			% ── Asse a sopra il corpo ──
			\draw[->, thick] (\xA, 2.10) -- (\xA, 2.85);
			\node[right=2pt] at (\xA, 2.85) {$\ab$};
			
			% ── Etichetta verticale "asse del carrello" ──
			\node[gray, rotate=90, anchor=south, font=\footnotesize] 
			at (\xA-0.12, 2.45) {asse del carrello};
			
			% ── Etichetta pannello ──
			\node[anchor=north] at (\xA+2, \yA-\thR) {$(a)$};
			
		\end{scope}
		
		
		% ═══════════════════════════════════════════════════════════
		% PANNELLO (b) — biella (realizzazione meccanica)
		% ═══════════════════════════════════════════════════════════
		\begin{scope}[xshift=5.5cm]
			
			% ── Parametri ──
			\def\xA{0.8}     % ascissa di A (e del perno al suolo)
			\def\yA{1.5}     % ordinata di A
			\def\yG{0.0}     % ordinata del perno al suolo
			\def\tw{0.30}    % semi-larghezza triangolo
			\def\thA{0.40}   % altezza triangolo cerniera al suolo
			\def\rn{0.10}    % raggio semicerchio in A
			\def\gw{0.50}    % semi-larghezza barra di terra
			
			% ── Asse della biella: retta verticale dashed continua ──
			%    (verrà parzialmente coperta da asta e triangolo)
			\draw[gray!60, line width=0.6pt] 
			(\xA, \yG-\thA-0.5) -- (\xA, 4.50);
			
			% ── Corpo di fantasia (alzato, senza fill) ──
			\draw[line width=1pt] 
			plot[smooth cycle, tension=0.7] coordinates 
			{(0.8, 1.5) (1.6, 1.65) (2.2, 2.2) (2.0, 3.2) 
				(1.0, 3.5) (-0.1, 3.2) (-0.6, 2.5) (-0.3, 1.8)};
			\node at (0.95, 2.6) {$\mathcal{C}$};
			
			% ── Cinematica finita: arco ──
			\draw[gray, line width=0.9pt] 
			(\xA, \yG) ++(60:1.5) 
			arc[start angle=60, end angle=120, radius=1.5];
			\node[gray, anchor=west, font=\footnotesize] 
			at (1.6, 0.85) {cinematica finita};
			
			% ── Cinematica infinitesima: tangente ──
			\draw[gray, dashed, line width=0.9pt] 
			(-0.7, \yA) -- (2.4, \yA);
			\node[gray, anchor=west, font=\footnotesize] 
			at (2.2, \yA+0.35) {cinematica infinitesima};
			
			% ── Asta della biella (copre il tratto centrale dell'asse) ──
			\draw[line width=0.8pt] (\xA, \yA) -- (\xA, \yG);
			
			% ── Cerniera al suolo: triangolo + barra + tratteggio ──
			\draw (\xA, \yG) -- ++(-\tw, -\thA) -- ++(2*\tw, 0) -- cycle;
			\draw (\xA-\gw, \yG-\thA) -- (\xA+\gw, \yG-\thA);
			\foreach \x in {-0.36,-0.18,0,0.18,0.36}
			\draw[thin] (\xA+\x, \yG-\thA) -- ++(-0.10, -0.13);
			
			% ── Anellini (in quest'ordine, dopo asta e triangolo) ──
			% Sommità della biella (in A): semicerchio bianco, concavità in basso
			\draw[fill=white] 
			(\xA-\rn, \yA) arc(180:360:\rn) -- cycle;
			% Cerniera al suolo: perno pieno
			\draw[fill=white] (\xA, \yG) circle (3.5pt);
			
			% ── Punto A (ridisegnato dopo il semicerchio) ──
			%\fill (\xA, \yA) circle (1.5pt);
			\node[above left=-1pt] at (\xA-0.02, \yA+0.05) {$A$};
			
			% ── Asse a sopra il corpo ──
			\draw[->, thick] (\xA, 3.60) -- (\xA, 4.30);
			\node[right=2pt] at (\xA, 4.30) {$\ab$};
			
			% ── Etichetta del luogo del CR ──
			\node[gray, anchor=west, font=\footnotesize] 
			at (\xA+0.08, 4.00) {$C \in$ asse della biella};
			
			% ── Etichetta pannello ──
			\node[anchor=north] at (\xA+2, \yG-\thA) {$(b)$};
			
		\end{scope}
		
	\end{tikzpicture}
	\caption{Vincolo semplice di carrello: simbolo convenzionale (a)
		e realizzazione meccanica tramite biella (b). L'equazione
		di vincolo $\ub(A)\cdot\ab = 0$ annulla la componente di
		spostamento di $A$ lungo l'asse $\ab$; le componenti libere
		--- traslazione lungo $\ab^\perp$ e rotazione --- non sono
		rappresentate. Per la biella, la cinematica finita è un
		arco di circonferenza centrato nel perno al suolo;
		nell'ipotesi di spostamenti infinitesimi essa coincide con
		la tangente in $A$, e le prestazioni cinematiche della
		biella si riducono a quelle del carrello. In entrambi i
		casi il \CRa\ giace sulla retta per $A$ in direzione $\ab$.}
	\label{fig:carrello}
\end{figure}

\paragraph{Caso esterno.}

\paragraph{Equazione cinematica.}
Il carrello impedisce la componente di spostamento del punto $A$
nella direzione $\ab$. L'equazione di vincolo è:
\begin{equation}\label{eq:carrello_cin}
	\ub (A) \cdot \ab = 0.
\end{equation}
Il punto $A$ può liberamente spostarsi nella direzione
perpendicolare, ma non lungo $\ab$.

\paragraph{Posizione del centro di rotazione assoluto.}
Sostituendo la~\eqref{eq:centro_campo} nella~\eqref{eq:carrello_cin}:
\begin{equation}
	\bigl(\bm{\upvartheta} \times \overrightarrow{CA}\bigr)
	\cdot \ab = 0.
\end{equation}
Applicando la proprietà ciclica del prodotto misto:
\begin{equation}
	\bm{\upvartheta} \cdot \bigl(\overrightarrow{CA} \times \ab\bigr)
	= 0.
\end{equation}
Poiché $\bm{\upvartheta} = \vartheta\,\ab_z$ e
$\overrightarrow{CA} \times \ab$ è parallelo ad $\ab_z$, la
condizione si riduce a richiedere che il prodotto vettoriale
$\overrightarrow{CA} \times \ab$ sia nullo, ovvero che
$\overrightarrow{CA}$ sia parallelo ad $\ab$.

\begin{proposizione}[Centro di rotazione --- carrello]
	\label{prop:CR_carrello}
	Il \CRa\ di un corpo vincolato da un carrello applicato in $A$
	con asse $\ab$ giace sulla retta passante per $A$ nella
	direzione di $\ab$:
	\begin{equation}\label{eq:CR_carrello}
		C = A + \lambda\,\ab, \qquad \lambda \in \mathbb{R}.
	\end{equation}
	La retta costituisce il \emph{luogo geometrico} dei centri
	compatibili con il vincolo; la posizione esatta lungo tale
	retta è determinata solo una volta assegnati ulteriori vincoli
	o lo spostamento del corpo.
\end{proposizione}


\paragraph{Reazione vincolare.}
Gli spostamenti compatibili con il carrello sono quelli per cui
$\ub(A)\cdot\ab = 0$: il punto $A$ si sposta liberamente in
direzione perpendicolare ad $\ab$. Per il principio dei lavori
virtuali, la reazione deve compiere lavoro nullo su ogni tale
spostamento: la forza reattiva è quindi parallela ad $\ab$ e la
coppia reattiva è nulla,
\begin{equation}\label{eq:reazione_carrello}
	\rb(A) = R(A)\,\ab,
	\qquad
	\mu(A) = 0,
	\qquad
	R(A) \in \mathbb{R}.
\end{equation}
Sono noti a priori il punto di applicazione $A$, la direzione
$\ab$ della forza reattiva e l'assenza di coppia reattiva;
l'unica incognita scalare è l'intensità con segno $R(A)$, il
cui segno determina se la reazione è concorde o discorde con
$\ab$, in accordo con la molteplicità uno. La
\FIGref{fig:carrello_stat} illustra la reazione per il simbolo
del carrello e per la sua realizzazione meccanica tramite
biella.




%%%%%%%%%%%%%%%%%%%%%%%%%%%%%%%%%%%%%%%%%%%
\begin{figure}[htbp]
	\centering        
	\begin{tikzpicture}[>=Stealth, line width=0.8pt]
		
		% ═══════════════════════════════════════════════════════════
		% PANNELLO (a) — Reazione del carrello
		% ═══════════════════════════════════════════════════════════
		\begin{scope}[xshift=0cm]
			
			% ── Parametri ──
			\def\xA{0.8}
			\def\yA{0.0}
			\def\tw{0.30}
			\def\thR{0.40}
			\def\rn{0.10}
			\def\bw{0.18}
			\def\gw{0.45}
			
			% ── Asse del carrello: retta verticale grigia continua ──
			\draw[gray!60, line width=0.6pt] 
			(\xA, \yA-\thR-0.5) -- (\xA, 3.85);
			
			% ── Corpo di fantasia ──
			\draw[line width=1pt] 
			plot[smooth cycle, tension=0.7] coordinates 
			{(0.8, 0.0) (1.6, 0.15) (2.2, 0.7) (2.0, 1.7) 
				(1.0, 2.0) (-0.1, 1.7) (-0.6, 1.0) (-0.3, 0.3)};
			\node at (1.5, 1.5) {$\mathcal{C}$};
			
			% ── Simbolo del carrello ──
			\draw (\xA, \yA) -- ++(-\tw, -\thR) -- ++(2*\tw, 0) -- cycle;
			\draw (\xA-\gw, \yA-\thR) -- (\xA+\gw, \yA-\thR);
			\draw (\xA-\bw, \yA-\thR-0.13) circle (3pt);
			\draw (\xA+\bw, \yA-\thR-0.13) circle (3pt);
			\draw (\xA-\gw, \yA-\thR-0.26) -- (\xA+\gw, \yA-\thR-0.26);
			\foreach \x in {-0.36,-0.18,0,0.18,0.36}
			\draw[thin] (\xA+\x, \yA-\thR-0.26) -- ++(-0.10, -0.13);
			
			% ── Semicerchio in A: rotazione consentita ──
			\draw[fill=white] 
			(\xA-\rn, \yA) arc(180:360:\rn) -- cycle;
			
			% ── REAZIONE: forza R(A) a, in rosso, da A verso l'alto ──
			\draw[->, line width=1.2pt] 
			(\xA, \yA+0.05) -- (\xA, \yA+1.05);
			\node[anchor=west, font=\small] 
			at (\xA, \yA+0.65) {$R(A)\,\ab$};
			
			% ── Etichetta A (a destra, fuori dal triangolo) ──
			\node[anchor=west, font=\small] at (\xA+0.18, \yA-0.08) {$A$};
			
			% ── Asse a sopra il corpo ──
			\draw[->, thick] (\xA, 2.10) -- (\xA, 2.85);
			\node[right=2pt] at (\xA, 2.85) {$\ab$};
			
			% ── Etichetta verticale "asse del carrello" ──
			\node[gray, rotate=90, anchor=south, font=\footnotesize] 
			at (\xA-0.12, 2.45) {asse del carrello};
			
			% ── Etichetta pannello ──
			\node[anchor=north] at (\xA+2, \yA-\thR) {$(a)$};
			
		\end{scope}
		
		
		% ═══════════════════════════════════════════════════════════
		% PANNELLO (b) — Reazione della biella
		% ═══════════════════════════════════════════════════════════
		\begin{scope}[xshift=5.5cm]
			
			% ── Parametri ──
			\def\xA{0.8}
			\def\yA{1.5}
			\def\yG{0.0}
			\def\tw{0.30}
			\def\thA{0.40}
			\def\rn{0.10}
			\def\gw{0.50}
			
			% ── Asse della biella ──
			\draw[gray!60, line width=0.6pt] 
			(\xA, \yG-\thA-0.5) -- (\xA, 4.50);
			
			% ── Corpo di fantasia ──
			\draw[line width=1pt] 
			plot[smooth cycle, tension=0.7] coordinates 
			{(0.8, 1.5) (1.6, 1.65) (2.2, 2.2) (2.0, 3.2) 
				(1.0, 3.5) (-0.1, 3.2) (-0.6, 2.5) (-0.3, 1.8)};
			\node at (1.5, 3.0) {$\mathcal{C}$};
			
			% ── Asta della biella ──
			\draw[line width=0.8pt] (\xA, \yA) -- (\xA, \yG);
			
			% ── Cerniera al suolo ──
			\draw (\xA, \yG) -- ++(-\tw, -\thA) -- ++(2*\tw, 0) -- cycle;
			\draw (\xA-\gw, \yG-\thA) -- (\xA+\gw, \yG-\thA);
			\foreach \x in {-0.36,-0.18,0,0.18,0.36}
			\draw[thin] (\xA+\x, \yG-\thA) -- ++(-0.10, -0.13);
			
			% ── Anellini: semicerchio in A, perno pieno al suolo ──
			\draw[fill=white] 
			(\xA-\rn, \yA) arc(180:360:\rn) -- cycle;
			\draw[fill=white] (\xA, \yG) circle (3.5pt);
			
			% ── REAZIONE: forza R(A) a, in rosso, da A verso l'alto ──
			\draw[->,  line width=1.2pt] 
			(\xA, \yA+0.05) -- (\xA, \yA+1.05);
			\node[anchor=west, font=\small] 
			at (\xA, \yA+0.65) {$R(A)\,\ab$};
			
			% ── Etichetta A (a sinistra dell'asta) ──
			\node[anchor=east, font=\small] at (\xA, \yA-0.24) {$A$};
			
			% ── Asse a sopra il corpo ──
			\draw[->, thick] (\xA, 3.60) -- (\xA, 4.30);
			\node[right=2pt] at (\xA, 4.30) {$\ab$};
			
			% ── Etichetta verticale "asse della biella" ──
			\node[gray, rotate=90, anchor=south, font=\footnotesize] 
			at (\xA-0.12, 3.40) {asse della biella};
			
			% ── Etichetta pannello ──
			\node[anchor=north] at (\xA+2, \yG-\thA) {$(b)$};
			
		\end{scope}
		
	\end{tikzpicture}
	\caption{Prestazione statica del carrello (a) e della biella
		(b): la reazione vincolare si riduce alla sola forza
		$R(A)\,\ab$ diretta lungo l'asse, in dualità con la
		componente di spostamento bloccata; le componenti reattive
		nulle (coppia e forza in $\ab^\perp$) non sono
		rappresentate, in accordo con la molteplicità uno. Nella
		biella, la forza è trasmessa al corpo $\mathcal{C}$
		dall'asta lungo il proprio asse.}
	\label{fig:carrello_stat}
\end{figure}
%%%%%%%%%%%%%%%%%%%%%%%%%%%%%%%%%%%%%%%%%%%


\paragraph{Caso interno.}
Quando il vincolo collega due corpi $\mathcal{C}_i$ e
$\mathcal{C}_j$, la condizione~\eqref{eq:carrello_cin} si
specializza in due forme distinte, secondo che il vincolo
sia puntiforme o esteso.


Nel \emph{carrello interno}\index{carrello!interno} i punti di applicazione $A_i$
e $A_j$ coincidono geometricamente, $A_i = A_j = A$, e la
condizione di uguaglianza degli spostamenti nella direzione
$\ab$ si scrive come salto attraverso $A$:
\begin{equation}\label{eq:carrello_int_punt}
	\llbracket \ub \rrbracket(A) \cdot \ab = 0,
	\qquad
	\llbracket \ub \rrbracket(A) \coloneqq \ub_j(A) -
	\ub_i(A),
\end{equation}
in accordo con la convenzione del
\PARref{sec:salto_incremento}. Nella \emph{biella di
	lunghezza finita}, invece, $A_i$ e $A_j$ sono punti
distinti dei due corpi e la condizione si esprime tramite
l'incremento dello spostamento fra di essi:
\begin{equation}\label{eq:carrello_int_est}
	\Delta\ub \cdot \ab = 0,
	\qquad
	\Delta\ub \coloneqq \ub(A_j) - \ub(A_i),
\end{equation}
con $\ab \coloneqq \overrightarrow{A_i A_j}/
|\overrightarrow{A_i A_j}|$ determinato dalla geometria
dell'asta. Il carrello si recupera come limite della
biella finita per $|\overrightarrow{A_i A_j}| \to 0$: i
due punti di applicazione collassano in un unico punto
$A$ e l'incremento $\Delta$ degenera nel salto
$\llbracket \cdot \rrbracket(A)$.


%%%%%%%%%%%%%%%%%%%%%%%%%%%%%%%%%%%%%%%%%%%
\begin{figure}[htbp]
	\centering
	\resizebox{\textwidth}{!}{%
		\begin{tikzpicture}[>=Stealth, line width=0.8pt]
			
			% ═══════════════════════════════════════════════════════════
			% PANNELLO (a) — Prestazione cinematica della biella interno
			% ═══════════════════════════════════════════════════════════
			\begin{scope}[xshift=0cm]
				
				% ── Parametri (carrello ruotato di 90° CCW) ──
				\def\xAi{-0.30}     % ascissa di A_i (bordo destro di C_i)
				\def\xAj{ 1.20}     % ascissa di A_j (apice triangolo)
				\def\yAx{ 0.00}     % asse a coincide con y=0
				\def\tw{0.30}       % semi-altezza triangolo (era larghezza)
				\def\thR{0.40}      % "altezza" triangolo = base - apice
				\def\rn{0.10}       % raggio semicerchio
				\def\bw{0.18}       % semi-distanza ruote (verticale)
				\def\gw{0.45}       % semi-altezza barre laterali
				
				%% ── Asse a (retta grigia continua) ──
				\draw[gray!60, line width=0.6pt] 
				(-2.5, \yAx) -- (3.40, \yAx);
				
				%% ── Corpo C_i a sinistra ──
				\draw[line width=1pt] 
				plot[smooth cycle, tension=0.7] coordinates 
				{(-0.30,  0.00) (-0.50,  0.80) (-1.40,  1.10) 
					(-2.40,  0.95) (-2.80,  0.20) (-2.60, -0.70) 
					(-1.70, -1.05) (-0.70, -0.80) (-0.45, -0.45)};
				\node at (-1.60, 0.30) {$\mathcal{C}_i$};
				
				%% ── Corpo C_j a destra ──
				\draw[line width=1pt] 
				plot[smooth cycle, tension=0.7] coordinates 
				{(1.20,  0.00) (1.40,  0.80) (2.20,  1.10) 
					(3.20,  0.95) (3.60,  0.20) (3.40, -0.70) 
					(2.50, -1.05) (1.60, -0.85) (1.35, -0.45)};
				\node at (2.40, 0.30) {$\mathcal{C}_j$};
				
				
				%% ── biella: asta orizzontale + semicerchi terminali ──
				\draw[line width=0.8pt] (\xAi, \yAx) -- (\xAj, \yAx);
				
				%% Semicerchio in A_i (concavità a sinistra, verso C_i)
				\draw[fill=white] 
				(\xAi, \yAx-\rn) arc(-90:90:\rn) -- cycle;
				
				%% Semicerchio in A_j (concavità a destra, verso C_j)
				\draw[fill=white] 
				(\xAj, \yAx+\rn) arc(90:270:\rn) -- cycle;
				
				%% ── Semicerchio in A_j (apice del triangolo) ──
				\draw[fill=white] 
				(\xAj, \yAx+\rn) arc(90:270:\rn) -- cycle;
				
				%% ── Etichette A_i e A_j ──
				\node[anchor=south east, font=\small] 
				at (\xAi-0.05, \yAx+0.05) {$A_i$};
				\node[anchor=south west, font=\small] 
				at (\xAj+0.10, \yAx+0.05) {$A_j$};
				
				
				%% ── Versore a, allineato con l'asse della biella ──
				\draw[->, line width=1pt] 
				(\xAj+1.40, \yAx) -- (\xAj+2.20, \yAx);
				\node[anchor=north, font=\small] 
				at (\xAj+1.70, \yAx-0.05) {$\ab$};
				
				%% ──────────────────────────────────────────────────
				%% VETTORI SPOSTAMENTO (linee spesse)
				%% ──────────────────────────────────────────────────
				
				% u(A_i): da A_i, direzione 70°, modulo |u_i|
				\def\uimx{ 0.41}    % cos(70°) * |u_i|
				\def\uimy{ 1.13}    % sin(70°) * |u_i|, scelti per ottenere d=0.41
				\draw[->, line width=1.4pt] 
				(\xAi, \yAx) -- ++(\uimx, \uimy);
				\node[anchor=south, font=\small] 
				at (\xAi+\uimx, \yAx+\uimy-0.05) {$\ub(A_i)$};
				
				% Linea tratteggiata di proiezione di u(A_i) sull'asse
				\draw[dashed, gray!70, line width=0.5pt] 
				(\xAi+\uimx, \yAx+\uimy) -- (\xAi+\uimx, \yAx);
				
				% u(A_j): da A_j, direzione -50°, modulo |u_j|
				\def\ujmx{ 0.41}    % cos(-50°) * |u_j| (stesso d!)
				\def\ujmy{-0.49}    % sin(-50°) * |u_j|
				\draw[->, line width=1.4pt] 
				(\xAj, \yAx) -- ++(\ujmx, \ujmy);
				\node[anchor=north west, font=\small] 
				at (\xAj+\ujmx-0.05, \yAx+\ujmy+0.2) {$\ub(A_j)$};
				
				% Linea tratteggiata di proiezione di u(A_j) sull'asse
				\draw[dashed, gray!70, line width=0.5pt] 
				(\xAj+\ujmx, \yAx+\ujmy) -- (\xAj+\ujmx, \yAx);
				
				%% ── Etichetta del luogo del CR relativo (sopra C_j) ──
				\node[font=\footnotesize, anchor=south] 
				at (2.50, 1.20) 
				{$C_{ij} \in$ asse della biella};
				
				
				%% ──────────────────────────────────────────────────
				%% TACCHE DI UGUAGLIANZA (doppia barretta)
				%% sui due segmenti di proiezione, sull'asse
				%% ──────────────────────────────────────────────────
				% Proiezione di u(A_i): segmento da (xAi, 0) a (xAi+uimx, 0)
				% Punto medio: (xAi + uimx/2, 0)
				% Proiezione di u(A_j): segmento da (xAj, 0) a (xAj+ujmx, 0)
				% Punto medio: (xAj + ujmx/2, 0)
				
				% Tacche su proiezione i (due barrette verticali)
				\draw[line width=0.7pt] 
				(\xAi+\uimx/2-0.04, -0.08) -- (\xAi+\uimx/2-0.04, 0.08);
				\draw[line width=0.7pt] 
				(\xAi+\uimx/2+0.04, -0.08) -- (\xAi+\uimx/2+0.04, 0.08);
				
				% Tacche su proiezione j
				\draw[line width=0.7pt] 
				(\xAj+\ujmx/2-0.04, -0.08) -- (\xAj+\ujmx/2-0.04, 0.08);
				\draw[line width=0.7pt] 
				(\xAj+\ujmx/2+0.04, -0.08) -- (\xAj+\ujmx/2+0.04, 0.08);
				
				%% ── Equazione in alto al centro ──
				\node[font=\small] at (0.45, 2.05) 
				{$\ub(A_i)\cdot\ab \;=\; \ub(A_j)\cdot\ab$};
				
				%% ── Etichetta pannello ──
				\node[anchor=north, font=\small] at (0.45, -1.0) {$(a)$};
				
			\end{scope}
			
			% ═══════════════════════════════════════════════════════════
			% PANNELLO (b) — Prestazione statica della biella interno
			% ═══════════════════════════════════════════════════════════
			\begin{scope}[xshift=8cm]   % o yshift=-3.5cm, a tua scelta
				
				% ── Parametri (identici a (a)) ──
				\def\xAi{-0.50}
				\def\xAj{ 1.40}
				\def\yAx{ 0.00}
				\def\rn{0.10}
				
				%% ── Asse a (retta grigia continua) ──
				\draw[gray!60, line width=0.6pt] 
				(-3.80, \yAx) -- (3.65, \yAx);
				
				%% ── Corpo C_i (identico ad (a)) ──
				\draw[line width=1pt] 
				plot[smooth cycle, tension=0.7] coordinates 
				{(-0.50,  0.00) (-0.70,  0.80) (-1.60,  1.10) 
					(-2.60,  0.95) (-3.00,  0.20) (-2.80, -0.70) 
					(-1.90, -1.05) (-0.90, -0.80) (-0.65, -0.45)};
				\node at (-1.80, 0.30) {$\mathcal{C}_i$};
				
				%% ── Corpo C_j (identico ad (a)) ──
				\draw[line width=1pt] 
				plot[smooth cycle, tension=0.7] coordinates 
				{(1.40,  0.00) (1.60,  0.80) (2.40,  1.10) 
					(3.40,  0.95) (3.80,  0.20) (3.60, -0.70) 
					(2.70, -1.05) (1.80, -0.85) (1.55, -0.45)};
				\node at (2.60, 0.30) {$\mathcal{C}_j$};
				
				%% ── biella: asta + semicerchi terminali ──
				\draw[line width=0.8pt] (\xAi, \yAx) -- (\xAj, \yAx);
				\draw[fill=white] 
				(\xAi, \yAx-\rn) arc(-90:90:\rn) -- cycle;
				\draw[fill=white] 
				(\xAj, \yAx+\rn) arc(90:270:\rn) -- cycle;
				
				%% ── Etichette A_i e A_j ──
				\node[anchor=south east, font=\small] 
				at (\xAi-0.05, \yAx+0.05) {$A_i$};
				\node[anchor=south west, font=\small] 
				at (\xAj+0.10, \yAx+0.05) {$A_j$};
				
				%% ── Versore a, allineato con l'asse della biella ──
				\draw[->, line width=1pt] 
				(\xAj+1.50, \yAx) -- (\xAj+2.10, \yAx);
				\node[anchor=north, font=\small] 
				at (\xAj+1.70, \yAx-0.05) {$\ab$};
				
				%% ──────────────────────────────────────────────────
				%% REAZIONI (linee spesse, entranti nei corpi)
				%% ──────────────────────────────────────────────────
				
				%% Reazione in A_i: -R a, da A_i verso sinistra (entra in C_i)
				\draw[->, line width=1.4pt] 
				(\xAi, \yAx) -- ++(-1.00, 0);
				\node[anchor=north, font=\small] 
				at (\xAi-0.80, \yAx-0.06) {$-R\,\ab$};
				
				%% Reazione in A_j: +R a, da A_j verso destra (entra in C_j)
				\draw[->, line width=1.4pt] 
				(\xAj, \yAx) -- ++(1.00, 0);
				\node[anchor=north, font=\small] 
				at (\xAj+0.50, \yAx-0.06) {$R\,\ab$};
				
				%% ── Etichetta pannello ──
				\node[anchor=north, font=\small] at (0.45, -1.0) {$(b)$};
				
			\end{scope}
			
		\end{tikzpicture}
	}
	\caption[Biella interna fra due corpi: prestazioni cinematica e statica.]{%
		Biella interna fra i corpi $\mathcal{C}_i$ e $\mathcal{C}_j$,
		con punti di applicazione $A_i$ e $A_j$ e versore d'asse
		$\ab = \overrightarrow{A_iA_j}/|\overrightarrow{A_iA_j}|$.
		Prestazione cinematica (a): l'equazione di vincolo
		$\ub(A_i)\cdot\ab = \ub(A_j)\cdot\ab$ impone l'uguaglianza
		delle componenti degli spostamenti di $A_i$ e $A_j$ lungo
		l'asse $\ab$; la componente trasversale e la rotazione
		relativa restano libere e non sono rappresentate.
		Prestazione statica (b): le reazioni scambiate fra i due
		corpi sono una coppia di forze uguali e opposte
		$\pm R\,\ab$ applicate in $A_i$ e $A_j$, in accordo con la
		molteplicità uno.}
	\label{fig:pendolo_interno}
\end{figure}
%%%%%%%%%%%%%%%%%%%%%%%%%%%%%%%%%%%%%%%%%%%


Il \CRr\ $C_{ij}$ giace sulla retta dell'asse del vincolo
--- la retta per $A$ in direzione $\ab$ nel carrello, la
retta per $A_i$ e $A_j$ nella biella finita. Le reazioni
scambiate fra i due corpi sono una coppia di forze uguali
e opposte applicate in $A_i$ e $A_j$, parallele ad $\ab$,
e nessuna coppia scambiata:
\begin{equation}\label{eq:reazione_carrello_int}
	\rb_i = -R(A)\,\ab,
	\qquad
	\rb_j = R(A)\,\ab,
	\qquad
	\mu(A) = 0,
	\qquad
	R(A) \in \mathbb{R}.
\end{equation}
Sono noti a priori i punti di applicazione $A_i$ e $A_j$,
la direzione $\ab$ delle forze e l'assenza di coppia
scambiata, $\mu(A) = 0$; l'unica incognita scalare è
l'intensità con segno $R(A)$, il cui segno determina se
la forza trasmessa da $\mathcal{C}_i$ a $\mathcal{C}_j$
è concorde o discorde con $\ab$. La
\FIGref{fig:pendolo_interno} illustra le prestazioni
cinematica e statica della biella interna fra due corpi.




\subsection{Il bipendolo}\label{sec:bipendolo}

\paragraph{Realizzazione meccanica e parametri caratterizzanti.}
Il bipendolo\index{bipendolo|textbf}\index{bipendolo|seealso{vincolo}}\index{pendolo!doppio|see{bipendolo}} è un vincolo che impedisce la rotazione del corpo
lasciandone libere entrambe le traslazioni nel piano. Una
realizzazione meccanica tipica è il \emph{pantografo}\index{pantografo|seealso{bipendolo}}: una
catena di aste articolate che trasmette al corpo qualsiasi
traslazione rigida senza trasmettere alcuna rotazione, come
illustrato nella \FIGref{fig:bipendolo_esterno}. Il vincolo
non è caratterizzato né da un punto di applicazione né da una
direzione preferenziale nel piano: l'unico parametro che lo
definisce è l'impedimento della rotazione attorno all'asse
uscente dal piano.


\paragraph{Caso esterno.}

\paragraph{Equazione cinematica.}
\begin{equation}\label{eq:bipendolo_cin}
	\vartheta = 0.
\end{equation}
Il campo di spostamento ammesso è una pura traslazione rigida:
$\ub(P) = \ub$ per ogni punto $P$ del corpo, con $\ub \in
\mathbb{R}^2$ arbitrario.

\paragraph{Posizione del centro di rotazione assoluto.}
Per $\vartheta = 0$ il centro di rotazione è all'infinito:
il corpo trasla senza ruotare e non esiste alcun punto fisso
nel piano.

\begin{proposizione}[Centro di rotazione --- bipendolo]
	Il \CRa\ di un corpo vincolato da un bipendolo è posto
	all'infinito. La direzione verso la quale si trova è
	perpendicolare alla traslazione residua, ed è indeterminata
	finché non si aggiungono ulteriori vincoli.
\end{proposizione}


%%%%%%%%%%%%%%%%%%%%%%%%%%%%%%%%%%%%%%%%%%%
\begin{figure}[htbp]
	\centering
	%\resizebox{\textwidth}{!}{%
		\begin{tikzpicture}[>=Stealth, line width=0.8pt, scale=0.85]
			
			% ═══════════════════════════════════════════════════════════
			% PANNELLO (a) — Prestazione cinematica del bipendolo
			% (realizzazione: pantografo a parallelogramma)
			% ═══════════════════════════════════════════════════════════
			\begin{scope}[xshift=0cm]
				
				%% ── Parametri ──
				\def\xAuno{-0.40}
				\def\xAdue{ 0.40}
				\def\yA{ 0.00}
				\def\Lbiella{0.70}      % era 0.90
				\def\angSup{-110}
				\def\angInf{-70}
				\def\rn{0.10}
				\def\tw{0.22}
				\def\thA{0.32}
				\def\gw{0.40}
				
				% Coordinate calcolate
				\pgfmathsetmacro{\dxSup}{\Lbiella*cos(\angSup)}
				\pgfmathsetmacro{\dySup}{\Lbiella*sin(\angSup)}
				\pgfmathsetmacro{\dxInf}{\Lbiella*cos(\angInf)}
				\pgfmathsetmacro{\dyInf}{\Lbiella*sin(\angInf)}
				% B_1, B_2 (nodi traversa)
				\pgfmathsetmacro{\xBuno}{\xAuno+\dxSup}
				\pgfmathsetmacro{\yB}{\yA+\dySup}
				\pgfmathsetmacro{\xBdue}{\xAdue+\dxSup}
				% G_1, G_2 (cerniere a terra)
				\pgfmathsetmacro{\xGuno}{\xBuno+\dxInf}
				\pgfmathsetmacro{\yG}{\yB+\dyInf}
				\pgfmathsetmacro{\xGdue}{\xBdue+\dxInf}
				
				
				%% ── Corpo C (più grande, agganciato a A_1=-0.40 e A_2=+0.40) ──
				\draw[line width=1pt] 
				plot[smooth cycle, tension=0.7] coordinates 
				{(-0.40, 0.00) (-1.30, 0.35) (-1.80, 1.15) 
					(-1.45, 2.05) (-0.55, 2.40) (0.55, 2.40) 
					(1.45, 2.05) (1.80, 1.15) (1.30, 0.35) (0.40, 0.00)};
				\node at (-0.65, 1.65) {$\mathcal{C}$};
				\node[font=\footnotesize] at (0.00, 0.75) {$\vartheta = 0$};
				
				%% ── Bielle superiori (parallelogramma 1) ──
				\draw[line width=0.8pt] (\xAuno, \yA) -- (\xBuno, \yB);
				\draw[line width=0.8pt] (\xAdue, \yA) -- (\xBdue, \yB);
				
				%% ── Traversa B_1 B_2 ──
				\draw[line width=0.8pt] (\xBuno, \yB) -- (\xBdue, \yB);
				
				%% ── Bielle inferiori (parallelogramma 2) ──
				\draw[line width=0.8pt] (\xBuno, \yB) -- (\xGuno, \yG);
				\draw[line width=0.8pt] (\xBdue, \yB) -- (\xGdue, \yG);
				
				%% ── Cerniere al suolo G_1, G_2 ──
				\draw (\xGuno, \yG) 
				-- ++(-\tw, -\thA) -- ++(2*\tw, 0) -- cycle;
				\draw (\xGuno-\gw, \yG-\thA) -- (\xGuno+\gw, \yG-\thA);
				\foreach \x in {-0.30,-0.15,0,0.15,0.30}
				\draw[thin] (\xGuno+\x, \yG-\thA) -- ++(-0.10, -0.13);
				
				\draw (\xGdue, \yG) 
				-- ++(-\tw, -\thA) -- ++(2*\tw, 0) -- cycle;
				\draw (\xGdue-\gw, \yG-\thA) -- (\xGdue+\gw, \yG-\thA);
				\foreach \x in {-0.30,-0.15,0,0.15,0.30}
				\draw[thin] (\xGdue+\x, \yG-\thA) -- ++(-0.10, -0.13);
				
				%% ── Anellini ──
				% A_1, A_2: semicerchi (corpo C ruota libero rispetto alla biella)
				\draw[fill=white] 
				(\xAuno-\rn, \yA) arc(180:360:\rn) -- cycle;
				\draw[fill=white] 
				(\xAdue-\rn, \yA) arc(180:360:\rn) -- cycle;
				% B_1, B_2: cerchi pieni (nodi del parallelogramma)
				\draw[fill=white] (\xBuno, \yB) circle (3pt);
				\draw[fill=white] (\xBdue, \yB) circle (3pt);
				% G_1, G_2: cerchi pieni (perni al suolo)
				\draw[fill=white] (\xGuno, \yG) circle (3pt);
				\draw[fill=white] (\xGdue, \yG) circle (3pt);
				
				%				
				%% ── Retta impropria r_∞ (luogo del CR all'infinito) ──
				\def\xRinf{-2.40}
				\def\yRinf{ 1.20}
				\def\angRinf{60}
				\def\Lrinf{1.20}
				\pgfmathsetmacro{\dxR}{\Lrinf*cos(\angRinf)}
				\pgfmathsetmacro{\dyR}{\Lrinf*sin(\angRinf)}
				\draw[gray!60, line width=0.6pt] 
				({\xRinf-\dxR}, {\yRinf-\dyR}) -- ({\xRinf+\dxR}, {\yRinf+\dyR});
				
				%% Etichetta "r_∞" all'estremità superiore della retta
				\node[gray, font=\footnotesize, anchor=south west] 
				at ({\xRinf+\dxR+0.05}, {\yRinf+\dyR+0.05}) {$r_\infty$};
				
				%% Etichetta "C ∈ r_∞" sopra la retta a metà, ruotata
				\node[gray, font=\footnotesize, rotate=\angRinf, anchor=south, 
				inner sep=2pt] 
				at (\xRinf, \yRinf) {$C \in r_\infty$};
				
				
				%% ── Etichetta pannello ──
				\node[anchor=north, font=\small] 
				at (-0.25, \yG-\thA-0.50) {$(a)$};
				
			\end{scope}
			
			% ═══════════════════════════════════════════════════════════
			% PANNELLO (b) — Prestazione statica del bipendolo
			% (realizzazione: pantografo a parallelogramma)
			% ═══════════════════════════════════════════════════════════
			\begin{scope}[xshift=5cm]
				
				%% ── Parametri (identici ad (a)) ──
				%% ── Parametri ──
				\def\xAuno{-0.40}
				\def\xAdue{ 0.40}
				\def\yA{ 0.00}
				\def\Lbiella{0.70}      % era 0.90
				\def\angSup{-110}
				\def\angInf{-70}
				\def\rn{0.10}
				\def\tw{0.22}
				\def\thA{0.32}
				\def\gw{0.40}
				
				% Coordinate calcolate
				\pgfmathsetmacro{\dxSup}{\Lbiella*cos(\angSup)}
				\pgfmathsetmacro{\dySup}{\Lbiella*sin(\angSup)}
				\pgfmathsetmacro{\dxInf}{\Lbiella*cos(\angInf)}
				\pgfmathsetmacro{\dyInf}{\Lbiella*sin(\angInf)}
				\pgfmathsetmacro{\xBuno}{\xAuno+\dxSup}
				\pgfmathsetmacro{\yB}{\yA+\dySup}
				\pgfmathsetmacro{\xBdue}{\xAdue+\dxSup}
				\pgfmathsetmacro{\xGuno}{\xBuno+\dxInf}
				\pgfmathsetmacro{\yG}{\yB+\dyInf}
				\pgfmathsetmacro{\xGdue}{\xBdue+\dxInf}
				
				%% ── Corpo C (identico ad (a)) ──
				\draw[line width=1pt] 
				plot[smooth cycle, tension=0.7] coordinates 
				{(-0.40, 0.00) (-1.30, 0.35) (-1.80, 1.15) 
					(-1.45, 2.05) (-0.55, 2.40) (0.55, 2.40) 
					(1.45, 2.05) (1.80, 1.15) (1.30, 0.35) (0.40, 0.00)};
				\node at (-0.65, 1.65) {$\mathcal{C}$};
				
				
				
				%% ── Bielle superiori, traversa, bielle inferiori (identici ad (a)) ──
				\draw[line width=0.8pt] (\xAuno, \yA) -- (\xBuno, \yB);
				\draw[line width=0.8pt] (\xAdue, \yA) -- (\xBdue, \yB);
				\draw[line width=0.8pt] (\xBuno, \yB) -- (\xBdue, \yB);
				\draw[line width=0.8pt] (\xBuno, \yB) -- (\xGuno, \yG);
				\draw[line width=0.8pt] (\xBdue, \yB) -- (\xGdue, \yG);
				
				%% ── Cerniere al suolo G_1, G_2 (identiche ad (a)) ──
				\draw (\xGuno, \yG) 
				-- ++(-\tw, -\thA) -- ++(2*\tw, 0) -- cycle;
				\draw (\xGuno-\gw, \yG-\thA) -- (\xGuno+\gw, \yG-\thA);
				\foreach \x in {-0.30,-0.15,0,0.15,0.30}
				\draw[thin] (\xGuno+\x, \yG-\thA) -- ++(-0.10, -0.13);
				\draw (\xGdue, \yG) 
				-- ++(-\tw, -\thA) -- ++(2*\tw, 0) -- cycle;
				\draw (\xGdue-\gw, \yG-\thA) -- (\xGdue+\gw, \yG-\thA);
				\foreach \x in {-0.30,-0.15,0,0.15,0.30}
				\draw[thin] (\xGdue+\x, \yG-\thA) -- ++(-0.10, -0.13);
				
				%% ── Anellini (identici ad (a)) ──
				\draw[fill=white] 
				(\xAuno-\rn, \yA) arc(180:360:\rn) -- cycle;
				\draw[fill=white] 
				(\xAdue-\rn, \yA) arc(180:360:\rn) -- cycle;
				\draw[fill=white] (\xBuno, \yB) circle (3pt);
				\draw[fill=white] (\xBdue, \yB) circle (3pt);
				\draw[fill=white] (\xGuno, \yG) circle (3pt);
				\draw[fill=white] (\xGdue, \yG) circle (3pt);
				
				%% ──────────────────────────────────────────────────
				%% COPPIA REATTIVA mu — arco con freccia sopra C
				%% ──────────────────────────────────────────────────
				\def\xMu{0.00}
				\def\yMu{0.50}
				\def\rMu{0.40}
				\draw[->, line width=1.4pt] ({\xMu+\rMu}, \yMu) arc[start angle=0, end angle=180, radius=\rMu];
				\node[font=\footnotesize] at (0.00, 0.5) {$\mu$};
				
				%% ── Etichetta pannello ──
				\node[anchor=north, font=\small] at (-0.25, \yG-\thA-0.50) {$(b)$};
				
			\end{scope}
			
		\end{tikzpicture}
		%	}
	
	\caption[Bipendolo esterno: prestazioni cinematica e statica.]{%
		Bipendolo esterno applicato al corpo $\mathcal{C}$,
		realizzato meccanicamente come pantografo a doppio
		parallelogramma. Prestazione cinematica (a): l'equazione
		di vincolo $\vartheta = 0$ annulla la rotazione del
		corpo, lasciando libere le due componenti di traslazione,
		non rappresentate; il \CRa\ giace sulla retta impropria
		$r_\infty$. Prestazione statica (b): la reazione vincolare
		si riduce alla sola coppia $\mu$, in dualità con la
		rotazione bloccata; la forza risultante è identicamente
		nulla, in accordo con la molteplicità uno.}
	\label{fig:bipendolo_esterno}
\end{figure}
%%%%%%%%%%%%%%%%%%%%%%%%%%%%%%%%%%%%%%%%%%%

\paragraph{Reazione vincolare.}
Gli spostamenti compatibili con il bipendolo sono tutte le
traslazioni rigide ($\vartheta = 0$, $\ub$ libero). Il lavoro
di una forza risultante $\rb$ su una traslazione $\ub$ è
$\rb\cdot\ub$: perché sia nullo per qualsiasi $\ub$, deve
essere $\rb = \zerovb$. La reazione è quindi una \emph{coppia
	pura}:
\begin{equation}\label{eq:reazione_bipendolo}
	\rb = \zerovb,
	\qquad
	\mu(A) \in \mathbb{R}.
\end{equation}
È noto a priori che il bipendolo non oppone alcuna forza
alle traslazioni, $\rb = \zerovb$, mentre reagisce a
qualsiasi tentativo di rotazione con una coppia di verso e
intensità non noti a priori; l'unica incognita scalare è
$\mu(A)$, il cui segno determina se la coppia è oraria o
antioraria, in accordo con la molteplicità uno.


\begin{oss}
	La dualità cinematica-statica è qui particolarmente
	evidente: la variabile cinematica vincolata è $\vartheta$
	(una rotazione) e la corrispondente reazione è $\mu$
	(una coppia), mentre le variabili libere sono le due
	traslazioni e le corrispondenti reazioni sono nulle.
\end{oss}



\paragraph{Caso interno.}
Il bipendolo interno\index{bipendolo!interno} fra $\mathcal{C}_i$ e $\mathcal{C}_j$
impone l'uguaglianza delle rotazioni dei due corpi,
scritta come annullamento del salto della rotazione,
\begin{equation}\label{eq:bipendolo_int}
	\llbracket \vartheta \rrbracket = 0,
	\qquad
	\llbracket \vartheta \rrbracket \coloneqq
	\vartheta_j - \vartheta_i,
\end{equation}
e i due corpi ruotano della stessa quantità. La condizione
non coinvolge punti, coerentemente col fatto che la
rotazione di un corpo rigido è una grandezza globale; il
\CRr\ $C_{ij}$ è all'infinito.

Le reazioni scambiate fra i due corpi sono una forza nulla
e una coppia di verso e intensità non noti a priori, uguali
e opposte:
\begin{equation}\label{eq:reazione_bipendolo_int}
	\rb_i = \rb_j = \zerovb,
	\qquad
	\mu(A) \in \mathbb{R}.
\end{equation}

%%%%%%%%%%%%%%%%%%%%%%%%%%%%%%%%%%%%%%%%%%%
\begin{figure}[htbp]
	\centering
	\begin{tikzpicture}[>=Stealth, line width=0.8pt, scale=0.85]
		
		% ═══════════════════════════════════════════════════════════
		% PANNELLO (a) — Prestazione cinematica del bipendolo interno
		% (realizzazione: pantografo a doppio parallelogramma)
		% ═══════════════════════════════════════════════════════════
		\begin{scope}[xshift=0cm]
			
			%% ── Parametri ──
			\def\xAi{ 0.00}
			\def\yAuno{ 0.35}
			\def\yAdue{-0.35}
			\def\Lbiella{0.65}
			\def\angInt{-25}
			\def\angExt{ 25}
			\def\rn{0.10}
			
			% Coordinate calcolate
			\pgfmathsetmacro{\dxInt}{\Lbiella*cos(\angInt)}
			\pgfmathsetmacro{\dyInt}{\Lbiella*sin(\angInt)}
			\pgfmathsetmacro{\dxExt}{\Lbiella*cos(\angExt)}
			\pgfmathsetmacro{\dyExt}{\Lbiella*sin(\angExt)}
			\pgfmathsetmacro{\xBuno}{\xAi+\dxInt}
			\pgfmathsetmacro{\yBuno}{\yAuno+\dyInt}
			\pgfmathsetmacro{\xBdue}{\xAi+\dxInt}
			\pgfmathsetmacro{\yBdue}{\yAdue+\dyInt}
			\pgfmathsetmacro{\xAj}{\xBuno+\dxExt}
			\pgfmathsetmacro{\yAjuno}{\yBuno+\dyExt}
			\pgfmathsetmacro{\yAjdue}{\yBdue+\dyExt}
			
			%% ── Corpo C_i (a sinistra, ingrandito, attaccato in y=±0.35) ──
			\draw[line width=1pt] 
			plot[smooth cycle, tension=0.7] coordinates 
			{(0.00, 0.35) (-0.35, 1.10) (-1.00, 1.55) 
				(-1.95, 1.35) (-2.40, 0.55) (-2.40, -0.55)
				(-1.95, -1.35) (-1.00, -1.55) (-0.35, -1.10) 
				(0.00, -0.35)};
			\node at (-1.25, 0.80) {$\mathcal{C}_i$};
			
			%% ── Corpo C_j (a destra, ingrandito, attaccato in y=±0.35) ──
			\draw[line width=1pt] 
			plot[smooth cycle, tension=0.7] coordinates 
			{(\xAj, 0.35) ({\xAj+0.35}, 1.10) ({\xAj+1.00}, 1.55) 
				({\xAj+1.95}, 1.35) ({\xAj+2.40}, 0.55) 
				({\xAj+2.40}, -0.55) ({\xAj+1.95}, -1.35) 
				({\xAj+1.00}, -1.55) ({\xAj+0.35}, -1.10) 
				(\xAj, -0.35)};
			\node at ({\xAj+1.25}, 0.80) {$\mathcal{C}_j$};
			
			
			%% ── Bielle interne (parallelogramma 1) ──
			\draw[line width=0.8pt] (\xAi, \yAuno) -- (\xBuno, \yBuno);
			\draw[line width=0.8pt] (\xAi, \yAdue) -- (\xBdue, \yBdue);
			
			%% ── Traversa B_1 B_2 (verticale) ──
			\draw[line width=0.8pt] (\xBuno, \yBuno) -- (\xBdue, \yBdue);
			
			%% ── Bielle esterne (parallelogramma 2) ──
			\draw[line width=0.8pt] (\xBuno, \yBuno) -- (\xAj, \yAjuno);
			\draw[line width=0.8pt] (\xBdue, \yBdue) -- (\xAj, \yAjdue);
			
			%% ── Anellini ──
			\draw[fill=white] (\xBuno, \yBuno) circle (3pt);
			\draw[fill=white] (\xBdue, \yBdue) circle (3pt);
			
			
			%% ── Anellini terminali ──
			% A^i (bordo destro di C_i, semicerchio sporge a destra verso pantografo)
			\draw[line width=0.8pt, fill=white] 
			(\xAi, \yAuno-\rn) arc(-90:90:\rn) -- cycle;
			\draw[line width=0.8pt, fill=white] 
			(\xAi, \yAdue-\rn) arc(-90:90:\rn) -- cycle;
			% A^j (bordo sinistro di C_j, semicerchio sporge a sinistra verso pantografo)
			\draw[line width=0.8pt, fill=white] 
			(\xAj, \yAjuno+\rn) arc(90:270:\rn) -- cycle;
			\draw[line width=0.8pt, fill=white] 
			(\xAj, \yAjdue+\rn) arc(90:270:\rn) -- cycle;
			% B_1, B_2: cerchi pieni
			\draw[fill=white] (\xBuno, \yBuno) circle (3pt);
			\draw[fill=white] (\xBdue, \yBdue) circle (3pt);
			
			
			%% ── Archi delle rotazioni ϑ_i e ϑ_j (semicerchio, 180°) ──
			\def\rrot{0.30}
			% In C_i
			\draw[->, line width=1pt] 
			(-1.25+\rrot, -0.20) arc(0:180:\rrot);
			\node[font=\footnotesize, anchor=west] 	at (-2.0+\rrot+0.05, -0.40) {$\vartheta_i$};
			% In C_j
			\draw[->, line width=1pt] ({\xAj+1.25+\rrot}, -0.20) arc(0:180:\rrot);
			\node[font=\footnotesize, anchor=west] 	at (+1.80+\rrot+0.05, -0.40) {$\vartheta_j$};
			
			%% ── Equazione di vincolo (in basso al centro) ──
			\node[font=\small] 
			at ({(\xAi+\xAj)/2}, 1.80) 
			{$\vartheta_j - \vartheta_i = 0$};
			
			%% ── Retta impropria r_∞ (in alto a sinistra) ──
			\def\xRinf{-1.80}
			\def\yRinf{ 2.20}
			\def\angRinf{15}
			\def\Lrinf{1.20}
			\pgfmathsetmacro{\dxR}{\Lrinf*cos(\angRinf)}
			\pgfmathsetmacro{\dyR}{\Lrinf*sin(\angRinf)}
			\draw[gray!60, line width=0.6pt] 
			({\xRinf-\dxR}, {\yRinf-\dyR}) 
			-- ({\xRinf+\dxR}, {\yRinf+\dyR});
			\node[gray, font=\footnotesize, anchor=south west] 
			at ({\xRinf+\dxR+0.05}, {\yRinf+\dyR+0.05}) {$r_\infty$};
			\node[gray, font=\footnotesize, rotate=\angRinf, 
			anchor=south, inner sep=2pt] 
			at (\xRinf, \yRinf) {$C_{ij} \in r_\infty$};
			
			%% ── Etichetta pannello ──
			\node[anchor=north, font=\small] 
			at ({(\xAi+\xAj)/2}, -1.50) {$(a)$};
			
		\end{scope}
		% ═══════════════════════════════════════════════════════════
		% PANNELLO (b) — Prestazione statica del bipendolo interno
		% (realizzazione: pantografo a doppio parallelogramma)
		% ═══════════════════════════════════════════════════════════
		\begin{scope}[xshift=7cm]
			
			%% ── Parametri ──
			\def\xAi{ 0.00}
			\def\yAuno{ 0.35}
			\def\yAdue{-0.35}
			\def\Lbiella{0.65}
			\def\angInt{-25}
			\def\angExt{ 25}
			\def\rn{0.10}
			
			% Coordinate calcolate
			\pgfmathsetmacro{\dxInt}{\Lbiella*cos(\angInt)}
			\pgfmathsetmacro{\dyInt}{\Lbiella*sin(\angInt)}
			\pgfmathsetmacro{\dxExt}{\Lbiella*cos(\angExt)}
			\pgfmathsetmacro{\dyExt}{\Lbiella*sin(\angExt)}
			\pgfmathsetmacro{\xBuno}{\xAi+\dxInt}
			\pgfmathsetmacro{\yBuno}{\yAuno+\dyInt}
			\pgfmathsetmacro{\xBdue}{\xAi+\dxInt}
			\pgfmathsetmacro{\yBdue}{\yAdue+\dyInt}
			\pgfmathsetmacro{\xAj}{\xBuno+\dxExt}
			\pgfmathsetmacro{\yAjuno}{\yBuno+\dyExt}
			\pgfmathsetmacro{\yAjdue}{\yBdue+\dyExt}
			
			%% ── Corpo C_i (a sinistra, ingrandito, attaccato in y=±0.35) ──
			\draw[line width=1pt] 
			plot[smooth cycle, tension=0.7] coordinates 
			{(0.00, 0.35) (-0.35, 1.10) (-1.00, 1.55) 
				(-1.95, 1.35) (-2.40, 0.55) (-2.40, -0.55)
				(-1.95, -1.35) (-1.00, -1.55) (-0.35, -1.10) 
				(0.00, -0.35)};
			\node at (-1.25, 0.80) {$\mathcal{C}_i$};
			
			%% ── Corpo C_j (a destra, ingrandito, attaccato in y=±0.35) ──
			\draw[line width=1pt] 
			plot[smooth cycle, tension=0.7] coordinates 
			{(\xAj, 0.35) ({\xAj+0.35}, 1.10) ({\xAj+1.00}, 1.55) 
				({\xAj+1.95}, 1.35) ({\xAj+2.40}, 0.55) 
				({\xAj+2.40}, -0.55) ({\xAj+1.95}, -1.35) 
				({\xAj+1.00}, -1.55) ({\xAj+0.35}, -1.10) 
				(\xAj, -0.35)};
			\node at ({\xAj+1.25}, 0.80) {$\mathcal{C}_j$};
			
			%% ── Bielle interne, traversa, bielle esterne ──
			\draw[line width=0.8pt] (\xAi, \yAuno) -- (\xBuno, \yBuno);
			\draw[line width=0.8pt] (\xAi, \yAdue) -- (\xBdue, \yBdue);
			\draw[line width=0.8pt] (\xBuno, \yBuno) -- (\xBdue, \yBdue);
			\draw[line width=0.8pt] (\xBuno, \yBuno) -- (\xAj, \yAjuno);
			\draw[line width=0.8pt] (\xBdue, \yBdue) -- (\xAj, \yAjdue);
			
			%% ── Anellini terminali (identici ad (a)) ──
			\draw[line width=0.8pt, fill=white] 
			(\xAi, \yAuno-\rn) arc(-90:90:\rn) -- cycle;
			\draw[line width=0.8pt, fill=white] 
			(\xAi, \yAdue-\rn) arc(-90:90:\rn) -- cycle;
			\draw[line width=0.8pt, fill=white] 
			(\xAj, \yAjuno+\rn) arc(90:270:\rn) -- cycle;
			\draw[line width=0.8pt, fill=white] 
			(\xAj, \yAjdue+\rn) arc(90:270:\rn) -- cycle;
			\draw[fill=white] (\xBuno, \yBuno) circle (3pt);
			\draw[fill=white] (\xBdue, \yBdue) circle (3pt);
			
			%% ──────────────────────────────────────────────────
			%% COPPIE REATTIVE μ (uguali e opposte)
			%% ──────────────────────────────────────────────────
			\def\rmu{0.40}
			\def\yMu{-0.20}
			\def\xMuI{-1.25}        % era -1.10
			\def\xMuJ{\xAj+1.25}    % era \xAj+1.10
			
			% Coppia in C_i: semicerchio antiorario
			\draw[<-, line width=1.4pt] (\xMuI+\rmu, \yMu) arc(0:180:\rmu);
			\node[font=\small] 	at (\xMuI, \yMu) {$\mu$};
			
			% Coppia in C_j: semicerchio orario
			\draw[->, line width=1.4pt] (\xMuJ+\rmu, \yMu) arc(0:180:\rmu);
			\node[font=\small] 	at (\xMuJ, \yMu) {$\mu$};
			
			%% ── Etichetta pannello ──
			\node[anchor=north, font=\small] 
			at ({(\xAi+\xAj)/2}, -1.50) {$(b)$};
			
		\end{scope}
	\end{tikzpicture}
	\caption[Bipendolo interno fra due corpi: prestazioni cinematica e statica.]{%
		Bipendolo interno fra i corpi $\mathcal{C}_i$ e
		$\mathcal{C}_j$, realizzato come pantografo a doppio
		parallelogramma. Prestazione cinematica (a): l'equazione
		di vincolo $\vartheta_j - \vartheta_i = 0$ impone l'uguaglianza
		delle rotazioni dei due corpi; le componenti di traslazione
		relativa restano libere e non sono rappresentate. Il \CRr\
		$C_{ij}$ giace sulla retta impropria $r_\infty$.
		Prestazione statica (b): i due corpi si scambiano
		una coppia pura di intensità $\mu$, in dualità con la
		rotazione relativa bloccata; la forza scambiata è
		identicamente nulla, in accordo con la molteplicità uno.}
	\label{fig:bipendolo_interno}
\end{figure}
%%%%%%%%%%%%%%%%%%%%%%%%%%%%%%%%%%%%%%%%%%%
Il punto $A$, secondo la convenzione introdotta
nell'osservazione del \PARref{sec:richiamo}, identifica
il giunto da cui la coppia proviene; questo è essenziale
sia per distinguerla da altre coppie reattive eventualmente
agenti sui due corpi, sia per localizzare correttamente
la discontinuità che essa produce nei diagrammi delle
caratteristiche di sollecitazione, quando i corpi
$\mathcal{C}_i$ e $\mathcal{C}_j$ sono tratti contigui di
una stessa trave. È noto a priori che il bipendolo interno
non scambia alcuna forza fra i due corpi,
$\rb_i = \rb_j = \zerovb$; l'unica incognita scalare è
$\mu(A)$. La \FIGref{fig:bipendolo_interno} illustra
le prestazioni cinematica e statica del bipendolo interno
fra due corpi.



%---------------------------------------------------------------------------------------------------------------------------- 
\section{Vincoli doppi: molteplicità due}\label{sec:doppi}
%---------------------------------------------------------------------------------------------------------------------------- 

I vincoli di molteplicità due\index{vincolo!doppio} rimuovono due gradi di libertà,
lasciando al corpo un solo grado di libertà residuo. Le due tipologie
fondamentali sono la \emph{cerniera} e il \emph{glifo}.


\subsection{La cerniera}\label{sec:cerniera}

\paragraph{Realizzazione meccanica e parametri caratterizzanti.}
La cerniera\index{cerniera|textbf}\index{cerniera|seealso{vincolo}} è un vincolo che fissa la posizione di un punto
del corpo lasciandone libera la rotazione. Il corpo può ruotare
attorno al punto di applicazione ma non può traslare in alcuna
direzione. Il vincolo è completamente caratterizzato dal solo
\emph{punto di applicazione} $A$.

\paragraph{Caso esterno.}

\paragraph{Equazione cinematica.}
La condizione di spostamento nullo del punto $A$ fornisce due
equazioni scalari:
\begin{equation}\label{eq:cerniera_cin}
	\ub(A) = \zerovb,
	\qquad\text{ovvero}\qquad
	u(A) = 0, \quad v(A) = 0.
\end{equation}

\paragraph{Posizione del centro di rotazione assoluto.}
Dalla~\eqref{eq:centro_campo}, imponendo $\ub(A) = \zerovb$:
\begin{equation}
	\bm{\upvartheta} \times \overrightarrow{CA} = \zerovb,
\end{equation}
che è soddisfatta se e solo se $C = A$.

\begin{proposizione}[Centro di rotazione --- cerniera]
	\label{prop:CR_cerniera}
	Il \CRa\ di un corpo vincolato da una cerniera in $A$
	coincide con il punto $A$ stesso:
	\begin{equation}\label{eq:CR_cerniera}
		C = A.
	\end{equation}
	Il centro è univocamente determinato dalla posizione
	della cerniera: qualsiasi spostamento compatibile è
	una rotazione rigida attorno ad $A$.
\end{proposizione}

\begin{oss}
	La cerniera può essere interpretata come la
	sovrapposizione di due carrelli ortogonali applicati
	in $A$: uno con asse $\ab_x$, che impone $u(A) = 0$,
	e uno con asse $\ab_y$, che impone $v(A) = 0$. Ciascun
	carrello colloca il centro su una retta passante per
	$A$; l'intersezione delle due rette --- ortogonali fra
	loro --- è il punto $A$ stesso.
\end{oss}

\paragraph{Reazione vincolare.}
Gli spostamenti compatibili con la cerniera sono le rotazioni
rigide attorno ad $A$: il punto $A$ è fisso e qualsiasi forza
concentrata in $A$ compie lavoro nullo. La reazione è quindi
una \emph{forza arbitraria} applicata in $A$, senza coppia:
\begin{equation}\label{eq:reazione_cerniera}
	\rb(A) = X(A)\,\ab_x + Y(A)\,\ab_y,
	\qquad
	\mu(A) = 0,
	\qquad
	X(A), Y(A) \in \mathbb{R}.
\end{equation}
Sono noti a priori il punto di applicazione $A$ e l'assenza
di coppia reattiva, $\mu(A) = 0$; la direzione, il verso e
l'intensità della forza reattiva non sono noti a priori e
sono sintetizzati dalle due componenti scalari $X(A)$ e
$Y(A)$, ciascuna con segno libero, in accordo con la
molteplicità due. La \FIGref{fig:cerniera_esterna} illustra
le prestazioni cinematica e statica della cerniera esterna.


%%%%%%%%%%%%%%%%%%%%%%%%%%%%%%%%%%%%%%%%%%%
\begin{figure}[htbp]
	\centering
	\begin{tikzpicture}[>=Stealth, line width=0.8pt]
		
		% ═══════════════════════════════════════════════════════════
		% PANNELLO (a) — Prestazione cinematica della cerniera esterna
		% ═══════════════════════════════════════════════════════════
		\begin{scope}[xshift=0cm]
			
			%% ── Parametri ──
			\def\xA{0.80}
			\def\yA{0.00}
			\def\tw{0.30}
			\def\thA{0.40}
			\def\gw{0.45}
			\def\rn{0.10}    % raggio semicerchio (anellino a metà)
			
			%% ── Terna x, y in basso a sinistra ──
			\draw[->, line width=0.8pt] (-2.20, 0.00) -- (-1.40, 0.00);
			\node[anchor=west, font=\footnotesize] at (-1.40, -0.05) {$x$};
			\draw[->, line width=0.8pt] (-2.20, 0.00) -- (-2.20, 0.80);
			\node[anchor=south, font=\footnotesize] at (-2.20, 0.80) {$y$};
			
			%% ── Equazione di vincolo (in alto) ──
			\node[font=\small] at (\xA-1.5, 2.0) {$\ub(A) = \zerovb$};
			
			
			
			%% ── Corpo C ──
			\draw[line width=1pt] 
			plot[smooth cycle, tension=0.7] coordinates 
			{(0.80, 0.00) (1.60, 0.15) (2.20, 0.70) (2.00, 1.70) 
				(1.00, 2.00) (-0.10, 1.70) (-0.60, 1.00) (-0.30, 0.30)};
			\node at (1.50, 1.10) {$\mathcal{C}$};
			
			
			%% ── Cerniera in A: triangolo + barra + tratteggio ──
			\draw (\xA, \yA) -- ++(-\tw, -\thA) -- ++(2*\tw, 0) -- cycle;
			\draw (\xA-\gw, \yA-\thA) -- (\xA+\gw, \yA-\thA);
			\foreach \x in {-0.36,-0.18,0,0.18,0.36}
			\draw[thin] (\xA+\x, \yA-\thA) -- ++(-0.10, -0.13);
			
			
			%% ── Semicerchio in A (corpo libero di ruotare sul perno) ──
			\draw[line width=0.8pt, fill=white] ({\xA-\rn}, \yA) arc(180:360:\rn) -- cycle;
			
			
			
			%% ── Etichetta A ──
			\node[above, font=\small] at (\xA, \yA+0.05) {$A$};
			
			%% ── Etichetta C ≡ A ──
			\node[font=\footnotesize, anchor=west] 
			at (\xA+0.50, \yA-0.20) {$C \equiv A$};
			
			%% ── Etichetta pannello ──
			\node[anchor=north, font=\small] 
			at (\xA-0.8, \yA-\thA-0.30) {$(a)$};
			
		\end{scope}
		
		% ═══════════════════════════════════════════════════════════
		% PANNELLO (b) — Prestazione statica della cerniera esterna
		% ═══════════════════════════════════════════════════════════
		\begin{scope}[xshift=6cm]
			
			%% ── Parametri (identici ad (a)) ──
			\def\xA{0.80}
			\def\yA{0.00}
			\def\tw{0.30}
			\def\thA{0.40}
			\def\gw{0.45}
			\def\rn{0.10}    % raggio semicerchio (anellino a metà)
			
			%% ── Terna x, y in basso a sinistra (identica ad (a)) ──
			\draw[->, line width=0.8pt] (-2.20, 0.00) -- (-1.40, 0.00);
			\node[anchor=west, font=\footnotesize] at (-1.40, -0.05) {$x$};
			\draw[->, line width=0.8pt] (-2.20, 0.00) -- (-2.20, 0.80);
			\node[anchor=south, font=\footnotesize] at (-2.20, 0.80) {$y$};
			
			%% ── Corpo C (identico ad (a)) ──
			\draw[line width=1pt] 
			plot[smooth cycle, tension=0.7] coordinates 
			{(0.80, 0.00) (1.60, 0.15) (2.20, 0.70) (2.00, 1.70) 
				(1.00, 2.00) (-0.10, 1.70) (-0.60, 1.00) (-0.30, 0.30)};
			\node at (1.50, 0.50) {$\mathcal{C}$};
			
			%% ── Cerniera in A: triangolo + barra + tratteggio (identica ad (a)) ──
			\draw (\xA, \yA) -- ++(-\tw, -\thA) -- ++(2*\tw, 0) -- cycle;
			\draw (\xA-\gw, \yA-\thA) -- (\xA+\gw, \yA-\thA);
			\foreach \x in {-0.36,-0.18,0,0.18,0.36}
			\draw[thin] (\xA+\x, \yA-\thA) -- ++(-0.10, -0.13);
			
			
			%% ── Semicerchio in A (corpo libero di ruotare sul perno) ──
			\draw[line width=0.8pt, fill=white] ({\xA-\rn}, \yA) arc(180:360:\rn) -- cycle;
			
			%% ──────────────────────────────────────────────────
			%% REAZIONE: due componenti X(A) e Y(A)
			%% ──────────────────────────────────────────────────
			
			%% Componente orizzontale X(A) a_x: da A verso destra
			\draw[->, line width=1.4pt] 
			(\xA+0.10, \yA) -- (\xA+1.10, \yA);
			\node[anchor=west, font=\small] 
			at (\xA+1.15, \yA) {$X(A)\,\ab_x$};
			
			%% Componente verticale Y(A) a_y: da A verso l'alto
			\draw[->, line width=1.4pt] 
			(\xA, \yA) -- (\xA, \yA+1.0);
			\node[anchor=south, font=\small] 
			at (\xA, \yA+1.15) {$Y(A)\,\ab_y$};
			
			%% ── Etichetta A (a sinistra del perno) ──
			\node[anchor=east, font=\small] at (\xA, \yA+0.3) {$A$};
			
			%% ── Etichetta pannello ──
			\node[anchor=north, font=\small] 
			at (\xA-0.8, \yA-\thA-0.30) {$(b)$};
			
		\end{scope}
		
	\end{tikzpicture}
	\caption[Cerniera esterna applicata in $A$: prestazioni cinematica e statica.]{%
		Cerniera esterna applicata al corpo $\mathcal{C}$ nel
		punto $A$. Prestazione cinematica (a): il punto $A$ è
		fisso, $\ub(A) = \zerovb$, e il corpo ruota liberamente
		attorno ad esso; il \CRa\ coincide con il punto di
		applicazione, $C \equiv A$. Prestazione statica (b): la
		reazione vincolare è una forza generica
		$X(A)\,\ab_x + Y(A)\,\ab_y$ applicata in $A$, in dualità
		con le due componenti di spostamento bloccate; la coppia
		reattiva è identicamente nulla, in accordo con la
		molteplicità due.}
	\label{fig:cerniera_esterna}
\end{figure}
%%%%%%%%%%%%%%%%%%%%%%%%%%%%%%%%%%%%%%%%%%%

\paragraph{Caso interno.}
La cerniera interna\index{cerniera!interna} è un vincolo puntiforme: i punti di
applicazione $A_i \in \mathcal{C}_i$ e $A_j \in \mathcal{C}_j$
coincidono geometricamente in un unico punto $A$. La
condizione di uguaglianza degli spostamenti dei due corpi in
$A$ si scrive come annullamento del salto attraverso $A$:
\begin{equation}\label{eq:cerniera_int}
	\llbracket \ub \rrbracket(A) = \zerovb,
	\qquad
	\llbracket \ub \rrbracket(A) \coloneqq \ub_j(A) - \ub_i(A),
\end{equation}
in accordo con la convenzione del
\PARref{sec:salto_incremento}. Il \CRr\ $C_{ij}$ coincide
con $A$. Le reazioni scambiate fra i due corpi sono forze
uguali e opposte applicate in $A$, di direzione, verso e
intensità non noti a priori, e nessuna coppia scambiata:
\begin{equation}\label{eq:reazione_cerniera_int}
	\rb_i = -X(A)\,\ab_x - Y(A)\,\ab_y,
	\qquad
	\rb_j = X(A)\,\ab_x + Y(A)\,\ab_y,
	\qquad
	\mu(A) = 0,
\end{equation}
con $X(A), Y(A) \in \mathbb{R}$. Sono noti a priori il punto
di applicazione $A$ e l'assenza di coppia scambiata,
$\mu(A) = 0$; le due incognite scalari $X(A)$ e $Y(A)$,
ciascuna con segno libero, sintetizzano direzione, verso e
intensità della forza trasmessa da $\mathcal{C}_i$ a
$\mathcal{C}_j$, in accordo con la molteplicità due. La
\FIGref{fig:cerniera_interna} illustra le prestazioni
cinematica e statica della cerniera interna fra due corpi.


%%%%%%%%%%%%%%%%%%%%%%%%%%%%%%%%%%%%%%%%%%%
\begin{figure}[htbp]
	\centering
	\begin{tikzpicture}[>=Stealth, line width=0.8pt, scale=0.85]
		
		% ═══════════════════════════════════════════════════════════
		% PANNELLO (a) — Prestazione cinematica della cerniera interna
		% ═══════════════════════════════════════════════════════════
		\begin{scope}[xshift=0cm]
			
			%% ── Parametri ──
			\pgfmathsetmacro{\rn}{0.15}
			\def\xA{0.00}
			\def\yA{0.00}
			
			%% ── Corpo C_i (bordo destro tangente al pallino) ──
			\draw[line width=1pt] 
			plot[smooth cycle, tension=0.7] coordinates 
			{(-0.15, 0.50) (-0.30, 0.85) (-0.65, 1.10) (-1.35, 1.40) 
				(-2.15, 1.20) (-2.55, 0.50) (-2.55, -0.50)
				(-2.15, -1.20) (-1.35, -1.40) (-0.65, -1.10) 
				(-0.30, -0.85) (-0.15, -0.50)};
			\node at (-1.55, 0.0) {$\mathcal{C}_i$};
			
			%% ── Corpo C_j (bordo sinistro tangente al pallino) ──
			\draw[line width=1pt] 
			plot[smooth cycle, tension=0.7] coordinates 
			{(0.15, 0.50) (0.30, 0.85) (0.65, 1.10) (1.35, 1.40) 
				(2.15, 1.20) (2.55, 0.50) (2.55, -0.50)
				(2.15, -1.20) (1.35, -1.40) (0.65, -1.10) 
				(0.30, -0.85) (0.15, -0.50)};
			\node at (1.55, 0.0) {$\mathcal{C}_j$};
			
			
			%% ── Pallino della cerniera in A ──
			\draw[line width=0.8pt, fill=white] (\xA, \yA) circle (\rn);
			
			%% ── Etichette dei punti di applicazione A_i, A_j ──
			\node[font=\footnotesize, anchor=south east] at ({-\rn-0.02}, 0.05) {$A_i$};
			\node[font=\footnotesize, anchor=south west] at ({\rn-0.02}, 0.0) {$A_j$};
			
			
			%% ── Equazione di vincolo (in alto) ──
			\node[font=\small] at (0.00, 1.80) 	{$\ub(A_j) - \ub(A_i) = \zerovb$};
			
			%% ── Archi delle rotazioni ϑ_i, ϑ_j (semicerchio, 180°) ──
			\def\rrot{0.30}
			
			%% ── Etichetta del centro di rotazione relativo (in basso) ──
			\node[font=\footnotesize] 
			at (0.00, -1.80) 
			{$C_{ij} \equiv A_i \equiv A_j$};
			
			%% ── Etichetta pannello ──
			\node[anchor=north, font=\small] 
			at (0.00, -2.30) {$(a)$};
			
		\end{scope}
		
		% ═══════════════════════════════════════════════════════════
		% PANNELLO (b) — Prestazione statica della cerniera interna
		% ═══════════════════════════════════════════════════════════
		\begin{scope}[xshift=8cm]
			
			%% ── Parametri (identici ad (a)) ──
			\pgfmathsetmacro{\rn}{0.15}
			\def\xA{0.00}
			\def\yA{0.00}
			
			%% ── Corpo C_i (identico ad (a)) ──
			\draw[line width=1pt] 
			plot[smooth cycle, tension=0.7] coordinates 
			{(-0.15, 0.50) (-0.30, 0.85) (-0.65, 1.10) (-1.35, 1.40) 
				(-2.15, 1.20) (-2.55, 0.50) (-2.55, -0.50)
				(-2.15, -1.20) (-1.35, -1.40) (-0.65, -1.10) 
				(-0.30, -0.85) (-0.15, -0.50)};
			\node at (-1.55, 0.60) {$\mathcal{C}_i$};
			
			%% ── Corpo C_j (identico ad (a)) ──
			\draw[line width=1pt] 
			plot[smooth cycle, tension=0.7] coordinates 
			{(0.15, 0.50) (0.30, 0.85) (0.65, 1.10) (1.35, 1.40) 
				(2.15, 1.20) (2.55, 0.50) (2.55, -0.50)
				(2.15, -1.20) (1.35, -1.40) (0.65, -1.10) 
				(0.30, -0.85) (0.15, -0.50)};
			\node at (1.55, 0.60) {$\mathcal{C}_j$};
			
			%% ── Pallino della cerniera in A ──
			\draw[line width=0.8pt, fill=white] (\xA, \yA) circle (\rn);
			
			%% ── Etichette A_i, A_j (identiche ad (a)) ──
			\node[font=\footnotesize, anchor=south east] 
			at ({-\rn-0.02}, 0.05) {$A_i$};
			\node[font=\footnotesize, anchor=south west] 
			at ({\rn-0.02}, 0.0) {$A_j$};
			
			%% ──────────────────────────────────────────────────
			%% REAZIONI: 4 frecce, azione e reazione
			%% ──────────────────────────────────────────────────
			
			%% Forze applicate a C_i (lato sinistro): -X(A) a_x, -Y(A) a_y
			% Componente orizzontale: da A verso sinistra
			\draw[->, line width=1.4pt] ({-\rn}, 0) -- ({-\rn-1.05}, 0);
			\node[anchor=north, font=\small] at ({-\rn-1.10}, -0.1) {$-X(A)\,\ab_x$};
			% Componente verticale: da A verso il basso
			\draw[->, line width=1.4pt] (-0.05, {-\rn}) -- (-0.05, {-\rn-1.0});
			\node[anchor=north, font=\small] at (0, {-\rn-1.10}) {$-Y(A)\,\ab_y$};
			
			%% Forze applicate a C_j (lato destro): +X(A) a_x, +Y(A) a_y
			% Componente orizzontale: da A verso destra
			\draw[->, line width=1.4pt] ({\rn}, 0) -- ({\rn+1.05}, 0);
			\node[anchor=north, font=\small] at ({\rn+1.10}, -0.1) {$X(A)\,\ab_x$};
			% Componente verticale: da A verso l'alto
			\draw[->, line width=1.4pt] (0.05, {\rn}) -- (0.05, {\rn+1.0});
			\node[anchor=south, font=\small] at (0, {\rn+1.10}) {$Y(A)\,\ab_y$};
			
			%% ── Terna x, y in basso a sinistra ──
			\draw[->, line width=0.8pt] (-3.20, -1.80) -- (-2.40, -1.80);
			\node[anchor=west, font=\footnotesize] at (-2.40, -1.85) {$x$};
			\draw[->, line width=0.8pt] (-3.20, -1.80) -- (-3.20, -1.00);
			\node[anchor=south, font=\footnotesize] at (-3.20, -1.00) {$y$};
			
			%% ── Etichetta pannello ──
			\node[anchor=north, font=\small] 
			at (0.00, -2.30) {$(b)$};
			
		\end{scope}
		
	\end{tikzpicture}
	\caption[Cerniera interna fra due corpi: prestazioni cinematica e statica.]{%
		Cerniera interna fra i corpi $\mathcal{C}_i$ e
		$\mathcal{C}_j$, applicata nel punto $A$ che si identifica
		con i punti di applicazione $A_i \in \mathcal{C}_i$ e
		$A_j \in \mathcal{C}_j$. Prestazione cinematica (a):
		l'equazione di vincolo $\ub(A_j) - \ub(A_i) = \zerovb$
		impone l'uguaglianza degli spostamenti dei due corpi in
		$A$; le rotazioni dei due corpi restano libere e non sono
		rappresentate. Il \CRr\ coincide con il punto di
		applicazione, $C_{ij} \equiv A_i \equiv A_j$.
		Prestazione statica (b): i due corpi si scambiano una
		forza generica $X(A)\,\ab_x + Y(A)\,\ab_y$ applicata in
		$A$, in dualità con le due componenti di spostamento
		relativo bloccate; la coppia scambiata è identicamente
		nulla, in accordo con la molteplicità due.}
	\label{fig:cerniera_interna}
\end{figure}
%%%%%%%%%%%%%%%%%%%%%%%%%%%%%%%%%%%%%%%%%%%



\subsection{Il glifo}\label{sec:glifo}

\paragraph{Realizzazione meccanica e parametri caratterizzanti.}
Il glifo\index{glifo|textbf}\index{glifo|seealso{vincolo}}\index{giunto prismatico|see{glifo}} è un giunto prismatico: una guida che costringe il
corpo a scorrere nella propria direzione mantenendone
inalterato l'orientamento. Impedisce quindi la traslazione
nella direzione $\ab$ e la rotazione, lasciando libera la
sola traslazione nella direzione perpendicolare $\ab^{\perp}$.
Il vincolo è caratterizzato da:
\begin{enumerate}[label=(\roman*)]
	\item il \emph{punto di applicazione} $A$, in cui la
	guida si connette al corpo;
	\item il \emph{versore dell'asse} $\ab$, ovvero la
	direzione lungo la quale la traslazione è impedita.
\end{enumerate}
A differenza del carrello, in cui $A$ identifica univocamente
il punto la cui traslazione è vincolata, qui la condizione
$\vartheta = 0$ rende lo spostamento $\ub$ uguale in ogni
punto del corpo: la scrittura $\ub(A)\cdot\ab = 0$ è quindi
equivalente a $\ub(P)\cdot\ab = 0$ per qualsiasi $P$, e $A$,
relativamente a questa equazione, agisce solo come etichetta
del giunto. Il punto $A$ recupera invece il proprio ruolo
fisico di punto di applicazione nella forza reattiva
$R(A)\,\ab$, dove la posizione di applicazione concorre
all'equilibrio dei momenti.


\paragraph{Caso esterno.}

\paragraph{Equazione cinematica.}
Le due equazioni scalari del glifo sono:
\begin{equation}
	\label{eq:glifo_cin}
	\ub(A) \cdot \ab = 0, \qquad
	\vartheta = 0.
\end{equation}
La seconda equazione impone che il campo di spostamento sia
una traslazione pura; la prima seleziona, fra tutte le
traslazioni, quelle nella direzione $\ab^{\perp}$.
Lo spostamento ammesso è pertanto:
\begin{equation}\label{eq:glifo_moto_ammesso}
	\ub(P) = \alpha\,\ab^{\perp}, \qquad \alpha \in \mathbb{R},
	\quad \forall\, P \in \mathcal{C}.
\end{equation}

\begin{oss}
	Il glifo può essere interpretato come la sovrapposizione
	di un carrello con asse $\ab$ e di un bipendolo: le due
	equazioni cinematiche sono indipendenti e si sommano.
	Questa prospettiva è sviluppata sistematicamente per
	tutti i vincoli composti nell'osservazione del
	\PARref{sec:incastro}.
\end{oss}

\paragraph{Posizione del centro di rotazione assoluto.}
Da $\vartheta = 0$ il centro è all'infinito. La traslazione
ammessa è nella direzione $\ab^{\perp}$; il centro si trova
quindi all'infinito nella direzione $\ab$, perpendicolare
alla traslazione:
\begin{equation}\label{eq:CR_glifo}
	C \longrightarrow \infty
	\quad \text{nella direzione } \ab.
\end{equation}

\begin{proposizione}[Centro di rotazione --- glifo]
	Il \CRa\ di un corpo vincolato da un glifo con asse $\ab$
	è situato all'infinito nella direzione $\ab$,
	perpendicolare alla direzione di scorrimento $\ab^{\perp}$.
\end{proposizione}


\begin{oss}[Il glifo come cerniera impropria]
	La proposizione precedente ammette una lettura proiettiva
	suggestiva. Il glifo può essere visto come una
	\emph{cerniera il cui centro è il punto improprio della
		direzione} $\ab$: la cinematica è una rotazione attorno
	a tale centro, che nel piano euclideo si manifesta come
	pura traslazione lungo $\ab^{\perp}$ (perpendicolare alla
	direzione del centro). Per distinguerle, in questa
	prospettiva si parla di \emph{cerniera propria} per la
	cerniera ordinaria — con centro $C = A$ al finito — e di
	\emph{cerniera impropria}\index{cerniera!impropria} per il glifo — con centro
	$C = \infty_\ab$. Questa classificazione è specifica del
	caso binario: il ``glifo multiplo'' — $k$ corpi che condividono
	in $A$ la componente di spostamento lungo $\ab$ e la
	rotazione — è una possibilità geometrica, ma non
	corrisponde a un dispositivo strutturale di uso comune;
	per questo nel \PARref{sec:vincoli_multipli} si trattano
	soltanto la cerniera multipla e l'incastro multiplo.
\end{oss}


%%%%%%%%%%%%%%%%%%%%%%%%%%%%%%%%%%%%%%%%%%%
\begin{figure}[htbp]
	\centering
	\begin{tikzpicture}[>=Stealth, line width=0.8pt]
		
		% ═══════════════════════════════════════════════════════════
		% PANNELLO (a) — Prestazione cinematica del glifo esterno
		% ═══════════════════════════════════════════════════════════
		\begin{scope}[xshift=0cm]
			
			%% ── Parametri ──
			\def\xA{0.00}
			\def\yA{0.00}
			\def\Lpiatto{0.60}    % semi-lunghezza piatto
			\def\bw{0.20}         % semi-distanza ruote
			\def\rrota{0.05}      % raggio ruote
			\def\gw{0.65}         % semi-larghezza barra di terra (un po' più del piatto)
			\def\rrota{0.10}     % raggio ruote (tangenti ai due piatti)
			
			%% ── Asse a (retta verticale grigia continua) ──
			\draw[gray!60, line width=0.6pt] 
			(\xA, -0.55) -- (\xA, 2.85);
			
			%% ── Etichetta del CR all'infinito ──
			\node[gray, font=\footnotesize, anchor=south west] 
			at (\xA+0.30, 2.20) {$C = \infty_\ab$};
			
			%% ── Versore a in alto ──
			\draw[->, line width=1pt] 	(\xA, 2.10) -- (\xA, 2.85);
			\node[anchor=west, font=\small] 
			at (\xA+0.05, 2.85) {$\ab$};
			
			%% ── Corpo C (con piatto orizzontale inferiore) ──
			\draw[line width=1pt] 
			(-\Lpiatto, 0.00) 
			.. controls (-1.55, 0.10) and (-1.80, 0.95) 
			.. (-1.70, 1.55)
			.. controls (-0.85, 2.40) and ( 0.35, 2.50) 
			.. ( 0.85, 2.15)
			.. controls ( 1.65, 1.70) and ( 1.80, 0.60) 
			.. (\Lpiatto, 0.00)
			-- cycle;
			\node at (0.50, 1.60) {$\mathcal{C}$};
			
			%% ── Equazioni di vincolo (in alto, accanto all'asse) ──
			\node[anchor=west, font=\small] at (\xA-2.40, 2.55) {$\ub(A)\cdot\ab = 0$};
			\node[anchor=west, font=\small] at (\xA-2.40, 2.10) {$\vartheta = 0$};
			
			
			%% ── Glifo: prima barra (sotto il piatto del corpo) ──
			%\draw (-\gw, -0.05) -- (\gw, -0.05);
			
			%% ── Glifo: due ruote ──
			\draw (-\bw, -0.13) circle (\rrota);
			\draw (\bw, -0.13) circle (\rrota);
			
			%% ── Glifo: due ruote (tangenti ai due piatti) ──
			\draw (-\bw, -0.155) circle (\rrota);
			\draw (\bw, -0.155) circle (\rrota);
			
			%% ── Glifo: seconda barra (al telaio) ──
			\draw (-\gw, -0.26) -- (\gw, -0.26);
			
			%% ── Glifo: tratteggio del suolo ──
			\foreach \x in {-0.50,-0.25,0,0.25,0.50}
			\draw[thin] (\x, -0.26) -- ++(-0.10, -0.13);
			
			%% ── Punto A e label ──
			\fill (\xA, \yA) circle (1.5pt);
			\node[anchor=south west, font=\small] at (\xA, \yA+0.05) {$A$};
			
			
			%% ── Etichetta verticale "asse del glifo" ──
			\node[gray, rotate=90, anchor=south, font=\footnotesize] at (\xA, 1.20) {asse del glifo};
			
			
			%% ── Etichetta pannello ──
			\node[anchor=north, font=\small] 	at (\xA-1.95, -0.25) {$(a)$};
			
		\end{scope}
		
		% ═══════════════════════════════════════════════════════════
		% PANNELLO (b) — Prestazione statica del glifo esterno
		% ═══════════════════════════════════════════════════════════
		\begin{scope}[xshift=5cm]
			
			%% ── Parametri (identici ad (a)) ──
			\def\xA{0.00}
			\def\yA{0.00}
			\def\Lpiatto{0.60}
			\def\bw{0.20}
			\def\rrota{0.10}
			\def\gw{0.65}
			
			%% ── Asse a (retta verticale grigia continua) ──
			\draw[gray!60, line width=0.6pt] 
			(\xA, -0.55) -- (\xA, 2.85);
			
			%% ── Versore a in alto ──
			\draw[->, line width=1pt] 
			(\xA, 2.10) -- (\xA, 2.85);
			\node[anchor=west, font=\small] 
			at (\xA+0.05, 2.85) {$\ab$};
			
			%% ── Corpo C (identico ad (a)) ──
			\draw[line width=1pt] 
			(-\Lpiatto, 0.00) 
			.. controls (-1.55, 0.10) and (-1.80, 0.95) 
			.. (-1.70, 1.55)
			.. controls (-0.85, 2.40) and ( 0.35, 2.50) 
			.. ( 0.85, 2.15)
			.. controls ( 1.65, 1.70) and ( 1.80, 0.60) 
			.. (\Lpiatto, 0.00)
			-- cycle;
			\node at (0.50, 1.60) {$\mathcal{C}$};
			
			%% ── Glifo: due ruote tangenti ──
			\draw (-\bw, -0.155) circle (\rrota);
			\draw (\bw, -0.155) circle (\rrota);
			
			%% ── Glifo: seconda barra (al telaio) ──
			\draw (-\gw, -0.26) -- (\gw, -0.26);
			
			%% ── Glifo: tratteggio del suolo ──
			\foreach \x in {-0.50,-0.25,0,0.25,0.50}
			\draw[thin] (\x, -0.26) -- ++(-0.10, -0.13);
			
			%% ── Punto A e label ──
			\fill (\xA, \yA) circle (1.5pt);
			\node[anchor=south west, font=\small] 	at (\xA+0.1, \yA-0.04) {$A$};
			
			%% ──────────────────────────────────────────────────
			%% REAZIONE: forza R(A) a + coppia μ
			%% ──────────────────────────────────────────────────
			
			%% Forza R(A) a: freccia verticale lungo l'asse, da A verso l'alto
			\draw[->, line width=1.4pt] 
			(\xA, \yA+0.05) -- (\xA, \yA+1.20);
			\node[anchor=west, font=\small] at (\xA+0.08, \yA+1.0) {$R(A)\,\ab$};
			
			%% Coppia reattiva μ(A): semicerchio centrato in A, sopra il punto
			\def\rmu{0.30}
			\draw[->, line width=1.4pt] ({\xA+\rmu}, {\yA+0.30}) arc(0:180:\rmu);
			\node[anchor=west, font=\small] at (\xA-1.30, {\yA+\rmu+0.5}) {$\mu(A)$};
			
			%% ── Etichetta pannello ──
			\node[anchor=north, font=\small] at (\xA-1.95, -0.25) {$(b)$};
			
		\end{scope}
		
	\end{tikzpicture}
	\caption[Glifo esterno con asse $\ab$: prestazioni cinematica e statica.]{%
		Glifo esterno applicato al corpo $\mathcal{C}$ in $A$,
		con asse $\ab$. Prestazione cinematica (a): le equazioni
		di vincolo $\ub(A)\cdot\ab = 0$ e $\vartheta = 0$
		impediscono la traslazione di $A$ lungo $\ab$ e la
		rotazione del corpo; il \CRa\ è all'infinito sulla retta
		dell'asse, $C = \infty_\ab$. Prestazione statica (b):
		la reazione vincolare consta di una forza $R(A)\,\ab$
		diretta lungo l'asse e di una coppia $\mu(A)$, in
		dualità rispettivamente con la traslazione lungo $\ab$
		e con la rotazione bloccate, in accordo con la
		molteplicità due.}
	\label{fig:glifo_esterno}
\end{figure}
%%%%%%%%%%%%%%%%%%%%%%%%%%%%%%%%%%%%%%%%%%%


\paragraph{Reazione vincolare.}
Gli spostamenti compatibili con il glifo sono della forma
$\ub(P) = \alpha\,\ab^{\perp}$, con $\vartheta = 0$.
Il lavoro della reazione su tali spostamenti deve essere nullo:
\begin{equation}
	\rb(A)\cdot(\alpha\,\ab^{\perp}) + \mu(A) \cdot 0 = 0
	\quad \forall\,\alpha
	\quad\Longrightarrow\quad
	\rb(A) \cdot \ab^{\perp} = 0,
\end{equation}
da cui $\rb(A)$ è diretta lungo $\ab$. La coppia $\mu(A)$
non è vincolata dalla condizione di lavoro nullo (essendo
$\vartheta = 0$) e rimane incognita libera. La reazione è
quindi una forza diretta lungo $\ab$ e una coppia di verso
e intensità non noti a priori:
\begin{equation}\label{eq:reazione_glifo}
	\rb(A) = R(A)\,\ab,
	\qquad
	\mu(A) \in \mathbb{R},
	\qquad
	R(A) \in \mathbb{R}.
\end{equation}
Sono noti a priori il punto di applicazione $A$ e la direzione $\ab$ della
forza reattiva. Per la forza, $A$ è il \emph{punto di
	applicazione} fisico — la sua posizione concorre
all'equilibrio dei momenti del corpo; per la coppia, $A$ è
\emph{etichetta} del giunto, secondo la convenzione del
\PARref{sec:richiamo} (la coppia è un vettore libero, il cui
effetto sul corpo è indipendente dal punto in cui viene
applicata). Le due incognite scalari sono $R(A)$, il cui
segno determina verso e intensità della forza, e $\mu(A)$,
il cui segno determina verso e intensità della coppia, in
accordo con la molteplicità due. La \FIGref{fig:glifo_esterno}
illustra le prestazioni cinematica e statica del glifo esterno.





\paragraph{Caso interno.}
Il glifo interno\index{glifo!interno} fra $\mathcal{C}_i$ e $\mathcal{C}_j$ è un
vincolo puntiforme: la condizione di uguaglianza, in $A$,
delle componenti dello spostamento dei due corpi nella
direzione $\ab$ e delle rotazioni si scrive come annullamento
dei salti corrispondenti,
\begin{equation}\label{eq:glifo_int}
	\llbracket \ub \rrbracket(A) \cdot \ab = 0,
	\qquad
	\llbracket \vartheta \rrbracket = 0.
\end{equation}
I due corpi hanno la stessa rotazione e la stessa componente
di spostamento lungo $\ab$ in $A$; possono scorrere l'uno
rispetto all'altro lungo $\ab^{\perp}$ mantenendo lo stesso
orientamento. Il \CRr\ $C_{ij}$ è all'infinito in direzione
$\ab$. Le reazioni scambiate fra i due corpi sono una forza
diretta lungo $\ab$ e una coppia di verso e intensità non
noti a priori, uguali e opposte:
\begin{equation}\label{eq:reazione_glifo_int}
	\rb_i = -R(A)\,\ab,
	\qquad
	\rb_j = R(A)\,\ab,
	\qquad
	R(A), \mu(A) \in \mathbb{R}.
\end{equation}
Sono noti a priori il punto di applicazione $A$ e la
direzione $\ab$ della forza scambiata; le due incognite
scalari sono $R(A)$, il cui segno determina il verso e
l'intensità della forza trasmessa da $\mathcal{C}_i$ a
$\mathcal{C}_j$, e $\mu(A)$, il cui segno determina il
verso e l'intensità della coppia, in accordo con la
molteplicità due. La \FIGref{fig:glifo_interno} illustra
le prestazioni cinematica e statica del glifo interno fra
due corpi.


%%%%%%%%%%%%%%%%%%%%%%%%%%%%%%%%%%%%%%%%%%%
\begin{figure}[htbp]
	\centering
	\begin{tikzpicture}[>=Stealth, line width=0.8pt, scale=0.85]
		
		% ═══════════════════════════════════════════════════════════
		% PANNELLO (a) — Prestazione cinematica del glifo interno
		% ═══════════════════════════════════════════════════════════
		\begin{scope}[xshift=0cm]
			
			%% ── Parametri ──
			\def\xAi{0.00}
			\def\yA{0.00}
			\def\Lpiatto{0.60}    % semi-lunghezza piatto verticale
			\def\bw{0.20}         % semi-distanza ruote
			\pgfmathsetmacro{\rrota}{0.105}
			\pgfmathsetmacro{\D}{2*\rrota}
			\pgfmathsetmacro{\xAj}{\xAi+\D}
			\pgfmathsetmacro{\xMid}{\xAi+\rrota}
			
			%% ── Asse a (retta orizzontale grigia continua) ──
			\draw[gray!60, line width=0.6pt]  (-4.0, \yA) -- (3.10, \yA);
			
			%% ── Versore a a destra ──
			\draw[->, line width=1pt] 	(1.80, \yA) -- (2.30, \yA);
			\node[anchor=north, font=\small] at (1.85, \yA-0.05) {$\ab$};
			
			%% ── Etichetta del CR all'infinito ──
			\node[gray, font=\footnotesize, anchor=south] 	at (-3.80, \yA+0.0) {$C_{ij} = \infty_\ab$};
			
			%% ── Corpo C_i (a sinistra, piatto verticale a x=0) ──
			\draw[line width=1pt] 
			(\xAi, -\Lpiatto)
			.. controls (-0.60, -1.55) and (-1.55, -1.65)
			.. (-2.15, -1.30)
			.. controls (-2.85, -0.70) and (-3.00,  0.50)
			.. (-2.50, 1.20)
			.. controls (-1.80, 1.65) and (-0.70, 1.55)
			.. (\xAi, \Lpiatto)
			-- cycle;
			\node at (-1.65, 0.70) {$\mathcal{C}_i$};
			
			%			%% ── Corpo C_j (a destra, piatto verticale a x=D) ──
			\draw[line width=1pt] 
			(\xAj, \Lpiatto)
			.. controls ({\xAj+0.70}, 1.55) and ({\xAj+1.80}, 1.65)
			.. ({\xAj+2.50}, 1.20)
			.. controls ({\xAj+3.00}, 0.50) and ({\xAj+2.85}, -0.70)
			.. ({\xAj+2.15}, -1.30)
			.. controls ({\xAj+1.55}, -1.65) and ({\xAj+0.60}, -1.55)
			.. (\xAj, -\Lpiatto)
			-- cycle;
			\node at ({\xAj+1.65}, 0.70) {$\mathcal{C}_j$};
			
			%% ── Equazioni di vincolo (in alto) ──
			\node[font=\small] at (\xMid, 2.0) {$\llbracket \ub \rrbracket(A)\cdot\ab = 0, \quad \llbracket \vartheta \rrbracket = 0$};
			
			
			%% ── Glifo: due ruote tangenti ai piatti ──
			\draw (\xMid, \bw) circle (\rrota);
			\draw (\xMid, -\bw) circle (\rrota);
			
			%% ── Punti A_i e A_j ──
			\fill (\xAi, \yA) circle (1.5pt);
			\fill (\xAj, \yA) circle (1.5pt);
			\node[anchor=east, font=\footnotesize] 
			at ({\xAi-0.05}, {\yA-0.18}) {$A_i$};
			\node[anchor=west, font=\footnotesize] 
			at ({\xAj+0.05}, {\yA-0.18}) {$A_j$};
			
			
			%			
			%% ── Etichetta "asse del glifo" sulla retta dell'asse ──
			\node[gray, font=\footnotesize, anchor=north] at (\xMid-4.0, \yA-0.05) {asse del glifo};
			
			%% ── Etichetta pannello ──
			\node[anchor=north, font=\small] 	at (\xMid, -1.50) {$(a)$};
			
		\end{scope}
		
		% ═══════════════════════════════════════════════════════════
		% PANNELLO (b) — Prestazione statica del glifo interno
		% ═══════════════════════════════════════════════════════════
		\begin{scope}[xshift=6.5cm]
			
			%% ── Parametri (identici ad (a)) ──
			\def\xAi{0.00}
			\def\yA{0.00}
			\def\Lpiatto{0.60}
			\def\bw{0.20}
			\pgfmathsetmacro{\rrota}{0.105}
			\pgfmathsetmacro{\D}{2*\rrota}
			\pgfmathsetmacro{\xAj}{\xAi+\D}
			\pgfmathsetmacro{\xMid}{\xAi+\rrota}
			
			%% ── Asse a (retta orizzontale grigia continua) ──
			\draw[gray!60, line width=0.6pt] (-2.0, \yA) -- (3.0, \yA);
			
			%% ── Versore a a destra ──
			\draw[->, line width=1pt] (3.30, \yA) -- (3.80, \yA);
			\node[anchor=south, font=\small] at (3.55, \yA+0.05) {$\ab$};
			
			%% ── Corpo C_i (identico ad (a)) ──
			%% ── Corpo C_i (a sinistra, piatto verticale a x=0) ──
			\draw[line width=1pt] 
			(\xAi, -\Lpiatto)
			.. controls (-0.60, -1.55) and (-1.55, -1.65)
			.. (-2.15, -1.30)
			.. controls (-2.85, -0.70) and (-3.00,  0.50)
			.. (-2.50, 1.20)
			.. controls (-1.80, 1.65) and (-0.70, 1.55)
			.. (\xAi, \Lpiatto)
			-- cycle;
			\node at (-1.75, 0.90) {$\mathcal{C}_i$};
			
			
			%% ── Corpo C_j (identico ad (a)) ──
			\draw[line width=1pt] 
			(\xAj, \Lpiatto)
			.. controls ({\xAj+0.70}, 1.55) and ({\xAj+1.80}, 1.65)
			.. ({\xAj+2.50}, 1.20)
			.. controls ({\xAj+3.00}, 0.50) and ({\xAj+2.85}, -0.70)
			.. ({\xAj+2.15}, -1.30)
			.. controls ({\xAj+1.55}, -1.65) and ({\xAj+0.60}, -1.55)
			.. (\xAj, -\Lpiatto)
			-- cycle;
			\node at ({\xAj+1.75}, 0.90) {$\mathcal{C}_j$};
			
			%% ── Glifo: due ruote tangenti ai piatti ──
			\draw (\xMid, \bw) circle (\rrota);
			\draw (\xMid, -\bw) circle (\rrota);
			
			%% ── Punti A_i e A_j ──
			\fill (\xAi, \yA) circle (1.5pt);
			\fill (\xAj, \yA) circle (1.5pt);
			\node[anchor=east, font=\footnotesize] 	at ({\xAi-0.05}, {\yA-0.28}) {$A_i$};
			\node[anchor=west, font=\footnotesize] 
			at ({\xAj+0.05}, {\yA-0.28}) {$A_j$};
			
			%% ──────────────────────────────────────────────────
			%% REAZIONI: 2 forze + 2 coppie, uguali e opposte
			%% ──────────────────────────────────────────────────
			
			%% Forza in A_i: -R a (verso sinistra, entra in C_i)
			\draw[->, line width=1.4pt] 
			({\xAi-0.05}, \yA) -- ({\xAi-1.05}, \yA);
			\node[anchor=south, font=\small] 
			at ({\xAi-0.55}, {\yA+0.05}) {$-R\,\ab$};
			
			%% Forza in A_j: +R a (verso destra, entra in C_j)
			\draw[->, line width=1.4pt] 
			({\xAj+0.05}, \yA) -- ({\xAj+1.05}, \yA);
			\node[anchor=south, font=\small] 
			at ({\xAj+0.55}, {\yA+0.05}) {$R\,\ab$};
			
			
			%% ── Coppie μ(A): semicerchi verticali, uguali e opposti ──
			\def\rmu{0.35}
			\def\yMu{0.0}              % quota delle coppie, dentro i corpi
			
			% Coppia in C_i: semicerchio orario (concavo verso destra, freccia in basso)
			\draw[->, line width=1.4pt] ({-1.65+\rmu}, \yMu-\rmu) arc(-90:-270:\rmu);
			\node[font=\small] at (-2.0, \yMu-0.5) {$\mu(A)$};
			
			% Coppia in C_j: semicerchio antiorario (concavo verso sinistra, freccia in basso)
			\draw[->, line width=1.4pt] ({\xAj+1.65-\rmu}, \yMu-\rmu) arc(-90:90:\rmu);
			\node[font=\small] at ({\xAj+2.0}, \yMu-0.5) {$\mu(A)$};
			
			%% ── Etichetta pannello ──
			\node[anchor=north, font=\small] 
			at (\xMid, -1.50) {$(b)$};
			
		\end{scope}
		
	\end{tikzpicture}
	\caption[Glifo interno con asse $\ab$ fra due corpi: prestazioni cinematica e statica.]{%
		Glifo interno fra i corpi $\mathcal{C}_i$ e $\mathcal{C}_j$
		con asse $\ab$. Prestazione cinematica (a): le equazioni
		di vincolo $\llbracket \ub \rrbracket(A) \cdot \ab = 0$ e
		$\llbracket \vartheta \rrbracket = 0$ impongono l'uguaglianza
		delle componenti di spostamento dei due corpi lungo $\ab$
		e l'uguaglianza delle rotazioni $\vartheta_i$ e $\vartheta_j$;
		lo scorrimento relativo lungo $\ab^\perp$ resta libero. Il
		\CRr\ è all'infinito sulla retta dell'asse,
		$C_{ij} = \infty_\ab$. Prestazione statica (b): i due corpi
		si scambiano una forza $R\,\ab$ diretta lungo l'asse e una
		coppia $\mu(A)$, uguali e opposte, in dualità rispettivamente
		con la componente di spostamento relativo lungo $\ab$ e con
		la rotazione relativa bloccate, in accordo con la
		molteplicità due.}
	\label{fig:glifo_interno}
\end{figure}
%%%%%%%%%%%%%%%%%%%%%%%%%%%%%%%%%%%%%%%%%%%



%---------------------------------------------------------------------------------------------------------------------------- 
\section{Vincolo triplo: l'incastro}\label{sec:incastro}
%---------------------------------------------------------------------------------------------------------------------------- 

\paragraph{Realizzazione meccanica e parametri caratterizzanti.}
L'incastro\index{incastro|textbf}\index{incastro|seealso{vincolo}}\index{vincolo!triplo} è l'unico vincolo di molteplicità tre: blocca
completamente il corpo impedendone sia le traslazioni che
la rotazione. Come la cerniera, è caratterizzato dal
\emph{punto di applicazione} $A$; a differenza di essa,
vincola anche la rotazione.

\paragraph{Caso esterno.}

\paragraph{Equazioni cinematiche.}
L'incastro in $A$ impone tre equazioni scalari:
\begin{equation}\label{eq:incastro_cin}
	\ub(A) = \zerovb, \qquad \vartheta = 0,
\end{equation}
equivalenti, tramite la~\eqref{eq:FFCI_piano_vett}, a
richiedere lo spostamento nullo di ogni punto del corpo:
\begin{equation}
	\ub(P) = \zerovb \qquad \forall\, P \in \mathcal{C}.
\end{equation}

\begin{oss}
	L'incastro può essere visto come la sovrapposizione di
	una cerniera in $A$ --- che impone $\ub(A) = \zerovb$
	e porta il centro in $A$ --- e di un bipendolo --- che
	impone $\vartheta = 0$ e porta il centro all'infinito.
	La combinazione dei due vincoli sovradetermina il centro:
	non esiste alcuno spostamento rigido non nullo compatibile
	con entrambi simultaneamente.
\end{oss}

\paragraph{Posizione del centro di rotazione.}
Non esiste alcuno spostamento rigido non nullo compatibile
con l'incastro: il concetto di centro di rotazione è privo
di significato. Il corpo è completamente bloccato.

\paragraph{Reazione vincolare.}
Non esistono spostamenti virtuali compatibili non nulli;
qualsiasi sistema di forze e coppie soddisfa banalmente
la condizione di lavoro virtuale nullo. La reazione
dell'incastro è pertanto del tutto arbitraria: una forza
e una coppia applicate in $A$, di direzione, verso e
intensità non noti a priori,
\begin{equation}\label{eq:reazione_incastro}
	\rb(A) = X(A)\,\ab_x + Y(A)\,\ab_y,
	\qquad
	\mu(A) \in \mathbb{R},
	\qquad
	X(A), Y(A) \in \mathbb{R}.
\end{equation}
È noto a priori il solo punto di applicazione $A$; le tre
incognite scalari sono $X(A)$ e $Y(A)$, ciascuna con segno
libero, che sintetizzano direzione, verso e intensità della
forza, e $\mu(A)$, il cui segno determina il verso e
l'intensità della coppia, in accordo con la molteplicità
tre. La \FIGref{fig:incastro} illustra le prestazioni
cinematica e statica dell'incastro esterno.


%%%%%%%%%%%%%%%%%%%%%%%%%%%%%%%%%%%%%%%%%%%
\begin{figure}[htbp]
	\centering
	\begin{tikzpicture}[>=Stealth, line width=0.8pt]
		
		% ═══════════════════════════════════════════════════════════
		% PANNELLO (a) — Prestazione cinematica dell'incastro
		% ═══════════════════════════════════════════════════════════
		\begin{scope}[xshift=0cm]
			
			%% ── Parametri ──
			\def\xA{0.00}
			\def\yA{0.00}
			\def\Lcamp{0.65}      % era 0.50, semi-altezza fascia campita
			
			%% ── Corpo C (ingrandito) ──
			\draw[line width=1pt] 
			(\xA, -\Lcamp)
			.. controls ( 0.65, -1.65) and ( 1.65, -1.75)
			.. ( 2.25, -1.40)
			.. controls ( 3.00, -0.75) and ( 3.10,  0.50)
			.. ( 2.65, 1.25)
			.. controls ( 1.90, 1.75) and ( 0.75, 1.65)
			.. (\xA, \Lcamp)
			-- cycle;
			\node at (1.80, 0.10) {$\mathcal{C}$};
			
			%% ── Campitura (tratteggio obliquo) lungo il bordo sinistro ──
			\foreach \y in {-0.55,-0.40,-0.25,-0.10,0.05,0.20,0.35,0.50}
			\draw[thin] (\xA, \y) -- ++(-0.13, -0.10);
			
			
			%% ── Punto A e label ──
			\fill (\xA, \yA) circle (1.5pt);
			\node[anchor=south west, font=\small] 
			at ({\xA+0.05}, {\yA+0.05}) {$A$};
			
			%% ── Equazioni di vincolo (a sinistra, fuori dal corpo) ──
			\node[font=\small, anchor=east] 
			at (-0.40, 0.50) {$\ub(A) = \zerovb$};
			\node[font=\small, anchor=east] 
			at (-0.40, 0.10) {$\vartheta = 0$};
			
			%% ── Etichetta pannello ──
			\node[anchor=north, font=\small] 
			at (-0.06, -1.0) {$(a)$};
			
		\end{scope}
		% ═══════════════════════════════════════════════════════════
		% PANNELLO (b) — Prestazione statica dell'incastro
		% ═══════════════════════════════════════════════════════════
		\begin{scope}[xshift=5cm]
			
			%% ── Parametri ──
			\def\xA{0.00}
			\def\yA{0.00}
			\def\Lcamp{0.65}      % era 0.50, semi-altezza fascia campita
			
			%% ── Corpo C (ingrandito) ──
			\draw[line width=1pt] 
			(\xA, -\Lcamp)
			.. controls ( 0.65, -1.65) and ( 1.65, -1.75)
			.. ( 2.25, -1.40)
			.. controls ( 3.00, -0.75) and ( 3.10,  0.50)
			.. ( 2.65, 1.25)
			.. controls ( 1.90, 1.75) and ( 0.75, 1.65)
			.. (\xA, \Lcamp)
			-- cycle;
			\node at (1.80, 0.10) {$\mathcal{C}$};
			
			%% ── Campitura (tratteggio obliquo) lungo il bordo sinistro ──
			\foreach \y in {-0.55,-0.40,-0.25,-0.10,0.05,0.20,0.35,0.50}
			\draw[thin] (\xA, \y) -- ++(-0.13, -0.10);
			
			%% ── Punto A e label ──
			\fill (\xA, \yA) circle (1.5pt);
			\node[anchor=south east, font=\small] 
			at ({\xA-0.05}, {\yA+0.05}) {$A$};
			
			%% ──────────────────────────────────────────────────
			%% REAZIONI: X(A), Y(A), μ(A)
			%% ──────────────────────────────────────────────────
			
			%% Forza X(A): da A verso destra
			\draw[->, line width=1.4pt] 
			(\xA+0.05, \yA) -- (\xA+1.05, \yA);
			\node[anchor=north, font=\small] at ({\xA+0.85}, {\yA-0.05}) {$X(A)\,\ab_x$};
			
			%% Forza Y(A): da A verso l'alto
			\draw[->, line width=1.4pt] 
			(\xA, \yA+0.05) -- (\xA, \yA+1.05);
			\node[anchor=south, font=\small] 	at ({\xA-0.05}, {\yA+1.0}) {$Y(A)\,\ab_y$};
			
			%% ── Terna x, y in basso a sinistra ──
			\draw[->, line width=0.8pt] (-1.50, -1.20) -- (-0.70, -1.20);
			\node[anchor=west, font=\footnotesize] at (-0.70, -1.25) {$x$};
			\draw[->, line width=0.8pt] (-1.50, -1.20) -- (-1.50, -0.40);
			\node[anchor=south, font=\footnotesize] at (-1.50, -0.40) {$y$};
			
			%% Coppia μ(A): arco antiorario dentro al corpo
			\def\xMu{1.10}     % era 0.80
			\def\yMu{0.50}     % era 0.80
			
			%% Coppia μ(A): arco di π/2 centrato in A, fra le due frecce di forza
			\def\rmu{0.55}
			\draw[->, line width=1.4pt] 
			({\xA+\rmu}, \yA) arc(0:90:\rmu);
			\node[font=\small] 
			at ({\xA+\rmu*1.7}, {\yA+\rmu*0.707+0.05}) {$\mu(A)$};
			
			
			%% ── Etichetta pannello ──
			\node[anchor=north, font=\small] 	at (-0.06, -1.0) {$(b)$};
			
		\end{scope}
	\end{tikzpicture}
	\caption[Incastro in $A$: prestazioni cinematica e statica.]{%
		Incastro applicato al corpo $\mathcal{C}$ nel punto $A$.
		Prestazione cinematica (a): nessuno spostamento è ammesso, 
		né traslazione né rotazione ($\ub(A) = \zerovb$, 
		$\vartheta = 0$); il corpo è completamente bloccato.
		Prestazione statica (b): la reazione vincolare è una forza
		generica $X(A)\,\ab_x + Y(A)\,\ab_y$ e una coppia $\mu(A)$,
		applicate in $A$, in accordo con la molteplicità tre.}
	\label{fig:incastro}
\end{figure}
%%%%%%%%%%%%%%%%%%%%%%%%%%%%%%%%%%%%%%%%%%%

\paragraph{Caso interno.}
\label{sec:incastro_interno}
L'incastro interno\index{incastro!interno} fra $\mathcal{C}_i$ e $\mathcal{C}_j$ è
un vincolo puntiforme: la condizione di uguaglianza, in $A$,
degli spostamenti e delle rotazioni dei due corpi si scrive
come annullamento di entrambi i salti,
\begin{equation}\label{eq:incastro_int}
	\llbracket \ub \rrbracket(A) = \zerovb,
	\qquad
	\llbracket \vartheta \rrbracket = 0.
\end{equation}
I due corpi si comportano come un unico corpo rigido nel
punto di collegamento: hanno lo stesso spostamento e la
stessa rotazione. Le reazioni scambiate fra i due corpi
sono una forza e una coppia, di direzione, verso e
intensità non noti a priori, uguali e opposte:
\begin{equation}\label{eq:reazione_incastro_int}
	\begin{aligned}
		&\rb_i = -X(A)\,\ab_x - Y(A)\,\ab_y,
		\qquad
		\rb_j = X(A)\,\ab_x + Y(A)\,\ab_y,
		\\
		&X(A), Y(A), \mu(A) \in \mathbb{R}.
	\end{aligned}
\end{equation}
È noto a priori il solo punto di applicazione $A$; le tre
incognite scalari sono $X(A)$ e $Y(A)$, ciascuna con segno
libero, che sintetizzano direzione, verso e intensità della
forza trasmessa da $\mathcal{C}_i$ a $\mathcal{C}_j$, e
$\mu(A)$, il cui segno determina il verso e l'intensità
della coppia, in accordo con la molteplicità tre.


\begin{oss}
	L'incastro interno è la forma di \emph{collegamento}
	del vincolo di rigidità introdotto nel
	\PARref{sec:vincoli_intro}: le tre condizioni
	cinematiche sono le stesse --- uguaglianza di
	spostamento e rotazione --- e le tre reazioni hanno
	la stessa struttura --- forza risultante e momento
	risultante. La differenza è di natura fisica: il
	vincolo di rigidità è una proprietà interna del
	corpo, diffusa sulla sezione di connessione;
	l'incastro interno è un dispositivo di collegamento
	concentrato nel punto $A$, che scambia forze e coppie
	concentrate. In entrambi i casi la struttura duale
	--- tre condizioni cinematiche, tre reazioni
	incognite --- è la stessa.
\end{oss}



%---------------------------------------------------------------------------------------------------------------------------- 
\section{Vincoli composti: giustapposizione di vincoli semplici}
\label{sec:vincoli_composti}
%---------------------------------------------------------------------------------------------------------------------------- 

I vincoli di molteplicità due e tre si possono interpretare 
come la \emph{giustapposizione}\index{vincolo!composto|textbf}\index{giustapposizione di vincoli|see{vincolo (composto)}} di vincoli semplici --- 
carrelli e bipendoli --- le cui equazioni cinematiche si 
sommano e le cui reazioni si combinano vettorialmente. 
Questa prospettiva non è solo mnemonica: chiarisce in modo 
unitario la struttura delle equazioni di vincolo, la 
posizione del centro di rotazione e la forma delle reazioni, 
rendendo immediato il conteggio della molteplicità.

\begin{oss}[Doppia accezione di ``vincolo elementare'']
	I cinque vincoli presentati nel catalogo --- carrello,
	bipendolo, cerniera, glifo, incastro --- sono ``elementari''
	nel senso didattico-tradizionale: costituiscono il
	repertorio di dispositivi di base con cui si descrivono
	le strutture piane. Dal punto di vista \emph{strutturale},
	però, i veri primitivi sono soltanto due --- carrello e
	bipendolo --- come mostra questa sezione: cerniera, glifo
	e incastro si ottengono per giustapposizione. Le due
	accezioni non sono in conflitto: la prima riflette l'uso
	consolidato del lessico tecnico, la seconda la struttura
	algebrica del problema di vincolo. Per evitare ambiguità,
	in questo capitolo si parla di vincoli \emph{semplici}
	(\PARref{sec:semplici}) quando si intende il livello
	atomico --- carrello e bipendolo --- e di vincoli
	\emph{composti} (questa sezione) quando si intende la
	loro giustapposizione; il termine \emph{elementare} resta
	disponibile per l'uso informale, riferito al catalogo dei
	cinque.
\end{oss}


\paragraph{Cinematica: teorema di intersezione.}
Le equazioni di vincolo della giustapposizione sono l'unione 
di quelle dei vincoli componenti; la loro indipendenza 
determina la molteplicità complessiva per somma. La 
posizione del \CRa\ obbedisce al \emph{teorema di 
	intersezione}\index{teorema!di intersezione}: il centro di un vincolo composto è il punto 
comune ai luoghi del centro dei vincoli componenti. La 
\FIGref{fig:teorema_intersezione} illustra i tre casi 
rilevanti, in cui i due luoghi si incontrano al finito 
(cerniera), all'infinito (glifo) o in nessun punto 
(incastro). In particolare:
\begin{itemize}[noitemsep]
	\item la \textbf{cerniera} è la giustapposizione di due
	carrelli con assi $\ab_1$ e $\ab_2$ \emph{non paralleli}.
	Il centro si trova all'intersezione delle due rette per
	$A$ in direzione $\ab_1$ e $\ab_2$, ossia nel punto $A$
	stesso. Il caso particolare $\ab_1 = \ab_x$,
	$\ab_2 = \ab_y$ corrisponde alla decomposizione
	cartesiana, ma la cerniera \emph{non dipende} dalla
	scelta degli assi dei carrelli componenti, purché non
	collineari;
	
	\item il \textbf{glifo} di asse $\ab$ è la giustapposizione
	di un carrello con asse $\ab$ e un bipendolo, oppure 
	equivalentemente di due carrelli con assi \emph{paralleli}
	ad $\ab$. Nel primo caso il carrello vincola il centro
	sulla retta per $A$ in direzione $\ab$, e il bipendolo lo
	spinge sulla retta impropria $r_\infty$; nel secondo caso,
	le due rette parallele si incontrano nel \emph{punto
		improprio} della loro direzione comune. In entrambe le 
	letture il centro è $C = \infty_\ab$;
	
	\item l'\textbf{incastro} è la giustapposizione di una
	cerniera e un bipendolo, ovvero di due carrelli non
	paralleli e un bipendolo. La cerniera porta il centro in
	$A$, il bipendolo lo richiederebbe sulla retta impropria:
	i due luoghi non hanno punti comuni e \emph{non esiste 
		alcun centro di rotazione}, in coerenza col fatto che 
	nessuno spostamento rigido infinitesimo è ammesso.
\end{itemize}

%%%%%%%%%%%%%%%%%%%%%%%%%%%%%%%%%%%%%%%%%%%
\begin{figure}[htbp]
	\centering
	\begin{tikzpicture}[>=Stealth, line width=0.8pt]
		
		% ═══════════════════════════════════════════════════════════
		% PANNELLO (a) — Cerniera equivalente: rette concorrenti
		% ═══════════════════════════════════════════════════════════
		\begin{scope}[xshift=0cm]
			
			%% ── Parametri ──
			\def\xAuno{-0.80}
			\def\yAuno{-0.65}     % quota di A_1 (più bassa per tangenza)
			\def\xAdue{ 0.80}
			\def\yAdue{-0.55}     % quota di A_2
			\def\anguno{70}
			\def\angdue{110}
			\def\Lasse{2.40}
			\def\tw{0.22}
			\def\thR{0.30}
			\def\bw{0.16}
			\def\rrota{2.4pt}
			\def\gw{0.40}
			\def\rn{0.10}
			
			\pgfmathsetmacro{\dxuno}{cos(\anguno)}
			\pgfmathsetmacro{\dyuno}{sin(\anguno)}
			\pgfmathsetmacro{\dxdue}{cos(\angdue)}
			\pgfmathsetmacro{\dydue}{sin(\angdue)}
			
			%% ── Carrello 1 in A_1 ──
			\begin{scope}[shift={(\xAuno, \yAuno)}, rotate=\anguno-90]
				\draw (0, 0) -- (-\tw, -\thR) -- (\tw, -\thR) -- cycle;
				\draw (-\gw, -\thR) -- (\gw, -\thR);
				\draw (-\bw, {-\thR-0.13}) circle (\rrota);
				\draw (\bw, {-\thR-0.13}) circle (\rrota);
				\draw (-\gw, {-\thR-0.26}) -- (\gw, {-\thR-0.26});
				\foreach \x in {-0.30,-0.15,0,0.15,0.30}
				\draw[thin] (\x, {-\thR-0.26}) -- ++(-0.10, -0.13);
			\end{scope}
			
			%% ── Carrello 2 in A_2 ──
			\begin{scope}[shift={(\xAdue, \yAdue)}, rotate=\angdue-90]
				\draw (0, 0) -- (-\tw, -\thR) -- (\tw, -\thR) -- cycle;
				\draw (-\gw, -\thR) -- (\gw, -\thR);
				\draw (-\bw, {-\thR-0.13}) circle (\rrota);
				\draw (\bw, {-\thR-0.13}) circle (\rrota);
				\draw (-\gw, {-\thR-0.26}) -- (\gw, {-\thR-0.26});
				\foreach \x in {-0.30,-0.15,0,0.15,0.30}
				\draw[thin] (\x, {-\thR-0.26}) -- ++(-0.10, -0.13);
			\end{scope}
			
			%% ── Corpo C (blob smooth) ──
			\draw[line width=1pt] 
			plot[smooth cycle, tension=0.7] coordinates 
			{(0.00, -0.50) (0.90, -0.40) (1.40, 0.10) 
				(1.20, 0.80) (0.30, 1.10) (-0.50, 1.00) 
				(-1.20, 0.50) (-1.40, -0.10) (-0.90, -0.50)};
			\node at (0.00, 0.40) {$\mathcal{C}$};
			
			%% ── Anellini in A_1, A_2 (tangenti al bordo) ──
			\draw[line width=0.8pt, fill=white] (\xAuno, \yAuno) circle (\rn);
			\draw[line width=0.8pt, fill=white] (\xAdue, \yAdue) circle (\rn);
			
			%% ── Rette degli assi ──
			\draw[gray!60, line width=0.6pt, dashed] 
			(\xAuno, \yAuno) -- ({\xAuno+\Lasse*\dxuno}, {\yAuno+\Lasse*\dyuno});
			\draw[gray!60, line width=0.6pt, dashed] 
			(\xAdue, \yAdue) -- ({\xAdue+\Lasse*\dxdue}, {\yAdue+\Lasse*\dydue});
			
			%% ── Etichette A_1, A_2 ──
			\node[anchor=east, font=\footnotesize] 
			at ({\xAuno-0.10}, \yAuno) {$A_1$};
			\node[anchor=west, font=\footnotesize] 
			at ({\xAdue+0.10}, \yAdue) {$A_2$};
			
			%% ── Centro C ──
			\fill (0.02, 1.61) circle (1.5pt);
			\node[anchor=south, font=\small] at (0.02, 1.66) {$C$};
			
			%% ── Etichetta pannello ──
			\node[anchor=north, font=\small] at (0, -2.10) {$(a)$};
			
		\end{scope}
		
		% ═══════════════════════════════════════════════════════════
		% PANNELLO (b) — Glifo: rette parallele, C all'infinito
		% ═══════════════════════════════════════════════════════════
		\begin{scope}[xshift=4.5cm]
			
			%% ── Parametri ──
			\def\xAuno{-0.80}
			\def\yAuno{-0.65}
			\def\xAdue{ 0.80}
			\def\yAdue{-0.55}
			\def\anguno{90}      % verticali, paralleli
			\def\angdue{90}
			\def\Lasse{2.40}
			\def\tw{0.22}
			\def\thR{0.30}
			\def\bw{0.16}
			\def\rrota{2.4pt}
			\def\gw{0.40}
			\def\rn{0.10}
			
			\pgfmathsetmacro{\dxuno}{cos(\anguno)}
			\pgfmathsetmacro{\dyuno}{sin(\anguno)}
			\pgfmathsetmacro{\dxdue}{cos(\angdue)}
			\pgfmathsetmacro{\dydue}{sin(\angdue)}
			
			%% ── Carrello 1 in A_1 ──
			\begin{scope}[shift={(\xAuno, \yAuno)}, rotate=\anguno-90]
				\draw (0, 0) -- (-\tw, -\thR) -- (\tw, -\thR) -- cycle;
				\draw (-\gw, -\thR) -- (\gw, -\thR);
				\draw (-\bw, {-\thR-0.13}) circle (\rrota);
				\draw (\bw, {-\thR-0.13}) circle (\rrota);
				\draw (-\gw, {-\thR-0.26}) -- (\gw, {-\thR-0.26});
				\foreach \x in {-0.30,-0.15,0,0.15,0.30}
				\draw[thin] (\x, {-\thR-0.26}) -- ++(-0.10, -0.13);
			\end{scope}
			
			%% ── Carrello 2 in A_2 ──
			\begin{scope}[shift={(\xAdue, \yAdue)}, rotate=\angdue-90]
				\draw (0, 0) -- (-\tw, -\thR) -- (\tw, -\thR) -- cycle;
				\draw (-\gw, -\thR) -- (\gw, -\thR);
				\draw (-\bw, {-\thR-0.13}) circle (\rrota);
				\draw (\bw, {-\thR-0.13}) circle (\rrota);
				\draw (-\gw, {-\thR-0.26}) -- (\gw, {-\thR-0.26});
				\foreach \x in {-0.30,-0.15,0,0.15,0.30}
				\draw[thin] (\x, {-\thR-0.26}) -- ++(-0.10, -0.13);
			\end{scope}
			
			%% ── Corpo C (identico a (a)) ──
			\draw[line width=1pt] 
			plot[smooth cycle, tension=0.7] coordinates 
			{(0.00, -0.50) (0.90, -0.40) (1.40, 0.10) 
				(1.20, 0.80) (0.30, 1.10) (-0.50, 1.00) 
				(-1.20, 0.50) (-1.40, -0.10) (-0.90, -0.50)};
			\node at (0.00, 0.40) {$\mathcal{C}$};
			
			%% ── Anellini ──
			\draw[line width=0.8pt, fill=white] (\xAuno, \yAuno) circle (\rn);
			\draw[line width=0.8pt, fill=white] (\xAdue, \yAdue) circle (\rn);
			
			%% ── Rette parallele degli assi ──
			\draw[gray!60, line width=0.6pt, dashed] 
			(\xAuno, \yAuno) -- ({\xAuno+\Lasse*\dxuno}, {\yAuno+\Lasse*\dyuno});
			\draw[gray!60, line width=0.6pt, dashed] 
			(\xAdue, \yAdue) -- ({\xAdue+\Lasse*\dxdue}, {\yAdue+\Lasse*\dydue});
			
			%% ── Etichette A_1, A_2 ──
			\node[anchor=east, font=\footnotesize] 
			at ({\xAuno-0.10}, \yAuno) {$A_1$};
			\node[anchor=west, font=\footnotesize] 
			at ({\xAdue+0.10}, \yAdue) {$A_2$};
			
			%% ── Etichetta del centro all'infinito ──
			\node[font=\small] at (0, 2.10) {$C = \infty_\ab$};
			
			%% ── Etichetta pannello ──
			\node[anchor=north, font=\small] at (0, -2.10) {$(b)$};
			
		\end{scope}
		
		% ═══════════════════════════════════════════════════════════
		% PANNELLO (c) — Incastro: cerniera + bipendolo (incompatibilità)
		% ═══════════════════════════════════════════════════════════
		\begin{scope}[xshift=9cm]
			
			%% ── Parametri ──
			\def\xAuno{-1.45}
			\def\yAuno{0.30}        % cerniera tangente al bordo sinistro
			\def\tw{0.22}
			\def\thR{0.30}
			\def\gw{0.40}
			\def\rn{0.10}
			
			%% ── Corpo C (identico ai pannelli a, b) ──
			\draw[line width=1pt] 
			plot[smooth cycle, tension=0.7] coordinates 
			{(0.00, -0.50) (0.90, -0.40) (1.40, 0.10) 
				(1.20, 0.80) (0.30, 1.10) (-0.50, 1.00) 
				(-1.20, 0.50) (-1.40, -0.10) (-0.90, -0.50)};
			\node at (0.10, 0.30) {$\mathcal{C}$};
			
			
			%% ── Cerniera in A_1 (sul bordo sinistro, geometria esplicita) ──
			%% Triangolo: apice in A_1, base verticale a sinistra
			\draw[line width=0.8pt] 
			(-1.45, 0.30) -- (-1.75, 0.05) -- (-1.75, 0.55) -- cycle;
			%% Barra di terra (verticale, sinistra del triangolo)
			\draw[line width=0.8pt] (-1.75, -0.05) -- (-1.75, 0.65);
			%% Tratteggio del suolo (a sinistra della barra)
			\foreach \y in {-0.05, 0.05, 0.15, 0.25, 0.35, 0.45, 0.55, 0.65}
			\draw[thin] (-1.75, \y) -- ++(-0.13, -0.10);
			%% Anellino in A_1
			\draw[line width=0.8pt, fill=white] (-1.45, 0.30) circle (\rn);
			%% Etichetta A_1 (sotto)
			\node[anchor=north, font=\footnotesize] at (-1.15, 0.30) {$A_1$};
			
			
			%% ── Pantografo (compatto, a destra del corpo) ──
			\def\xBuno{1.50}      % poco oltre il bordo destro del corpo
			\def\yBuno{0.50}
			\def\xBdue{1.50}
			\def\yBdue{0.05}      % più ravvicinato in verticale
			\def\Lbiella{0.35}    % bielle più corte
			
			\pgfmathsetmacro{\xMuno}{\xBuno+\Lbiella}
			\pgfmathsetmacro{\yMuno}{\yBuno-0.18}
			\pgfmathsetmacro{\xMdue}{\xBdue+\Lbiella}
			\pgfmathsetmacro{\yMdue}{\yBdue-0.18}
			\pgfmathsetmacro{\xGuno}{\xMuno+\Lbiella}
			\pgfmathsetmacro{\yGuno}{\yMuno+0.18}
			\pgfmathsetmacro{\xGdue}{\xMdue+\Lbiella}
			\pgfmathsetmacro{\yGdue}{\yMdue+0.18}
			
			
			%% ── Bielle del pantografo ──
			\draw[line width=0.8pt] (\xBuno, \yBuno) -- (\xMuno, \yMuno);
			\draw[line width=0.8pt] (\xBdue, \yBdue) -- (\xMdue, \yMdue);
			\draw[line width=0.8pt] (\xMuno, \yMuno) -- (\xMdue, \yMdue);
			\draw[line width=0.8pt] (\xMuno, \yMuno) -- (\xGuno, \yGuno);
			\draw[line width=0.8pt] (\xMdue, \yMdue) -- (\xGdue, \yGdue);
			
			%% ── Anellini ──
			\draw[line width=0.8pt, fill=white] (\xBuno, \yBuno) circle (\rn);
			\draw[line width=0.8pt, fill=white] (\xBdue, \yBdue) circle (\rn);
			\draw[line width=0.8pt, fill=white] (\xMuno, \yMuno) circle (3pt);
			\draw[line width=0.8pt, fill=white] (\xMdue, \yMdue) circle (3pt);
			\draw[line width=0.8pt, fill=white] (\xGuno, \yGuno) circle (3pt);
			\draw[line width=0.8pt, fill=white] (\xGdue, \yGdue) circle (3pt);
			
			%% ── Parete del telaio (a destra) ──
			\draw[line width=0.8pt] 
			({\xGuno+0.10}, {\yGuno+0.15}) -- ({\xGdue+0.10}, {\yGdue-0.15});
			\foreach \y in {-0.20,0,0.20,0.40,0.60}
			\draw[thin] ({\xGdue+0.10}, \y) -- ++(0.13, -0.10);
			
			%% ── Annotazione su due righe (centrata, sotto il corpo) ──
			\node[font=\footnotesize, anchor=north] at (0, -1.10) {cerniera: $C \in A_1$, \quad bipendolo: $C \in r_\infty$};
			\node[font=\small, anchor=north] 	at (0, -1.45) {$\Rightarrow \quad \nexists\, C$};
			
			%% ── Etichetta pannello ──
			\node[anchor=north, font=\small] at (0, -2.05) {$(c)$};
			
		\end{scope}
		
	\end{tikzpicture}
	\caption[Teorema di intersezione.]{%
		Il \CRa\ di un vincolo composto come intersezione dei luoghi 
		del centro dei vincoli componenti: (a) cerniera, (b) glifo, 
		(c) incastro.}
	\label{fig:teorema_intersezione}
\end{figure}
%%%%%%%%%%%%%%%%%%%%%%%%%%%%%%%%%%%%%%%%%%%

\paragraph{Statica: somma vettoriale delle reazioni.}
Le reazioni di un vincolo composto sono la somma delle 
reazioni dei componenti. Ogni carrello contribuisce con 
una forza scalare diretta lungo il proprio asse, ogni 
bipendolo con una coppia: un vincolo a molteplicità $m$ 
ha quindi esattamente $m$ incognite scalari, in dualità 
con le $m$ equazioni cinematiche.
\begin{itemize}[noitemsep]
	\item \textbf{cerniera}: due carrelli con assi $\ab_1, \ab_2$
	contribuiscono con le forze scalari $R_1\,\ab_1$ e
	$R_2\,\ab_2$. La reazione complessiva è la somma vettoriale
	$\rb(A) = R_1\,\ab_1 + R_2\,\ab_2$. Per $\ab_1, \ab_2$
	ortogonali si ottiene la decomposizione cartesiana
	$\rb(A) = X(A)\,\ab_x + Y(A)\,\ab_y$ usata in
	\FIGref{fig:cerniera_esterna};
	
	\item \textbf{glifo}: il carrello componente contribuisce
	con una forza $R(A)\,\ab$ in $A$, il bipendolo con la
	coppia $\mu(A)$. La reazione complessiva è la sovrapposizione
	dei due contributi e coincide con quella riportata in
	\FIGref{fig:glifo_esterno};
	
	\item \textbf{incastro}: la cerniera componente contribuisce
	con la forza vettoriale $\rb(A) = X(A)\,\ab_x + Y(A)\,\ab_y$,
	il bipendolo componente con la coppia $\mu(A)$. Le tre
	incognite scalari $X(A), Y(A), \mu(A)$ sono in dualità con
	le tre componenti scalari della cinematica
	$u(A), v(A), \vartheta$: l'incastro chiude la sequenza dei
	vincoli del piano, esaurendo i tre gradi di libertà del
	corpo rigido.
\end{itemize}

Per il glifo, in particolare, la decomposizione in carrello e 
bipendolo ammette una lettura più ricca, basata sulla 
decomposizione modale delle reazioni e sulla loro dualità con 
le componenti del cedimento vincolare; questo sviluppo è 
trattato nel \PARref{sec:dualita_composti}, dopo
l'introduzione del cedimento.

\paragraph{Svincoli.}
La decomposizione rende immediato il significato degli 
\emph{svincoli}: introdurre uno svincolo equivale a rimuovere 
uno dei vincoli semplici componenti. Da un incastro, la 
rimozione del bipendolo libera la rotazione e il vincolo si 
riduce a una cerniera; la rimozione di un carrello da una 
cerniera lascia libero uno scorrimento e il vincolo si riduce 
al carrello restante. Lo stesso vale per i vincoli interni e 
sarà sviluppato sistematicamente per la trave rigida nel \PARref{sec:svincoli_interni}.



%---------------------------------------------------------------------------------------------------------------------------- 
\section{Tavola riepilogativa}\label{sec:tavola}
%---------------------------------------------------------------------------------------------------------------------------- 

Il filo conduttore comune a tutti i vincoli esaminati è il
\emph{principio dei lavori virtuali}: la reazione vincolare
non compie lavoro su alcuno spostamento virtuale compatibile
con il vincolo. Questo principio stabilisce una corrispondenza
biunivoca fra le equazioni cinematiche imposte dal vincolo e
le incognite statiche che esso introduce. Ad ogni componente
di spostamento vincolata corrisponde una e una sola componente
reattiva incognita: la forza reattiva $\rb(A)$, applicata nel
punto $A$, è duale dello spostamento $\ub(A)$ vincolato; la
coppia reattiva $\mu(A)$, applicata nello stesso punto, è
duale della rotazione $\vartheta$ vincolata.

Le componenti della forza reattiva si indicano con $X(A)$ e
$Y(A)$, dove il punto di applicazione è esplicitato fra
parentesi e le lettere identificano le direzioni $\ab_x$ e
$\ab_y$ rispettivamente; per i vincoli che esplicano una sola
componente scalare nella direzione $\ab$ si usa $R(A)$.\nomenclature[L]{$R(A)$}{intensità scalare della reazione vincolare lungo l'asse del vincolo}\nomenclature[L]{$X(A),\ Y(A)$}{componenti cartesiane della reazione vincolare}
Le coppie attive (assegnate) si indicano con $\mu_a(A)$,
quelle reattive (incognite) con $\mu_r(A)$.\nomenclature[G]{$\mu_a(A),\ \mu_r(A)$}{coppia attiva (assegnata) e coppia reattiva (incognita)}

\begin{nota}
	Negli Esempi, per alleggerire la notazione, il punto
	di applicazione viene scritto come pedice anziché tra
	parentesi: $X_A$, $Y_A$, $R_A$, $\mu_{A}$ in luogo
	di $X(A)$, $Y(A)$, $R(A)$, $\mu(A)$. Le due notazioni
	sono equivalenti.
\end{nota}

\medskip

\noindent
La tavola seguente riassume, per ciascun vincolo esterno, le
equazioni cinematiche, il luogo del centro di rotazione
assoluto e le reazioni vincolari.

\medskip

\begin{table}[H]
	\centering
	\resizebox{\textwidth}{!}{%
		\renewcommand{\arraystretch}{1.8}
		\small
		\begin{tabular}{>{\bfseries}lcccc}
			\toprule
			Vincolo & Molt.
			& Eq.\ cinematica
			& Centro $C$
			& Reazione\\
			\midrule
			Carrello & 1
			& $\ub(A)\cdot\ab=0$
			& retta per $A$ in direzione $\ab$
			& $R(A)\,\ab$\\
			Bipendolo & 1
			& $\vartheta=0$
			& $\infty$
			& $\mu(A)$\\
			Cerniera & 2
			& $\ub(A)=\zerovb$
			& $C=A$
			& $X(A)\,\ab_x+Y(A)\,\ab_y$\\
			Glifo & 2
			& $\ub(A)\cdot\ab=0,\;\vartheta=0$
			& $\infty$ in direzione $\ab$
			& $R(A)\,\ab,\,\,\mu(A)$\\
			Incastro & 3
			& $\ub(A)=\zerovb,\;\vartheta=0$
			& nessuno
			& $X(A)\,\ab_x+Y(A)\,\ab_y,\,\,\mu(A)$\\
			\bottomrule
		\end{tabular}%
	}
	\caption[Tavola riepilogativa dei vincoli piani esterni]{%
		Prestazioni cinematiche e statiche dei vincoli piani
		esterni: molteplicità, equazioni cinematiche, centro
		di rotazione assoluto e reazioni vincolari.}
	\label{tab:vincoli_esterni}
\end{table}


\medskip

\noindent
\medskip

\noindent
La tavola analoga per i vincoli \textbf{interni} si ottiene
sostituendo, per i vincoli puntiformi:
\begin{itemize}[noitemsep]
	\item $\ub(A)$ con il salto
	$\llbracket \ub \rrbracket(A) = \ub_j(A) - \ub_i(A)$;
	\item $\vartheta$ con il salto
	$\llbracket \vartheta \rrbracket = \vartheta_j - \vartheta_i$;
	\item il centro assoluto $C$ con il centro relativo
	$C_{ij}$;
	\item le reazioni con coppie di forze e momenti uguali
	e opposti scambiati fra i due corpi nel punto $A$.
\end{itemize}
Per la biella di lunghezza finita, unico vincolo interno
non puntiforme, l'equazione cinematica
$\Delta\ub \cdot \ab = 0$ usa l'incremento fra i due punti
di applicazione $A_i$ e $A_j$ anziché il salto, e le
reazioni sono coppie di forze uguali e opposte applicate
in $A_i$ e $A_j$ separatamente.


%---------------------------------------------------------------------------------------------------------------------------- 
\section{Il cedimento vincolare}\label{sec:cedimento}
%---------------------------------------------------------------------------------------------------------------------------- 

\begin{definizione}[Cedimento vincolare]\label{def:cedimento}
	Si chiama \emph{cedimento vincolare}\index{cedimento!vincolare|textbf}\index{cedimento|seealso{distorsione}} la prescrizione di
	uno spostamento assegnato --- anziché nullo --- nella
	direzione vincolata. Il cedimento è una condizione
	cinematica attiva: come il vincolo perfetto, prescrive
	una componente di spostamento; a differenza di esso,
	la prescrive a un valore $\bar{s} \neq 0$\nomenclature[L]{$\bar{s}$}{cedimento: spostamento (o scorrimento relativo) assegnato lungo la direzione vincolata}\nomenclature[G]{$\bar{\vartheta}$}{cedimento rotazionale: rotazione (relativa) assegnata} assegnato
	dall'esterno.
\end{definizione}

\begin{oss}
	Il cedimento vincolare si distingue sia dal vincolo
	perfetto (spostamento nullo) sia dal carico esterno
	(forza assegnata). È una prescrizione cinematica, non
	dinamica: non introduce alcuna forza nota nel sistema,
	ma modifica la configurazione ammissibile del corpo.
	Nella pratica ingegneristica i cedimenti vincolari
	modellano, ad esempio, gli abbassamenti dei supporti
	di una trave per effetto del consolidamento del terreno
	o delle deformazioni elastiche dei dispositivi di
	appoggio.
\end{oss}

\paragraph{Generalizzazione delle equazioni cinematiche.}
Per ciascun vincolo, l'introduzione del cedimento si traduce
nella sostituzione del valore zero nel membro destro
dell'equazione cinematica con il valore assegnato $\bar{s}$.
La struttura delle equazioni rimane invariata:

\medskip

\begin{center}
	\renewcommand{\arraystretch}{1.6}
	\small
	\begin{tabular}{>{\bfseries}lll}
		\toprule
		Vincolo
		& Eq.\ cinematica perfetta
		& Eq.\ cinematica con cedimento\\
		\midrule
		Carrello
		& $\ub(A)\cdot\ab = 0$
		& $\ub(A)\cdot\ab = \bar{s}(A)$\\
		Bipendolo
		& $\vartheta = 0$
		& $\vartheta = \bar{\vartheta}(A)$\\
		Cerniera
		& $\ub(A) = \zerovb$
		& $\ub(A) = \bar{u}(A)\,\ab_x
		+ \bar{v}(A)\,\ab_y$\\
		Glifo
		& $\ub(A)\cdot\ab = 0,\;\vartheta=0$
		& $\ub(A)\cdot\ab = \bar{s}(A),\;
		\vartheta = \bar{\vartheta}(A)$\\
		Incastro
		& $\ub(A)=\zerovb,\;\vartheta=0$
		& $\ub(A) = \bar{u}(A)\,\ab_x
		+ \bar{v}(A)\,\ab_y,\;
		\vartheta = \bar{\vartheta}(A)$\\
		\bottomrule
	\end{tabular}
\end{center}

\smallskip

\paragraph{Effetto sul centro di rotazione.}
Il cedimento modifica la posizione del centro di rotazione
ma ne preserva la struttura geometrica. Riprendiamo il caso
del carrello: con cedimento $\bar{s}(A)$, la condizione
$\ub(A)\cdot\ab = \bar{s}(A)$ sostituita
nella~\eqref{eq:centro_campo} dà
\begin{equation}
	\bigl(\bm{\upvartheta}\times\overrightarrow{CA}\bigr)
	\cdot\ab = \bar{s}(A),
\end{equation}
da cui, con la proprietà ciclica del prodotto misto:
\begin{equation}
	\vartheta\,\bigl(\overrightarrow{CA}\times\ab\bigr)
	\cdot\ab_z = \bar{s}(A).
\end{equation}
Il centro non giace più sulla retta per $A$ in direzione
$\ab$, ma su una retta parallela, traslata di una quantità
$\bar{s}(A)/\vartheta$ nella direzione perpendicolare
$\ab^{\perp}$:
\begin{equation}\label{eq:CR_carrello_cedimento}
	C = A + \frac{\bar{s}(A)}{\vartheta}\,\ab^{\perp}
	+ \lambda\,\ab,
	\qquad \lambda \in \mathbb{R}.
\end{equation}
Per $\bar{s}(A) = 0$ si ritrova il
risultato~\eqref{eq:CR_carrello}. Lo stesso ragionamento
si applica agli altri vincoli: il cedimento trasla il luogo
del centro senza modificarne la direzione.


\paragraph{Effetto sulla reazione vincolare.}
La struttura della reazione vincolare non è modificata dal
cedimento. Il principio dei lavori virtuali impone che la
reazione compia lavoro nullo sugli \emph{spostamenti
	virtuali} compatibili con il vincolo, non sullo spostamento
assoluto: gli spostamenti virtuali sono variazioni
infinitesime attorno alla configurazione corrente, e la
presenza del cedimento non ne modifica la struttura. Quindi:

\begin{center}
	\renewcommand{\arraystretch}{1.6}
	\small
	\begin{tabular}{>{\bfseries}ll}
		\toprule
		Vincolo & Reazione (con o senza cedimento)\\
		\midrule
		Carrello  & $\rb(A) = R(A)\,\ab$\\
		Bipendolo & $\rb = \zerovb,\quad \mu$\\
		Cerniera  & $\rb(A) = X(A)\,\ab_x + Y(A)\,\ab_y$\\
		Glifo     & $\rb(A) = R(A)\,\ab,\quad \mu$\\
		Incastro  & $\rb(A) = X(A)\,\ab_x + Y(A)\,\ab_y,
		\quad \mu$\\
		\bottomrule
	\end{tabular}
\end{center}

\begin{oss}
	Il fatto che la reazione non dipenda dal cedimento è
	un risultato non banale: il cedimento modifica la
	\emph{cinematica} del sistema (posizione del centro
	di rotazione, spostamenti dei punti) ma non la
	\emph{struttura delle incognite statiche}. Le
	componenti della reazione rimangono le stesse
	grandezze incognite del caso perfetto; sono i loro
	valori numerici a cambiare, non la loro natura.
	Questo principio è alla base del metodo degli
	spostamenti imposti nell'analisi strutturale.
\end{oss}

\paragraph{Cedimento nei vincoli interni.}
Per i vincoli interni il cedimento ha un significato
analogo: anziché prescrivere uno spostamento relativo
nullo, si prescrive uno spostamento relativo assegnato.
Le equazioni cinematiche si generalizzano sostituendo
le grandezze cinematiche relative --- il salto
$\llbracket\ub\rrbracket(A)$ e $\llbracket\vartheta\rrbracket$
per i vincoli puntiformi, l'incremento $\Delta\ub$ per la
biella di lunghezza finita --- con i valori assegnati
corrispondenti:

\medskip

\begin{center}
	\renewcommand{\arraystretch}{1.6}
	\small
	\begin{tabular}{>{\bfseries}lll}
		\toprule
		Vincolo
		& Eq.\ interna perfetta
		& Eq.\ interna con cedimento\\
		\midrule
		Biella
		& $\Delta\ub\cdot\ab = 0$
		& $\Delta\ub\cdot\ab = \bar{s}$\\
		Carrello
		& $\llbracket\ub\rrbracket(A)\cdot\ab = 0$
		& $\llbracket\ub\rrbracket(A)\cdot\ab = \bar{s}(A)$\\
		Bipendolo
		& $\llbracket\vartheta\rrbracket = 0$
		& $\llbracket\vartheta\rrbracket = \bar{\vartheta}(A)$\\
		Cerniera
		& $\llbracket\ub\rrbracket(A) = \zerovb$
		& $\llbracket\ub\rrbracket(A) = \bar{u}(A)\,\ab_x
		+ \bar{v}(A)\,\ab_y$\\
		Glifo
		& $\llbracket\ub\rrbracket(A)\cdot\ab = 0,\;
		\llbracket\vartheta\rrbracket = 0$
		& $\llbracket\ub\rrbracket(A)\cdot\ab = \bar{s}(A),\;
		\llbracket\vartheta\rrbracket = \bar{\vartheta}(A)$\\
		Incastro
		& $\llbracket\ub\rrbracket(A) = \zerovb,\;
		\llbracket\vartheta\rrbracket = 0$
		& $\llbracket\ub\rrbracket(A) = \bar{u}(A)\,\ab_x
		+ \bar{v}(A)\,\ab_y,\;
		\llbracket\vartheta\rrbracket = \bar{\vartheta}(A)$\\
		\bottomrule
	\end{tabular}
\end{center}


\medskip

\noindent
Il cedimento relativo $\bar{s}(A)$ rappresenta lo
scorrimento imposto fra i due corpi nella direzione
vincolata; $\bar{\vartheta}$ rappresenta la rotazione
relativa imposta. Come nel caso esterno, la struttura
delle reazioni scambiate fra i due corpi non è
modificata dalla presenza del cedimento: cambiano i
valori delle incognite, non la loro natura.

\begin{oss}
	Un cedimento relativo in una cerniera interna
	($\llbracket\ub\rrbracket(A) \neq \zerovb$) modella, ad esempio,
	il gioco meccanico in un giunto: i due corpi non
	sono più esattamente coincidenti nel punto $A$, ma
	sono spostati l'uno rispetto all'altro di una
	quantità nota e assegnata. Analogamente, un cedimento
	relativo in un bipendolo interno
	($\llbracket\vartheta\rrbracket \neq 0$) modella una rotazione
	relativa imposta fra i due corpi, come quella
	prodotta da un attuatore rotazionale interposto nel
	giunto.
\end{oss}
% %

\paragraph{Distorsioni del vincolo di rigidità.}
Il vincolo di rigidità ammette anch'esso condizioni
non nulle, dette \emph{distorsioni} (o distorsioni di
Volterra): una o più delle tre condizioni cinematiche
di rigidità --- uguaglianza delle due componenti di
spostamento e della rotazione --- sono sostituite da
valori assegnati non nulli. Come per i cedimenti dei
vincoli di collegamento, la struttura delle reazioni
interne resta invariata: cambiano i valori delle
incognite, non la loro natura.
La trattazione esplicita delle distorsioni --- con la
loro espressione nel riferimento intrinseco al corpo e
la distinzione fra distorsioni longitudinali, trasversali
e rotazionali --- è sviluppata nei capitoli dedicati
alla trave, sia rigida (cfr.~\CAPref{cap:Trave_Rigida}) che deformabile (Vol.~II.).


\section{Dualità fra vincoli composti e componenti}
\label{sec:dualita_composti}

Riprendiamo la decomposizione dei vincoli composti
introdotta nel \PARref{sec:vincoli_composti}, ora che
disponiamo del concetto di cedimento. Ai cedimenti dei
vincoli componenti corrispondono altrettante incognite
statiche; le componenti del moto rigido residuo sono in
dualità con le componenti delle reazioni. Il \PLV,
richiedendo lavoro nullo sui moti compatibili, fornisce
la verifica diretta della dualità. Per cerniera e incastro
la verifica è immediata; per il glifo, invece, la
decomposizione produce una struttura modale (sym/skw) di
particolare interesse, che è il principale contenuto di
questa sezione.

\paragraph{Cerniera.}
Detti $s_1, s_2$ i cedimenti dei due carrelli componenti e
$R_1, R_2$ le rispettive intensità, il \PLV\ fornisce
\begin{equation}
	R_1\,s_1 + R_2\,s_2 = \rb(A) \cdot \ub(A),
\end{equation}
con $\rb(A) = R_1\,\ab_1 + R_2\,\ab_2$ e $\ub(A)$ lo
spostamento di $A$. La dualità è perfetta: per
$\ab_1, \ab_2$ ortogonali si recupera la decomposizione
cartesiana $\rb(A) = X(A)\,\ab_x + Y(A)\,\ab_y$ duale a
$\ub(A) = u(A)\,\ab_x + v(A)\,\ab_y$.

\paragraph{Glifo: decomposizione modale sym/skw.}
Letto come giustapposizione di due carrelli con assi
paralleli ad $\ab$ a distanza $d$, applicati in punti
$A_1, A_2$, il glifo trasmette al corpo le forze
$R_1\,\ab, R_2\,\ab$ duali ai cedimenti $s_1, s_2$.

Le grandezze \emph{naturali} per il vincolo composto non
sono $R_1, R_2$ né $s_1, s_2$ separatamente, ma le loro
combinazioni \emph{simmetrica} e \emph{antisimmetrica}:
\begin{equation}\label{eq:sym_skw_R}
	R_\text{sym} = \frac{R_1 + R_2}{2},
	\qquad
	R_\text{skw} = \frac{R_2 - R_1}{2},
\end{equation}
\begin{equation}\label{eq:sym_skw_s}
	s_\text{sym} = \frac{s_1 + s_2}{2},
	\qquad
	\vartheta = \frac{s_2 - s_1}{d}.
\end{equation}
Le forze del glifo composto si scrivono in termini di
queste combinazioni come
\begin{equation}\label{eq:glifo_R_mu}
	R(A) = 2\,R_\text{sym},
	\qquad
	\mu(A) = d\,R_\text{skw},
\end{equation}
dove $R(A)$ è la forza risultante (lungo $\ab$) e $\mu(A)$
è la coppia di trasporto. Dualmente, $s_\text{sym}$ è la
traslazione comune dei due punti di applicazione lungo
$\ab$ ed è la grandezza cinematica vincolata dal carrello
componente; $\vartheta$ è la rotazione del corpo ed è la
grandezza vincolata dal bipendolo componente. La
conservazione del lavoro virtuale si verifica direttamente:
\begin{equation}\label{eq:lavoro_glifo}
	R_1\,s_1 + R_2\,s_2
	= R(A)\,s_\text{sym} + \mu(A)\,\vartheta,
\end{equation}
mostrando che $R(A)$ è duale alla traslazione comune
$s_\text{sym}$ (equazione cinematica del carrello
componente) e $\mu(A)$ è duale alla rotazione $\vartheta$
(equazione cinematica del bipendolo componente).

A differenza della cerniera --- la cui decomposizione
duale produce una forza vettoriale in $A$ --- nel glifo
essa produce una forza e una coppia. Il motivo è
geometrico: gli assi dei due componenti sono
\emph{paralleli} anziché concorrenti, e i due modi
cinematici indipendenti del vincolo includono una
rotazione anziché due traslazioni.

\paragraph{Incastro.}
La cerniera componente contribuisce con la forza
vettoriale $\rb(A) = X(A)\,\ab_x + Y(A)\,\ab_y$, duale
allo spostamento $\ub(A) = u(A)\,\ab_x + v(A)\,\ab_y$;
il bipendolo componente con la coppia $\mu(A)$, duale
alla rotazione $\vartheta$. Il \PLV\ assume la forma
\begin{equation}\label{eq:lavoro_incastro}
	\rb(A) \cdot \ub(A) + \mu(A)\,\vartheta
	= X(A)\,u(A) + Y(A)\,v(A) + \mu(A)\,\vartheta,
\end{equation}
che chiude la sequenza dei vincoli del piano: tre
incognite scalari $X(A), Y(A), \mu(A)$ in dualità con
le tre componenti cinematiche $u(A), v(A), \vartheta$,
esaurendo i gradi di libertà del corpo rigido.



%---------------------------------------------------------------------------------------------------------------------------- 
\section{Vincoli che collegano più corpi}\label{sec:vincoli_multipli}
%---------------------------------------------------------------------------------------------------------------------------- 

Nelle sezioni precedenti i vincoli interni sono stati
descritti nel caso più semplice: un vincolo che collega
due corpi $\mathcal{C}_i$ e $\mathcal{C}_j$. In molte
applicazioni strutturali --- nodi di telaio, archi a tre
cerniere, sistemi reticolari --- uno stesso vincolo collega
simultaneamente $k \geq 3$ corpi in un unico punto. In
questi casi la molteplicità del vincolo non è semplicemente
quella del vincolo binario, ma va calcolata contando le
equazioni cinematiche indipendenti che il vincolo impone
sull'insieme dei corpi collegati.

Il principio generale è il seguente: fissato un corpo di
riferimento $\mathcal{C}_1$, il vincolo impone che ciascuno
degli altri $k-1$ corpi soddisfi le stesse condizioni
cinematiche rispetto a $\mathcal{C}_1$. Ogni condizione
binaria contribuisce con la propria molteplicità, e la
molteplicità totale del vincolo multiplo è il prodotto
della molteplicità binaria per $(k-1)$.

Ci limiteremo qui ai due casi di interesse applicativo
diretto: la \emph{cerniera multipla} e l'\emph{incastro
	multiplo}, considerandoli sia come vincoli puramente
interni, sia come vincoli interni il cui nodo è a sua
volta collegato al suolo da un ulteriore vincolo esterno.

\subsection{La cerniera multipla}
\label{sec:cerniera_multipla}

Una \emph{cerniera multipla}\index{cerniera!multipla|textbf} collegante $k$ corpi in un
unico punto $A$ impone che tutti i corpi abbiano lo stesso
spostamento in $A$:
\begin{equation}\label{eq:cerniera_multipla}
	\ub_2(A) = \ub_1(A), \quad
	\ub_3(A) = \ub_1(A), \quad
	\ldots, \quad
	\ub_k(A) = \ub_1(A).
\end{equation}
Ciascuna delle $k-1$ condizioni di uguaglianza fornisce
2 equazioni scalari indipendenti. La molteplicità totale
del vincolo interno è quindi:
\begin{equation}\label{eq:molt_cerniera_multipla}
	m_\text{int} = 2(k-1).
\end{equation}

\begin{oss}
	Per $k=2$ si ritrova la cerniera binaria di molteplicità
	2. Per $k=3$ la molteplicità è 4: la cerniera tripla
	non è equivalente a due cerniere binarie applicate
	separatamente, ma introduce tante equazioni quante ne
	introdurrebbero due cerniere binarie applicate nello
	stesso punto.
\end{oss}

\paragraph{Prestazione statica.}
Le forze scambiate fra i corpi sono applicate in $A$. Per il
principio dei lavori virtuali, gli spostamenti virtuali
compatibili con la cerniera multipla sono le rotazioni
rigide indipendenti $\vartheta_1, \ldots, \vartheta_k$ di
ciascun corpo attorno ad $A$. Detta $\rb_i(A)$ la forza che
il vincolo trasmette al corpo $\mathcal{C}_i$, l'equilibrio
del nodo si scrive:
\begin{equation}
	\sum_{i=1}^{k} \rb_i(A) = \zerovb.
\end{equation}
Le $2(k-1)$ incognite statiche indipendenti sono le
componenti delle forze scambiate fra i corpi, in accordo
con la molteplicità. La \FIGref{fig:cerniera_tripla}
illustra le prestazioni cinematica e statica della cerniera
multipla nel caso $k = 3$.



%%%%%%%%%%%%%%%%%%%%%%%%%%%%%%%%%%%%%%%%%%%
\begin{figure}[htbp]
	\centering
	\begin{tikzpicture}[>=Stealth, line width=0.8pt, scale=0.95]
		
		% ═══════════════════════════════════════════════════════════
		% PANNELLO (a) — Prestazione cinematica della cerniera tripla
		% ═══════════════════════════════════════════════════════════
		\begin{scope}[xshift=0cm]
			
			%% ── Parametri ──
			\def\rA{0.28}   % raggio del pallino della cerniera in A
			\def\rcorpo{1.10}   % distanza del centro dei corpi da A
			\def\rrot{0.32}  % raggio degli archi di rotazione
			
			
			% Angoli di disposizione dei tre corpi (in gradi)
			\def\angCuno{90}       % C_1 in alto
			\def\angCdue{-30}      % C_2 in basso a destra
			\def\angCtre{-150}     % C_3 in basso a sinistra
			
			% Coordinate dei centri dei corpi
			\pgfmathsetmacro{\xCuno}{\rcorpo*cos(\angCuno)}
			\pgfmathsetmacro{\yCuno}{\rcorpo*sin(\angCuno)}
			\pgfmathsetmacro{\xCdue}{\rcorpo*cos(\angCdue)}
			\pgfmathsetmacro{\yCdue}{\rcorpo*sin(\angCdue)}
			\pgfmathsetmacro{\xCtre}{\rcorpo*cos(\angCtre)}
			\pgfmathsetmacro{\yCtre}{\rcorpo*sin(\angCtre)}
			
			%% ── Corpo C_1 (in alto) ──
			\begin{scope}[shift={(\xCuno, \yCuno)}, rotate=\angCuno-90]
				\draw[line width=1pt] 
				plot[smooth cycle, tension=0.7] coordinates 
				{(0.00, -0.85) (0.50, -0.65) (0.80, -0.10) 
					(0.70, 0.55) (0.20, 0.90) (-0.40, 0.80) 
					(-0.75, 0.30) (-0.70, -0.30) (-0.40, -0.75)};
				\node at (0.10, 0.40) {$\mathcal{C}_1$};
			\end{scope}
			
			
			
			
			%% ── Corpo C_2 ──
			\begin{scope}[shift={(\xCdue, \yCdue)}, rotate=\angCdue-90]
				\draw[line width=1pt] 
				plot[smooth cycle, tension=0.7] coordinates 
				{(0.00, -0.85) (0.50, -0.65) (0.80, -0.10) 
					(0.70, 0.55) (0.20, 0.90) (-0.40, 0.80) 
					(-0.75, 0.30) (-0.70, -0.30) (-0.40, -0.75)};
				\node at (0.20, 0.50) {$\mathcal{C}_2$};
			\end{scope}
			
			%% ── Corpo C_3 ──
			\begin{scope}[shift={(\xCtre, \yCtre)}, rotate=\angCtre-90]
				\draw[line width=1pt] 
				plot[smooth cycle, tension=0.7] coordinates 
				{(0.00, -0.85) (0.50, -0.65) (0.80, -0.10) 
					(0.70, 0.55) (0.20, 0.90) (-0.40, 0.80) 
					(-0.75, 0.30) (-0.70, -0.30) (-0.40, -0.75)};
				\node at (-0.30, 0.50) {$\mathcal{C}_3$};
			\end{scope}
			
			%% ── Pallino della cerniera in A ──
			\draw[line width=0.8pt, fill=white] (0, 0) circle (\rA);
			\node[font=\small] at (0, 0) {$A$};
			
			%% ── Equazioni di vincolo (a sinistra del pannello) ──
			\node[font=\footnotesize, anchor=west] 	at (-3.4, 1.30) {$\ub_2(A) = \ub_1(A)$};
			\node[font=\footnotesize, anchor=west] 	at (-3.40, 0.80) {$\ub_3(A) = \ub_1(A)$};
			
			%% ── Etichetta pannello ──
			\node[anchor=north, font=\small] 
			at (0, -1.40) {$(a)$};
			
		\end{scope}
		
		% ═══════════════════════════════════════════════════════════
		% PANNELLO (b) — Prestazione statica della cerniera tripla
		% ═══════════════════════════════════════════════════════════
		\begin{scope}[xshift=6cm]
			
			%% ── Parametri (identici ad (a)) ──
			\def\rA{0.28}
			\def\rcorpo{1.10}
			
			\def\angCuno{90}
			\def\angCdue{-30}
			\def\angCtre{-150}
			
			\pgfmathsetmacro{\xCuno}{\rcorpo*cos(\angCuno)}
			\pgfmathsetmacro{\yCuno}{\rcorpo*sin(\angCuno)}
			\pgfmathsetmacro{\xCdue}{\rcorpo*cos(\angCdue)}
			\pgfmathsetmacro{\yCdue}{\rcorpo*sin(\angCdue)}
			\pgfmathsetmacro{\xCtre}{\rcorpo*cos(\angCtre)}
			\pgfmathsetmacro{\yCtre}{\rcorpo*sin(\angCtre)}
			
			%% ── Corpo C_1 (identico ad (a)) ──
			\begin{scope}[shift={(\xCuno, \yCuno)}, rotate=\angCuno-90]
				\draw[line width=1pt] 
				plot[smooth cycle, tension=0.7] coordinates 
				{(0.00, -0.85) (0.50, -0.65) (0.80, -0.10) 
					(0.70, 0.55) (0.20, 0.90) (-0.40, 0.80) 
					(-0.75, 0.30) (-0.70, -0.30) (-0.40, -0.75)};
				\node at (0.10, 0.40) {$\mathcal{C}_1$};
			\end{scope}
			
			%% ── Corpo C_2 (identico ad (a)) ──
			\begin{scope}[shift={(\xCdue, \yCdue)}, rotate=\angCdue-90]
				\draw[line width=1pt] 
				plot[smooth cycle, tension=0.7] coordinates 
				{(0.00, -0.85) (0.50, -0.65) (0.80, -0.10) 
					(0.70, 0.55) (0.20, 0.90) (-0.40, 0.80) 
					(-0.75, 0.30) (-0.70, -0.30) (-0.40, -0.75)};
				\node at (0.30, 0.40) {$\mathcal{C}_2$};
			\end{scope}
			
			%% ── Corpo C_3 (identico ad (a)) ──
			\begin{scope}[shift={(\xCtre, \yCtre)}, rotate=\angCtre-90]
				\draw[line width=1pt] 
				plot[smooth cycle, tension=0.7] coordinates 
				{(0.00, -0.85) (0.50, -0.65) (0.80, -0.10) 
					(0.70, 0.55) (0.20, 0.90) (-0.40, 0.80) 
					(-0.75, 0.30) (-0.70, -0.30) (-0.40, -0.75)};
				\node at (-0.30, 0.50) {$\mathcal{C}_3$};
			\end{scope}
			
			%% ── Pallino della cerniera in A ──
			\draw[line width=0.8pt, fill=white] (0, 0) circle (\rA);
			\node[font=\small] at (0, 0) {$A$};
			
			%% ──────────────────────────────────────────────────
			%% REAZIONI: tre forze r_1, r_2, r_3 uscenti da A
			%% radialmente verso i rispettivi corpi
			%% ──────────────────────────────────────────────────
			
			\def\Lfreccia{0.50}
			
			% Reazione r_1 (verso C_1, in alto)
			\draw[->, line width=1.4pt] 
			({\rA*cos(\angCuno)}, {\rA*sin(\angCuno)}) 
			-- ++({\Lfreccia*cos(\angCuno)}, {\Lfreccia*sin(\angCuno)});
			\node[font=\small, anchor=west] 
			at ({(\rA+\Lfreccia)*cos(\angCuno)+0.10}, 
			{(\rA+\Lfreccia)*sin(\angCuno)}) {$\rb_1$};
			
			% Reazione r_2 (verso C_2, in basso a destra)
			\draw[->, line width=1.4pt] 
			({\rA*cos(\angCdue)}, {\rA*sin(\angCdue)}) 
			-- ++({\Lfreccia*cos(\angCdue)}, {\Lfreccia*sin(\angCdue)});
			\node[font=\small, anchor=west] 
			at ({(\rA+\Lfreccia)*cos(\angCdue)+0.10}, 
			{(\rA+\Lfreccia)*sin(\angCdue)-0.05}) {$\rb_2$};
			
			% Reazione r_3 (verso C_3, in basso a sinistra)
			\draw[->, line width=1.4pt] 
			({\rA*cos(\angCtre)}, {\rA*sin(\angCtre)}) 
			-- ++({\Lfreccia*cos(\angCtre)}, {\Lfreccia*sin(\angCtre)});
			\node[font=\small, anchor=east] 
			at ({(\rA+\Lfreccia)*cos(\angCtre)-0.10}, 
			{(\rA+\Lfreccia)*sin(\angCtre)-0.05}) {$\rb_3$};
			
			%% ── Equazione di equilibrio del nodo (a sinistra) ──
			\node[font=\footnotesize, anchor=west] 
			at (-3.40, 0.5) {$\rb_1 + \rb_2 + \rb_3 = \zerovb$};
			
			%% ── Etichetta pannello ──
			\node[anchor=north, font=\small] 
			at (0, -1.40) {$(b)$};
			
		\end{scope}
		
	\end{tikzpicture}
	\caption[Cerniera multipla  fra tre corpi]{%
		Cerniera multipla  fra i corpi $\mathcal{C}_1$,
		$\mathcal{C}_2$, $\mathcal{C}_3$ nel punto $A$.
		Prestazione cinematica (a):  le equazioni di vincolo
		$\ub_2(A) = \ub_1(A)$ e $\ub_3(A) = \ub_1(A)$ impongono
		l'uguaglianza degli spostamenti dei tre corpi in $A$;
		le rotazioni dei tre corpi restano libere e non sono
		rappresentate. Il \CRr\ di ogni coppia di corpi coincide
		con $A$. Prestazione statica (b): i tre corpi si scambiano in $A$ tre forze
		$\rb_1$, $\rb_2$, $\rb_3$ di direzione arbitraria,
		soggette all'equilibrio del nodo $\sum_i \rb_i = \zerovb$;
		nessuna coppia scambiata. Le incognite indipendenti
		sono le componenti di $\rb_2$ e $\rb_3$ (per un totale
		di $2(k-1) = 4$), in accordo con la molteplicità.}
	\label{fig:cerniera_tripla}
\end{figure}
%%%%%%%%%%%%%%%%%%%%%%%%%%%%%%%%%%%%%%%%%%%

\begin{oss}[Lettura simmetrica delle forze interne]
	Le equazioni cinematiche \eqref{eq:cerniera_multipla}
	sono scritte rispetto al corpo $\mathcal{C}_1$, scelto
	come riferimento, e impongono che ciascuno degli altri
	$k-1$ corpi abbia in $A$ lo stesso spostamento di
	$\mathcal{C}_1$. Coerentemente, la lettura delle reazioni
	è binaria: per ciascun $j = 2, \ldots, k$, la coppia
	azione-reazione $(\rb_j, -\rb_j)$ rappresenta la forza
	scambiata in $A$ fra il corpo $\mathcal{C}_j$ e il corpo
	$\mathcal{C}_1$, con $\rb_j$ agente su $\mathcal{C}_j$ e
	$-\rb_j$ agente su $\mathcal{C}_1$. Le incognite
	indipendenti sono le $k-1$ forze $\rb_2, \ldots, \rb_k$,
	ciascuna con due componenti scalari libere, per un totale
	di $2(k-1)$ in accordo con la molteplicità.
	
	La forza risultante che il corpo $\mathcal{C}_1$ riceve
	dal vincolo è la somma delle reazioni provenienti dai
	rimanenti corpi:
	\begin{equation}\label{eq:cerniera_multipla_r1}
		\rb_1 = -\sum_{j=2}^{k} \rb_j,
	\end{equation}
	che coincide con l'equazione di equilibrio del nodo
	$\sum_{i=1}^{k} \rb_i = \zerovb$ del paragrafo
	precedente. La scelta di $\mathcal{C}_1$ come corpo di
	riferimento è arbitraria: qualsiasi altro corpo può
	svolgere lo stesso ruolo, con un opportuno
	riordinamento degli indici e una corrispondente
	ridefinizione delle coppie azione-reazione.
\end{oss}

\paragraph{Forza esterna applicata al nodo.}
Quando una forza esterna $\fb$ --- attiva o reattiva --- è
applicata direttamente al nodo, l'equilibrio del nodo si
modifica in
\begin{equation}\label{eq:eq_nodo_F}
	\sum_{i=1}^{k} \rb_i(A) = \fb,
\end{equation}
e la forza risultante sul corpo $\mathcal{C}_1$ diventa
\begin{equation}
	\rb_1 = \fb - \sum_{j=2}^{k} \rb_j(A).
\end{equation}
Due i casi di interesse. Se $\fb$ è una forza
\emph{attiva} nota, essa è un dato del problema e non
incrementa il numero delle incognite: queste rimangono le
$2(k-1)$ componenti delle reazioni interne $\rb_2, \ldots,
\rb_k$, esattamente come nel caso $\fb = \zerovb$. Se
invece $\fb$ è la \emph{reazione} di un ulteriore vincolo
esterno applicato al nodo, le sue $m_\text{ext}$
componenti scalari libere si aggiungono alle incognite
del problema; questo caso è trattato nel paragrafo
successivo.



\paragraph{Cerniera multipla vincolata al suolo.}
Quando il nodo $A$ è collegato al suolo da un ulteriore
vincolo esterno, le sue equazioni cinematiche si sommano
a quelle interne: queste ultime coinvolgono spostamenti
relativi (i salti $\ub_j - \ub_1$), mentre quelle esterne
coinvolgono lo spostamento assoluto di $\mathcal{C}_1$ in
$A$. I due gruppi di equazioni sono quindi indipendenti
e la molteplicità complessiva è
\begin{equation}\label{eq:molt_cerniera_multipla_tot}
	m = 2(k-1) + m_\text{ext},
\end{equation}
con $m_\text{ext} \in \{1, 2\}$ a seconda del vincolo
esterno (rispettivamente carrello o cerniera).

Sul lato statico, la situazione corrisponde al caso del
paragrafo precedente con $\fb = \rb_\text{ext}$, dove
$\rb_\text{ext}$ è la reazione che il vincolo esterno
trasmette al nodo. L'equilibrio del nodo è allora
\begin{equation}
	\sum_{i=1}^{k} \rb_i(A) = \rb_\text{ext}(A),
\end{equation}
con $\rb_\text{ext}$ a $m_\text{ext}$ componenti scalari
libere ($R\,\ab$ se carrello, $X\,\ab_x + Y\,\ab_y$ se
cerniera). A differenza del caso con forza attiva, qui
$\rb_\text{ext}$ è incognita, e le sue componenti si
aggiungono alle $2(k-1)$ delle reazioni interne, per un
totale di $2(k-1) + m_\text{ext}$ incognite indipendenti
in accordo con la \eqref{eq:molt_cerniera_multipla_tot}.

\subsection{L'incastro multiplo}
\label{sec:incastro_multiplo}

Un \emph{incastro multiplo}\index{incastro!multiplo|textbf} collegante $k$ corpi impone
che tutti abbiano lo stesso spostamento e la stessa
rotazione:
\begin{align}
	\ub_2(A) &= \ub_1(A), \quad
	\ldots, \quad
	\ub_k(A) = \ub_1(A),\\
	\vartheta_2 &= \vartheta_1, \quad
	\ldots, \quad
	\vartheta_k = \vartheta_1.
\end{align}
Ciascuna delle $k-1$ condizioni fornisce 3 equazioni
scalari indipendenti. La molteplicità totale del vincolo
interno è:
\begin{equation}\label{eq:molt_incastro_multiplo}
	m_\text{int} = 3(k-1).
\end{equation}

\paragraph{Prestazione statica.}
Le forze e le coppie scambiate fra i corpi sono applicate
in $A$. Per il principio dei lavori virtuali, gli spostamenti
virtuali compatibili con l'incastro multiplo sono le
traslazioni e rotazioni rigide \emph{comuni} ai $k$ corpi
(che si comportano in $A$ come un unico corpo): qualunque
sistema di forze e coppie applicato in $A$ con risultante
e momento risultante nulli compie lavoro nullo su tali
spostamenti. Detta $\rb_i(A)$ la forza e $\mu_i(A)$ la
coppia che il vincolo trasmette al corpo $\mathcal{C}_i$,
l'equilibrio del nodo si scrive:
\begin{equation}
	\sum_{i=1}^{k} \rb_i(A) = \zerovb,
	\qquad
	\sum_{i=1}^{k} \mu_i(A) = 0.
\end{equation}
Le $3(k-1)$ incognite statiche indipendenti sono le
componenti delle forze e delle coppie scambiate fra i
corpi, in accordo con la molteplicità. La
\FIGref{fig:incastro_triplo} illustra le prestazioni
cinematica e statica dell'incastro multiplo nel caso
$k = 3$.

\begin{oss}[Lettura simmetrica delle reazioni interne]
	Come per la cerniera multipla, le equazioni cinematiche
	dell'incastro multiplo sono scritte rispetto al corpo
	$\mathcal{C}_1$, scelto come riferimento, e impongono
	che ciascuno degli altri $k-1$ corpi abbia in $A$ lo
	stesso spostamento e la stessa rotazione di
	$\mathcal{C}_1$. Coerentemente, la lettura delle
	reazioni è binaria: per ciascun $j = 2, \ldots, k$, la
	coppia azione-reazione $(\rb_j, -\rb_j)$ rappresenta la
	forza scambiata fra $\mathcal{C}_j$ e $\mathcal{C}_1$,
	e $(\mu_j, -\mu_j)$ la coppia scambiata. Le incognite
	indipendenti sono le $k-1$ forze $\rb_2, \ldots, \rb_k$
	e le $k-1$ coppie $\mu_2, \ldots, \mu_k$, per un totale
	di $3(k-1)$ in accordo con la molteplicità.
	
	La forza e la coppia risultanti che il corpo
	$\mathcal{C}_1$ riceve dal vincolo sono le somme delle
	reazioni provenienti dagli altri corpi:
	\begin{equation}\label{eq:incastro_multiplo_r1_mu1}
		\rb_1 = -\sum_{j=2}^{k} \rb_j,
		\qquad
		\mu_1 = -\sum_{j=2}^{k} \mu_j,
	\end{equation}
	che coincidono con le equazioni di equilibrio del nodo
	del paragrafo precedente.
\end{oss}


%%%%%%%%%%%%%%%%%%%%%%%%%%%%%%%%%%%%%%%%%%%
\begin{figure}[htbp]
	\centering
	\begin{tikzpicture}[>=Stealth, line width=0.8pt, scale=0.95]
		
		% ═══════════════════════════════════════════════════════════
		% PANNELLO (a) — Prestazione cinematica dell'incastro tripla
		% ═══════════════════════════════════════════════════════════
		\begin{scope}[xshift=0cm]
			
			%% ── Parametri ──
			\def\Lquadr{0.32}    % semi-lato del quadrato dell'incastro
			\def\rcorpo{1.10}
			\def\rrot{0.32}
			
			\def\angCuno{90}
			\def\angCdue{-30}
			\def\angCtre{-150}
			
			\pgfmathsetmacro{\xCuno}{\rcorpo*cos(\angCuno)}
			\pgfmathsetmacro{\yCuno}{\rcorpo*sin(\angCuno)}
			\pgfmathsetmacro{\xCdue}{\rcorpo*cos(\angCdue)}
			\pgfmathsetmacro{\yCdue}{\rcorpo*sin(\angCdue)}
			\pgfmathsetmacro{\xCtre}{\rcorpo*cos(\angCtre)}
			\pgfmathsetmacro{\yCtre}{\rcorpo*sin(\angCtre)}
			
			%% ── Corpi C_1, C_2, C_3 (identici alla cerniera multipla) ──
			\begin{scope}[shift={(\xCuno, \yCuno)}, rotate=\angCuno-90]
				\draw[line width=1pt] 
				plot[smooth cycle, tension=0.7] coordinates 
				{(0.00, -0.85) (0.50, -0.65) (0.80, -0.10) 
					(0.70, 0.55) (0.20, 0.90) (-0.40, 0.80) 
					(-0.75, 0.30) (-0.70, -0.30) (-0.40, -0.75)};
				\node at (0.10, -0.30) {$\mathcal{C}_1$};
			\end{scope}
			
			\begin{scope}[shift={(\xCdue, \yCdue)}, rotate=\angCdue-90]
				\draw[line width=1pt] 
				plot[smooth cycle, tension=0.7] coordinates 
				{(0.00, -0.85) (0.50, -0.65) (0.80, -0.10) 
					(0.70, 0.55) (0.20, 0.90) (-0.40, 0.80) 
					(-0.75, 0.30) (-0.70, -0.30) (-0.40, -0.75)};
				\node at (0.20, 0.30) {$\mathcal{C}_2$};
			\end{scope}
			
			\begin{scope}[shift={(\xCtre, \yCtre)}, rotate=\angCtre-90]
				\draw[line width=1pt] 
				plot[smooth cycle, tension=0.7] coordinates 
				{(0.00, -0.85) (0.50, -0.65) (0.80, -0.10) 
					(0.70, 0.55) (0.20, 0.90) (-0.40, 0.80) 
					(-0.75, 0.30) (-0.70, -0.30) (-0.40, -0.75)};
				\node at (-0.30, 0.50) {$\mathcal{C}_3$};
			\end{scope}
			
			%% ── Quadrato grigio dell'incastro in A ──
			\draw[line width=0.8pt, fill=gray!30] 
			(-\Lquadr, -\Lquadr) rectangle (\Lquadr, \Lquadr);
			\node[font=\small] at (0, 0) {$A$};
			
			%% ── Archi di rotazione ϑ_1=ϑ_2=ϑ_3 (semicerchi, 180°) ──
			\draw[->, line width=1pt] 
			({\xCuno+\rrot}, {\yCuno+0.40}) arc(0:180:\rrot);
			\node[font=\small] 
			at ({\xCuno-\rrot+0.3}, {\yCuno+0.40}) {$\vartheta_1$};
			
			\draw[->, line width=1pt] 
			({\xCdue+\rrot}, {\yCdue+0.20}) arc(0:180:\rrot);
			\node[font=\small] 
			at ({\xCdue+\rrot-0.30}, {\yCdue+0.20}) {$\vartheta_2$};
			
			\draw[->, line width=1pt] 
			({\xCtre+\rrot}, {\yCtre+0.20}) arc(0:180:\rrot);
			\node[font=\small] 
			at ({\xCtre-\rrot+0.30}, {\yCtre+0.20}) {$\vartheta_3$};
			
			%% ── Equazioni di vincolo (a sinistra del pannello) ──
			\node[font=\footnotesize, anchor=west] 
			at (-3.40, 1.50) {$\ub_2(A) = \ub_1(A)$};
			\node[font=\footnotesize, anchor=west] 
			at (-3.40, 1.00) {$\ub_3(A) = \ub_1(A)$};
			\node[font=\footnotesize, anchor=west] 
			at (-3.40, 0.50) {$\vartheta_2 = \vartheta_1$};
			\node[font=\footnotesize, anchor=west] 
			at (-3.40, 0.00) {$\vartheta_3 = \vartheta_1$};
			
			\node[anchor=north, font=\small] 
			at (0, -1.40) {$(a)$};
			
		\end{scope}
		
		% ═══════════════════════════════════════════════════════════
		% PANNELLO (b) — Prestazione statica dell'incastro tripla
		% ═══════════════════════════════════════════════════════════
		\begin{scope}[xshift=6cm]
			
			\def\Lquadr{0.32}
			\def\rcorpo{1.10}
			\def\rmu{0.30}
			
			\def\angCuno{90}
			\def\angCdue{-30}
			\def\angCtre{-150}
			
			\pgfmathsetmacro{\xCuno}{\rcorpo*cos(\angCuno)}
			\pgfmathsetmacro{\yCuno}{\rcorpo*sin(\angCuno)}
			\pgfmathsetmacro{\xCdue}{\rcorpo*cos(\angCdue)}
			\pgfmathsetmacro{\yCdue}{\rcorpo*sin(\angCdue)}
			\pgfmathsetmacro{\xCtre}{\rcorpo*cos(\angCtre)}
			\pgfmathsetmacro{\yCtre}{\rcorpo*sin(\angCtre)}
			
			
			
			%% ── Corpi (identici al pannello (a)) ──
			\begin{scope}[shift={(\xCuno, \yCuno)}, rotate=\angCuno-90]
				\draw[line width=1pt] 
				plot[smooth cycle, tension=0.7] coordinates 
				{(0.00, -0.85) (0.50, -0.65) (0.80, -0.10) 
					(0.70, 0.55) (0.20, 0.90) (-0.40, 0.80) 
					(-0.75, 0.30) (-0.70, -0.30) (-0.40, -0.75)};
				\node at (-0.30, 0.0) {$\mathcal{C}_1$};
			\end{scope}
			
			\begin{scope}[shift={(\xCdue, \yCdue)}, rotate=\angCdue-90]
				\draw[line width=1pt] 
				plot[smooth cycle, tension=0.7] coordinates 
				{(0.00, -0.85) (0.50, -0.65) (0.80, -0.10) 
					(0.70, 0.55) (0.20, 0.90) (-0.40, 0.80) 
					(-0.75, 0.30) (-0.70, -0.30) (-0.40, -0.75)};
				\node at (0.30, 0.40) {$\mathcal{C}_2$};
			\end{scope}
			
			\begin{scope}[shift={(\xCtre, \yCtre)}, rotate=\angCtre-90]
				\draw[line width=1pt] 
				plot[smooth cycle, tension=0.7] coordinates 
				{(0.00, -0.85) (0.50, -0.65) (0.80, -0.10) 
					(0.70, 0.55) (0.20, 0.90) (-0.40, 0.80) 
					(-0.75, 0.30) (-0.70, -0.30) (-0.40, -0.75)};
				\node at (-0.30, 0.50) {$\mathcal{C}_3$};
			\end{scope}
			
			%% ── Quadrato grigio dell'incastro in A ──
			\draw[line width=0.8pt, fill=gray!30] 
			(-\Lquadr, -\Lquadr) rectangle (\Lquadr, \Lquadr);
			\node[font=\small] at (0, 0) {$A$};
			
			%% ── Forze r_1, r_2, r_3 (radiali, dal bordo del quadrato) ──
			\def\Lfreccia{0.50}
			
			% Reazione r_1
			\draw[->, line width=1.4pt] 
			({\Lquadr*cos(\angCuno)}, {\Lquadr*sin(\angCuno)}) 
			-- ++({\Lfreccia*cos(\angCuno)}, {\Lfreccia*sin(\angCuno)});
			\node[font=\small, anchor=west] 
			at ({(\Lquadr+\Lfreccia)*cos(\angCuno)+0.10}, 
			{(\Lquadr+\Lfreccia)*sin(\angCuno)}) {$\rb_1$};
			
			% Reazione r_2
			\draw[->, line width=1.4pt] 
			({\Lquadr*cos(\angCdue)}, {\Lquadr*sin(\angCdue)}) 
			-- ++({\Lfreccia*cos(\angCdue)}, {\Lfreccia*sin(\angCdue)});
			\node[font=\small, anchor=north] 
			at ({(\Lquadr+\Lfreccia)*cos(\angCdue)+0.0}, 
			{(\Lquadr+\Lfreccia)*sin(\angCdue)-0.05}) {$\rb_2$};
			
			% Reazione r_3
			\draw[->, line width=1.4pt] 
			({\Lquadr*cos(\angCtre)}, {\Lquadr*sin(\angCtre)}) 
			-- ++({\Lfreccia*cos(\angCtre)}, {\Lfreccia*sin(\angCtre)});
			\node[font=\small, anchor=west] 
			at ({(\Lquadr+\Lfreccia)*cos(\angCtre)-0.10}, 
			{(\Lquadr+\Lfreccia)*sin(\angCtre)-0.15}) {$\rb_3$};
			
			%% ── Coppie μ_1, μ_2, μ_3 (semicerchi dentro ai corpi) ──
			% Coppia in C_1 (in alto, dentro al corpo, sotto la label)
			
			% Coppia μ_1
			\draw[->, line width=1.4pt] 
			({\xCuno+\rmu}, {\yCuno+0.40}) arc(0:180:\rmu);
			\node[font=\small] 
			at ({\xCuno-\rmu+0.30}, {\yCuno+0.40}) {$\mu_1$};
			
			% Coppia μ_2
			\draw[->, line width=1.4pt] 
			({\xCdue+\rmu}, {\yCdue+0.20}) arc(0:180:\rmu);
			\node[font=\small] 
			at ({\xCdue+\rmu-0.30}, {\yCdue+0.20}) {$\mu_2$};
			
			% Coppia μ_3
			\draw[->, line width=1.4pt] 
			({\xCtre+\rmu}, {\yCtre+0.20}) arc(0:180:\rmu);
			\node[font=\small] 
			at ({\xCtre-\rmu+0.30}, {\yCtre+0.20}) {$\mu_3$};
			
			
			
			%% ── Equazioni di equilibrio del nodo (a sinistra) ──
			\node[font=\footnotesize, anchor=west] 
			at (-3.40, 0.80) {$\sum_i \rb_i = \zerovb$};
			\node[font=\footnotesize, anchor=west] 
			at (-3.40, 0.30) {$\sum_i \mu_i = 0$};
			
			\node[anchor=north, font=\small] 
			at (0, -1.40) {$(b)$};
			
		\end{scope}
		
	\end{tikzpicture}
	\caption[Incastro multiplo fra tre corpi]{%
		Incastro multiplo fra i corpi $\mathcal{C}_1$, $\mathcal{C}_2$,
		$\mathcal{C}_3$ nel punto $A$. Prestazione cinematica (a):
		le equazioni di vincolo $\ub_2(A) = \ub_1(A)$,
		$\ub_3(A) = \ub_1(A)$, $\vartheta_2 = \vartheta_1$,
		$\vartheta_3 = \vartheta_1$ impongono l'uguaglianza degli
		spostamenti e delle rotazioni dei tre corpi in $A$; i tre
		corpi si comportano come un unico corpo rigido nel nodo.
		Prestazione statica (b): i tre corpi si scambiano in $A$
		forze $\rb_1, \rb_2, \rb_3$ e coppie $\mu_1, \mu_2, \mu_3$
		di direzione e verso arbitrari, soggette all'equilibrio
		del nodo $\sum_i \rb_i = \zerovb$ e $\sum_i \mu_i = 0$.
		Le incognite indipendenti sono le componenti delle
		coppie e delle forze relative a $\mathcal{C}_2$ e
		$\mathcal{C}_3$ (per un totale di $3(k-1) = 6$), in
		accordo con la molteplicità.}
	\label{fig:incastro_triplo}
\end{figure}
%%%%%%%%%%%%%%%%%%%%%%%%%%%%%%%%%%%%%%%%%%%

\paragraph{Forza e coppia esterne applicate al nodo.}
Quando una forza esterna $\fb$ e/o una coppia esterna $\mu$
--- attive o reattive --- sono applicate direttamente al
nodo, l'equilibrio del nodo si modifica in
\begin{equation}\label{eq:eq_nodo_FC}
	\sum_{i=1}^{k} \rb_i(A) = \fb,
	\qquad
	\sum_{i=1}^{k} \mu_i(A) = \mu,
\end{equation}
e la forza e la coppia risultanti sul corpo
$\mathcal{C}_1$ diventano
\begin{equation}
	\rb_1 = \fb - \sum_{j=2}^{k} \rb_j(A),
	\qquad
	\mu_1 = \mu - \sum_{j=2}^{k} \mu_j(A).
\end{equation}
Se $\fb$ e $\mu$ sono \emph{attive} note, esse sono dati del
problema e non incrementano il numero delle incognite,
che rimangono $3(k-1)$. Se invece sono le \emph{reazioni}
di un ulteriore vincolo esterno applicato al nodo, le
loro componenti scalari libere si aggiungono alle
incognite del problema; questo caso è trattato nel
paragrafo successivo.

\paragraph{Incastro multiplo vincolato al suolo.}
Quando il nodo $A$ è collegato al suolo da un ulteriore
vincolo esterno, le sue equazioni cinematiche si sommano
a quelle interne, in modo del tutto analogo al caso della
cerniera multipla. La molteplicità complessiva è
\begin{equation}\label{eq:molt_incastro_multiplo_tot}
	m = 3(k-1) + m_\text{ext},
\end{equation}
con $m_\text{ext} \in \{1, 2, 3\}$ a seconda che il nodo
sia vincolato al suolo da un carrello, una cerniera o un
incastro esterno.

Sul lato statico, la situazione corrisponde al caso del
paragrafo precedente con $\fb = \rb_\text{ext}$ e
$C = \mu_\text{ext}$, dove $\rb_\text{ext}$ e
$\mu_\text{ext}$ sono la forza e la coppia trasmesse dal
vincolo esterno al nodo. L'equilibrio del nodo è allora
\begin{equation}
	\sum_{i=1}^{k} \rb_i(A) = \rb_\text{ext}(A),
	\qquad
	\sum_{i=1}^{k} \mu_i(A) = \mu_\text{ext}(A),
\end{equation}
con complessivamente $m_\text{ext}$ componenti scalari
libere ($R\,\ab$ per il carrello, $X\,\ab_x + Y\,\ab_y$
per la cerniera, $X\,\ab_x + Y\,\ab_y$ e $\mu$ per
l'incastro esterno). A differenza del caso con forza
attiva, qui $\rb_\text{ext}$ e $\mu_\text{ext}$ sono
incognite, e le loro componenti si aggiungono alle
$3(k-1)$ delle reazioni interne, per un totale di
$3(k-1) + m_\text{ext}$ incognite indipendenti in accordo
con la \eqref{eq:molt_incastro_multiplo_tot}.


\medskip

\noindent
La tavola seguente riassume le molteplicità dei vincoli
multipli in funzione del numero $k$ di corpi collegati.

\begin{table}[H]
	\centering
	\resizebox{\textwidth}{!}{%
		\renewcommand{\arraystretch}{1.8}
		\small
		\begin{tabular}{>{\bfseries}lccc}
			\toprule
			Vincolo & Vincolo esterno
			& Molteplicità
			& Reazione scambiata\\
			\midrule
			Cerniera multipla & assente
			& $2(k-1)$
			& $k-1$ forze indipendenti\\
			Cerniera multipla & carrello
			& $2(k-1) + 1$
			& $k-1$ forze interne $+$ $R\,\ab$ esterna\\
			Cerniera multipla & cerniera
			& $2(k-1) + 2$
			& $k-1$ forze interne $+$ forza esterna\\
			\midrule
			Incastro multiplo & assente
			& $3(k-1)$
			& $k-1$ forze e coppie indipendenti\\
			Incastro multiplo & carrello
			& $3(k-1) + 1$
			& precedenti $+$ $R\,\ab$ esterna\\
			Incastro multiplo & cerniera
			& $3(k-1) + 2$
			& precedenti $+$ forza esterna\\
			Incastro multiplo & incastro
			& $3(k-1) + 3$
			& precedenti $+$ forza e coppia esterne\\
			\bottomrule
		\end{tabular}%
	}
	\caption[Molteplicità dei vincoli multipli]{%
		Molteplicità dei vincoli che collegano $k$ corpi
		simultaneamente in un punto $A$, con eventuale vincolo
		esterno che collega il nodo al suolo.}
	\label{tab:vincoli_multipli}
\end{table}

\begin{oss}
	L'equazione $m = m_\text{int} + m_\text{ext}$ ha una
	lettura statica meno immediata di quella cinematica.
	Sul lato cinematico la somma è trasparente: le
	equazioni interne riguardano i salti
	$\ub_j - \ub_1$ e $\vartheta_j - \vartheta_1$, quelle
	esterne lo spostamento assoluto e la rotazione del
	nodo, e i due gruppi sono manifestamente indipendenti.
	Sul lato statico, invece, l'aggiunta del vincolo
	esterno non aggiunge semplicemente $m_\text{ext}$
	incognite: ne aggiunge $m_\text{ext}$ \emph{e
		contemporaneamente} libera altrettante equazioni di
	equilibrio del nodo. In assenza di vincolo esterno,
	l'equilibrio $\sum_i \rb_i = \zerovb$ riduce le $2k$
	componenti delle forze interne a $2(k-1)$ indipendenti;
	in presenza, $\sum_i \rb_i + \rb_\text{ext} = \zerovb$
	è soddisfatto identicamente nelle $m_\text{ext}$
	componenti libere di $\rb_\text{ext}$ e impone solo le
	rimanenti $2 - m_\text{ext}$ relazioni sulle incognite
	interne. Il conteggio finale, $2(k-1) + m_\text{ext}$,
	è lo stesso che si otterrebbe per somma diretta, ma il
	bilancio è diverso: il vincolo esterno non si limita ad
	aggiungere incognite, libera quelle interne dall'obbligo
	di equilibrarsi da sole.
\end{oss}

\begin{oss}
	I risultati di questa sezione costituiscono la base
	per la formulazione del problema cinematico e statico
	di sistemi di corpi rigidi, che sarà affrontata nel
	\CAPref{cap:problemi_cin_stat_SCR}. Il dato fondamentale
	che si trasferisce è la \emph{molteplicità} di ciascun
	vincolo: essa conta simultaneamente il numero di
	equazioni cinematiche indipendenti introdotte dal
	vincolo e il numero di incognite statiche che esso
	introduce. La classificazione del sistema --- labile,
	isostatico o iperstatico --- dipende dal confronto fra
	il numero totale di gradi di libertà dei corpi e la
	somma delle molteplicità di tutti i vincoli, interni
	ed esterni. In questo conteggio i vincoli multipli vanno
	trattati con la molteplicità corretta ricavata in
	questa sezione: sostituire una cerniera tripla con due
	cerniere binarie applicate separatamente, o trascurare
	l'aggiunta di un vincolo esterno al nodo, conduce a
	conteggi errati e a classificazioni sbagliate del
	sistema.
\end{oss}

\begin{oss}
	La formulazione adottata in questa sezione --- al netto
	dell'eventuale vincolo esterno --- esprime le equazioni
	di vincolo con riferimento al corpo $\mathcal{C}_1$:
	gli spostamenti degli altri corpi
	$\mathcal{C}_2, \ldots, \mathcal{C}_k$ sono vincolati
	a coincidere con quelli di $\mathcal{C}_1$ nel punto
	(o nella direzione) comune. Questa scelta è conveniente
	dal punto di vista formale ma nasconde una asimmetria:
	i cedimenti relativi fra due corpi $\mathcal{C}_i$ e
	$\mathcal{C}_j$ con $i,j \neq 1$ non possono essere
	imposti direttamente, ma solo indirettamente attraverso
	i cedimenti rispetto a $\mathcal{C}_1$. La formulazione
	più generale richiede di scrivere le equazioni di
	vincolo per ogni coppia $(\mathcal{C}_i, \mathcal{C}_j)$,
	verificando che le equazioni risultanti siano
	indipendenti.
	
	La stessa asimmetria si riflette nella descrizione delle
	reazioni. Nel caso della cerniera multipla pura (senza
	vincolo esterno), dette $\rb_j$ le forze reattive che il
	vincolo esercita sui corpi $\mathcal{C}_j$ con $j = 2,
	\ldots, k$, la forza reattiva sul corpo $\mathcal{C}_1$
	è determinata dall'equilibrio del nodo:
	\begin{equation}
		\rb_1 = -\sum_{j=2}^{k} \rb_j.
	\end{equation}
	Le incognite statiche indipendenti sono quindi le $k-1$
	forze $\rb_2, \ldots, \rb_k$, ciascuna con due
	componenti scalari libere, per un totale di $2(k-1)$
	incognite. La forza fra due corpi $\mathcal{C}_i$ e
	$\mathcal{C}_j$ con $i,j \neq 1$ non è direttamente una
	delle incognite, ma si ricava per differenza
	$\rb_i - \rb_j$ una volta note le reazioni rispetto a
	$\mathcal{C}_1$. In presenza di vincolo esterno
	l'equilibrio del nodo coinvolge anche $\rb_\text{ext}$,
	come illustrato nel paragrafo precedente.
\end{oss}


%---------------------------------------------------------------------------------------------------------------------------- 
\section{Cenno ai vincoli spaziali}\label{sec:spaziali}
%---------------------------------------------------------------------------------------------------------------------------- 

Nel caso spaziale il corpo rigido libero possiede sei gradi
di libertà: tre componenti di traslazione e tre componenti
di rotazione. I vincoli posizionali\index{vincolo!spaziale} hanno quindi molteplicità
$m \in \{1,\ldots,6\}$ e le reazioni corrispondenti sono
sistemi di forze e coppie in tre dimensioni.

I vincoli piani studiati nelle sezioni precedenti si
ritrovano come casi particolari di quelli spaziali,
ottenuti confinando il moto in un piano. La struttura
concettuale rimane invariata: ad ogni equazione cinematica
scalare corrisponde una e una sola componente reattiva,
determinata dal principio dei lavori virtuali.

Vediamo i vincoli spaziali fondamentali.

\paragraph{Piano di scorrimento (molteplicità 1).}\index{piano di scorrimento}\index{piano di scorrimento|seealso{vincolo}}
Impedisce la traslazione del punto di applicazione $A$
nella direzione del versore $\ab$, lasciando libere le
due traslazioni nel piano perpendicolare e le tre
rotazioni. È l'analogo spaziale del carrello piano.
La reazione è una forza $R\,\ab$.

\paragraph{Biella (molteplicità 1).}
Impedisce la traslazione del punto di applicazione $A$
lungo l'asse dell'asta. La reazione è una forza $R\,\ab$
diretta lungo l'asse della biella. Coincide con il piano
di scorrimento nel regime infinitesimo, in analogia con
l'equivalenza biella-carrello già osservata nel caso
piano (\FIGref{fig:carrello}).

\paragraph{Bipendolo (molteplicità 1).}
Impedisce la rotazione del corpo attorno a un asse
assegnato $\ab$, lasciando libere le tre traslazioni e
le due rotazioni attorno agli assi perpendicolari. La
reazione è una coppia $\mu\,\ab$ diretta lungo $\ab$.

\paragraph{Cerniera sferica (molteplicità 3).}\index{cerniera!sferica}
Fissa la posizione del punto di applicazione $A$
lasciando libere tutte e tre le rotazioni. È
l'analogo spaziale della cerniera piana. La reazione
è una forza arbitraria:
\begin{equation}
	\rb(A) = X(A)\,\ab_x + Y(A)\,\ab_y + Z(A)\,\ab_z.
\end{equation}

\paragraph{Perno scorrevole (molteplicità 4).}\index{perno scorrevole}\index{perno scorrevole|seealso{vincolo}}
Impedisce le due traslazioni perpendicolari all'asse
$\ab$ del perno e le due rotazioni attorno agli assi perpendicolari
ad $\ab$, lasciando libere la traslazione lungo $\ab$ e
la rotazione attorno ad $\ab$. Realizza il vincolo
classico ``albero in un foro cilindrico'' in cui l'albero
può scorrere lungo l'asse del foro e ruotare attorno ad
esso. Posto $\ab = \ab_z$, la reazione è una forza nel
piano perpendicolare ed una coppia nello stesso piano:
\begin{equation}
	\rb(A) = X(A)\,\ab_x + Y(A)\,\ab_y,
	\qquad
	\bm{\mu} = \mu_x\,\ab_x + \mu_y\,\ab_y.
\end{equation}

\paragraph{Glifo (molteplicità 5).}
Impedisce le due traslazioni perpendicolari all'asse $\ab$
e tutte e tre le rotazioni, lasciando libera la sola
traslazione lungo $\ab$. Realizza una \emph{guida
	prismatica}\index{guida prismatica|see{glifo}}: il corpo scorre lungo l'asse del vincolo
mantenendo il proprio orientamento. Si ottiene aggiungendo
al perno scorrevole un terzo bipendolo con asse $\ab$, che
blocca anche la rotazione attorno ad $\ab$. Posto $\ab =
\ab_z$, la reazione è una forza nel piano perpendicolare
ed una coppia arbitraria:
\begin{equation}
	\rb(A) = X(A)\,\ab_x + Y(A)\,\ab_y,
	\qquad
	\bm{\mu} = \mu_x\,\ab_x + \mu_y\,\ab_y + \mu_z\,\ab_z.
\end{equation}


\paragraph{Cerniera cilindrica (molteplicità 5).}\index{cerniera!cilindrica}
Impedisce tutte e tre le traslazioni e le due rotazioni
attorno agli assi perpendicolari all'asse $\ab$,
lasciando libera la sola rotazione attorno ad $\ab$. È
il vincolo del perno scorrevole con l'aggiunta del
blocco della traslazione lungo $\ab$, ossia un cardine
in cui l'albero, oltre a non potersi sfilare, ruota
soltanto attorno al proprio asse. Posto $\ab = \ab_z$,
la reazione è una forza arbitraria ed una coppia nel
piano perpendicolare:
\begin{equation}
	\rb(A) = X(A)\,\ab_x + Y(A)\,\ab_y + Z(A)\,\ab_z,
	\qquad
	\bm{\mu} = \mu_x\,\ab_x + \mu_y\,\ab_y.
\end{equation}

\paragraph{Incastro (molteplicità 6).}
Blocca completamente il corpo: impedisce tutte e tre
le traslazioni e tutte e tre le rotazioni. Non esiste
alcun grado di libertà residuo. La reazione è un
sistema generico di forza e coppia:
\begin{equation}
	\rb(A) = X(A)\,\ab_x + Y(A)\,\ab_y + Z(A)\,\ab_z,
	\qquad
	\bm{\mu} = \mu_x\,\ab_x + \mu_y\,\ab_y +
	\mu_z\,\ab_z.
\end{equation}

\medskip

\noindent
La tavola seguente riassume le prestazioni dei vincoli
spaziali, come estensione della
tavola~\ref{tab:vincoli_esterni}.

\begin{table}[H]
	\centering
	\resizebox{\textwidth}{!}{%
		\renewcommand{\arraystretch}{1.8}
		\small
		\begin{tabular}{>{\bfseries}lccll}
			\toprule
			Vincolo & Molt.
			& GdL res.
			& Forza reattiva
			& Coppia reattiva\\
			\midrule
			Piano di scorrimento & 1
			& 5
			& $R\,\ab$
			& nessuna\\
			Biella & 1
			& 5
			& $R\,\ab$
			& nessuna\\
			Bipendolo & 1
			& 5
			& nessuna
			& $\mu\,\ab$\\
			Cerniera sferica & 3
			& 3
			& $X\,\ab_x+Y\,\ab_y+Z\,\ab_z$
			& nessuna\\
			Perno scorrevole & 4
			& 2
			& $X\,\ab_x+Y\,\ab_y$
			& $\mu_x\,\ab_x+\mu_y\,\ab_y$\\
			Glifo & 5
			& 1
			& $X\,\ab_x+Y\,\ab_y$
			& $\mu_x\,\ab_x+\mu_y\,\ab_y+\mu_z\,\ab_z$\\			
			Cerniera cilindrica & 5
			& 1
			& $X\,\ab_x+Y\,\ab_y+Z\,\ab_z$
			& $\mu_x\,\ab_x+\mu_y\,\ab_y$\\
			Incastro & 6
			& 0
			& $X\,\ab_x+Y\,\ab_y+Z\,\ab_z$
			& $\mu_x\,\ab_x+\mu_y\,\ab_y+\mu_z\,\ab_z$\\
			\bottomrule
		\end{tabular}%
	}
	\caption[Cenno ai vincoli spaziali]{%
		Vincoli spaziali fondamentali: molteplicità,
		gradi di libertà residui, forza reattiva e
		coppia reattiva. L'asse del vincolo è denotato
		con $\ab = \ab_z$.}
	\label{tab:vincoli_spaziali}
\end{table}

\begin{oss}
	Il piano di scorrimento e la biella hanno la stessa
	prestazione cinematica e statica nel regime
	infinitesimo: entrambi impongono
	$\ub(A)\cdot\ab = 0$ e reagiscono con una forza
	$R\,\ab$. La differenza è geometrica: nel piano di
	scorrimento la direzione $\ab$ è fissa nello spazio
	(normale al piano di scorrimento); nella biella la
	direzione $\ab$ coincide con l'asse dell'asta e
	ruota solidalmente con essa per spostamenti finiti.
	Nel regime infinitesimo, che è quello adottato in
	questo corso, i due vincoli sono equivalenti e
	possono essere usati indifferentemente come
	realizzazioni meccaniche l'uno dell'altro.
\end{oss}

\begin{oss}[Costruzione per giustapposizione]
	Anche i vincoli spaziali si possono interpretare come
	giustapposizioni di vincoli semplici, prendendo come
	mattoni elementari la \emph{biella} (molteplicità 1,
	una condizione di traslazione lungo una direzione
	assegnata) e il \emph{bipendolo} (molteplicità 1,
	una condizione di rotazione attorno a un asse
	assegnato). La logica è quella sviluppata nel
	\PARref{sec:vincoli_composti} per il caso piano: le
	equazioni cinematiche dei componenti si sommano e
	la molteplicità totale è la somma delle molteplicità.
	Per esempio:
	\begin{itemize}[noitemsep]
		\item la \textbf{cerniera sferica} (molt.~3) è
		la giustapposizione di tre bielle in $A$ con
		direzioni indipendenti, che bloccano le tre
		componenti di $\ub(A)$ lasciando libere le tre
		rotazioni;
		
		\item il \textbf{perno scorrevole} (molt.~4),
		con asse $\ab = \ab_z$, è la giustapposizione di
		due bielle in $A$ di direzioni $\ab_x, \ab_y$
		(che bloccano le traslazioni perpendicolari ad
		$\ab$) e di due bipendoli con assi $\ab_x, \ab_y$
		(che bloccano le rotazioni attorno agli assi
		perpendicolari ad $\ab$);
		
		\item il \textbf{glifo} (molt.~5), con asse $\ab = \ab_z$,
		si ottiene aggiungendo al perno scorrevole un terzo
		bipendolo con asse $\ab_z$, che blocca anche la rotazione
		attorno ad $\ab$;
		
		\item la \textbf{cerniera cilindrica} (molt.~5)
		si ottiene aggiungendo al perno scorrevole una
		terza biella di direzione $\ab_z$ (che blocca
		anche la traslazione lungo l'asse), oppure
		equivalentemente partendo da una cerniera sferica
		e aggiungendo due bipendoli con assi $\ab_x,
		\ab_y$;
		
		\item l'\textbf{incastro} (molt.~6) è la
		giustapposizione di una cerniera sferica e di
		tre bipendoli con assi indipendenti.
	\end{itemize}
	
	\smallskip
	
	\noindent
	Le decomposizioni proposte non sono uniche: ogni bipendolo,
	che blocca una rotazione del corpo, può essere sostituito
	da una biella applicata in un punto $B$ distinto, con
	direzione opportunamente scelta. Per esempio, la cerniera
	cilindrica di asse $\ab_z$ può essere realizzata
	equivalentemente come cerniera sferica in $A$ con due
	bipendoli di assi $\ab_x, \ab_y$, oppure come cerniera
	sferica in $A$ con due bielle in un secondo punto $B$ sul
	medesimo asse $\ab_z$, con direzioni $\ab_x$ e $\ab_y$:
	le bielle bloccano lo spostamento di $B$ nel piano
	perpendicolare ad $\ab_z$, equivalente al blocco delle
	rotazioni attorno ad $\ab_x$ e $\ab_y$. La libertà nella
	scelta della decomposizione riflette il fatto che la stessa
	prestazione cinematica può essere realizzata da
	configurazioni meccaniche diverse, ed è fonte di
	flessibilità nelle applicazioni.
	
	I due ``mattoni'' elementari del caso piano --- biella
	e bipendolo --- sono dunque sufficienti a costruire
	tutti i vincoli posizionali dello spazio, e la
	classificazione si estende in modo naturale dalla
	dimensione due alla dimensione tre.
\end{oss}




%-------------------------------------------------------------------------------
\section{Esercizi proposti}
%-------------------------------------------------------------------------------

\begin{esercizio}[Carrello]
	\label{ex:C04_carrello}
	
	Un corpo rigido piano è vincolato da un carrello con asse $\ab$ applicato nel punto $A$.
	(a) Scrivere l'equazione di vincolo~\eqref{eq:carrello_cin} e determinare il luogo geometrico dei centri di rotazione compatibili.
	(b) Scrivere la reazione vincolare $\rb(A)$ e identificare le incognite statiche.
	(c) Caso con cedimento $\bar{s}(A)$ nella direzione $\ab$, $\ub(A)\cdot\ab = \bar{s}(A)$: determinare il luogo dei centri di rotazione compatibili e confrontarlo con il caso $\bar{s}(A) = 0$.
	
	\risultato (a) Sostituendo $\ub(A) = \bm{\upvartheta} \times \overrightarrow{CA}$ in $\ub(A)\cdot\ab = 0$ si ottiene $(\bm{\upvartheta} \times \overrightarrow{CA})\cdot\ab = 0$, ovvero $\overrightarrow{CA}$ parallelo ad $\ab$. Il centro $C$ giace dunque sulla retta per $A$ parallela ad $\ab$. (b) Il PLV richiede $\rb(A)\cdot\ub(A) = 0$ per ogni $\ub(A)$ tale che $\ub(A)\cdot\ab = 0$, da cui $\rb(A) = R(A)\,\ab$, con $R(A) \in \mathbb{R}$ unica incognita scalare. (c) L'equazione $(\bm{\upvartheta} \times \overrightarrow{CA})\cdot\ab = \bar{s}(A)$ fornisce, sviluppando il prodotto misto nel piano, $\vartheta\,\overrightarrow{AC}\cdot\ab^{\perp} = \bar{s}(A)$. Il luogo dei centri compatibili è dunque la retta parallela ad $\ab$ traslata, rispetto al caso senza cedimento, di una distanza $\bar{s}(A)/\vartheta$ in direzione $\ab^{\perp}$. La struttura della reazione --- che dipende dalla sola prestazione cinematica --- non cambia.
	
\end{esercizio}

\begin{esercizio}[Cerniera]
	\label{ex:C04_cerniera}
	
	Un corpo rigido piano è vincolato da una cerniera in $A$.
	(a) Scrivere le equazioni di vincolo~\eqref{eq:cerniera_cin} e determinare il centro di rotazione $C$.
	(b) Scrivere la reazione $\rb(A)$ e identificare le incognite.
	(c) Mostrare che la cerniera può essere interpretata come la sovrapposizione di due carrelli ortogonali, verificando che i luoghi dei centri dei due carrelli si intersecano in $A$.
	
	\risultato (a) La cerniera impone $\ub(A) = \zerovb$, ossia $u(A) = 0$ e $v(A) = 0$. Il centro di rotazione coincide con il punto fisso, dunque $C \equiv A$. (b) Il PLV non vincola la rotazione (la cerniera consente $\vartheta$ arbitraria), quindi la coppia reattiva è nulla, $\mu = 0$; non vincola lo spostamento di $A$ (ammesso solo $\ub(A) = \zerovb$, su cui ogni reazione compie lavoro nullo), quindi $\rb(A)$ ha entrambe le componenti libere: $\rb(A) = X(A)\,\ab_x + Y(A)\,\ab_y$, con $X(A)$ e $Y(A) \in \mathbb{R}$ incognite scalari. (c) Considerando due carrelli con assi $\ab_x$ e $\ab_y$, entrambi applicati in $A$: il primo ha luogo dei centri sulla retta per $A$ parallela ad $\ab_x$; il secondo, sulla retta per $A$ parallela ad $\ab_y$. L'unico punto comune ai due luoghi è $A$ stesso. Le reazioni dei due carrelli ($R_1\,\ab_x$ e $R_2\,\ab_y$) si sommano nelle due componenti $X(A) = R_1$ e $Y(A) = R_2$ della reazione della cerniera, in piena coerenza con la molteplicità due.
	
\end{esercizio}

\begin{esercizio}[Glifo]
	\label{ex:C04_glifo}
	
	Un corpo rigido piano è vincolato da un glifo con asse $\ab$ applicato in $A$.
	(a) Scrivere le equazioni di vincolo~\eqref{eq:glifo_cin} e descrivere lo spostamento ammesso~\eqref{eq:glifo_moto_ammesso}.
	(b) Scrivere la forza reattiva $\rb(A)$ e la coppia reattiva $\mu(A)$, motivando perché il glifo introduce due incognite scalari.
	(c) Verificare che il lavoro della reazione sullo spostamento ammesso è identicamente nullo.
	(d) Caso con cedimenti assegnati $\ub(A)\cdot\ab = \bar{s}(A)$ e $\vartheta = \bar{\vartheta}$: descrivere lo spostamento totale del corpo e interpretare geometricamente $\bar{s}(A)$.
	
	\risultato (a) Equazioni del glifo: $\ub(A)\cdot\ab = 0$ e $\vartheta = 0$. Lo spostamento ammesso è $\ub(P) = \alpha\,\ab^{\perp}$ per ogni $P$ del corpo, con $\alpha \in \mathbb{R}$: traslazione rigida lungo $\ab^{\perp}$. Il centro di rotazione è all'infinito nella direzione $\ab$. (b) Per il PLV: $\rb(A)\cdot(\alpha\,\ab^{\perp}) + \mu(A)\cdot 0 = 0$ per ogni $\alpha$, da cui $\rb(A)\cdot\ab^{\perp} = 0$, ovvero $\rb(A) = R(A)\,\ab$. La coppia $\mu(A)$ non è vincolata dal PLV (essendo $\vartheta = 0$) e resta incognita libera. Le due incognite scalari $R(A)$ e $\mu(A)$ corrispondono ai due movimenti bloccati dal vincolo (traslazione lungo $\ab$ e rotazione), in accordo con la molteplicità due. (c) Il lavoro è $\rb(A)\cdot(\alpha\,\ab^{\perp}) + \mu(A)\cdot 0 = R(A)(\ab\cdot\ab^{\perp})\,\alpha = 0$, per ortogonalità di $\ab$ e $\ab^{\perp}$. (d) Lo spostamento totale del corpo si decompone in: una traslazione $\bar{s}(A)\,\ab$ del punto $A$ lungo l'asse bloccato (cedimento lineare imposto), una traslazione libera $\alpha\,\ab^{\perp}$ ammessa dal vincolo, e una rotazione rigida $\bar{\vartheta}$ attorno al centro di rotazione (ora finito). Il cedimento $\bar{s}(A)$ rappresenta lo scostamento di $A$ in direzione perpendicolare al piano di scorrimento del glifo --- cioè uno spostamento nella direzione che il vincolo, in condizioni ideali, impedirebbe. La struttura delle reazioni $\rb(A) = R(A)\,\ab$ e $\mu(A)$ resta invariata.
	
\end{esercizio}

\begin{esercizio}[Biella interna]
	\label{ex:C04_pendolo_interno}
	
	Due corpi $\mathcal{C}_i$ e $\mathcal{C}_j$ sono connessi da una biella interna con punti di applicazione $A_i$ e $A_j$ e versore $\ab \coloneqq \overrightarrow{A_iA_j}/|\overrightarrow{A_iA_j}|$.
	(a) Scrivere l'equazione di vincolo~\eqref{eq:carrello_int_punt}. Dati i centri assoluti $C_i$ e $C_j$, determinare il centro relativo $C_{ij}$ usando il teorema di allineamento del \CAPref{cap:cinematica_infinitesima}.
	(b) Scrivere le reazioni scambiate fra i due corpi in $A_i$ e $A_j$.
	(c) Caso con cedimento $\bar{s}(A)$, ovvero $\Delta\ub(A)\cdot\ab = \bar{s}(A)$ con $\Delta\ub(A) \coloneqq \ub(A_j) - \ub(A_i)$: interpretare fisicamente $\bar{s}(A)$ e determinare il luogo del centro relativo $C_{ij}$.
	
	\risultato (a) L'equazione di vincolo della biella interno è $\Delta\ub(A)\cdot\ab = 0$: il vincolo annulla la componente dello spostamento relativo nella direzione della biella. Per il teorema di allineamento, $C_i$, $C_j$ e $C_{ij}$ sono allineati. Per la cinematica della biella --- analoga a quella del carrello esterno --- $C_{ij}$ giace inoltre sulla retta passante per $A_i$ e $A_j$. Pertanto $C_{ij}$ è l'\textit{intersezione fra la retta $C_iC_j$ e la retta della biella} $A_iA_j$. (b) Le reazioni scambiate sono dirette lungo l'asse della biella, uguali in modulo e opposte in verso: $\rb_i(A_i) = -R\,\ab$ su $\mathcal{C}_i$ e $\rb_j(A_j) = +R\,\ab$ su $\mathcal{C}_j$, con $R \in \mathbb{R}$ unica incognita scalare. (c) Il cedimento $\bar{s}(A)$ rappresenta uno scorrimento relativo dei due corpi lungo l'asse della biella: la biella si comporta come se fosse stato allungata (o accorciata) di $\bar{s}(A)$. L'allineamento $C_i, C_j, C_{ij}$ persiste --- è una proprietà geometrica indipendente dal cedimento --- ma $C_{ij}$ non giace più sulla retta $A_iA_j$: si trasla in direzione perpendicolare ad essa di una distanza pari a $\bar{s}(A)/\vartheta_{ij}$, dove $\vartheta_{ij} = \vartheta_j - \vartheta_i$ è la rotazione relativa fra i due corpi.
	
\end{esercizio}

\begin{esercizio}[Identificazione del vincolo]
	\label{ex:C04_identificazione}
	
	Sono assegnati gli spostamenti di due punti di un corpo rigido piano:
	\begin{equation*}
		\ub(A) = \bar{u}(A)\,\ab_x, \qquad
		\ub(B) = \bar{v}(B)\,\ab_y.
	\end{equation*}
	(a) Ricavare la rotazione $\vartheta$ e lo spostamento $\ub(O)$ di un polo $O$ a piacere dalla formula fondamentale~\eqref{eq:FFCI_piano_vett} del \CAPref{cap:cinematica_infinitesima}.
	(b) Determinare il centro di rotazione $C$.
	(c) Identificare una combinazione di vincoli con cedimento compatibile con i dati assegnati.
	
	\risultato (a) Dalla formula fondamentale: $\ub(B) - \ub(A) = \bm{\upvartheta} \times \overrightarrow{AB}$. In componenti, con $\overrightarrow{AB} = (x_B - x_A,\,y_B - y_A)^\top$: la prima componente fornisce $-\bar{u}(A) = -\vartheta(y_B - y_A)$, da cui $\vartheta = \bar{u}(A)/(y_B - y_A)$; la seconda componente fornisce $\bar{v}(B) = \vartheta(x_B - x_A)$, da cui $\vartheta = \bar{v}(B)/(x_B - x_A)$. La compatibilità richiede $\bar{u}(A)(x_B - x_A) = \bar{v}(B)(y_B - y_A)$: se vale, $\vartheta$ è determinata; altrimenti i dati non sono compatibili con un moto rigido. Scegliendo $O = A$ come polo: $\ub(O) = \ub(A) = \bar{u}(A)\,\ab_x$. (b) Centro di rotazione: $\overrightarrow{AC} = (0,\,\bar{u}(A)/\vartheta)^\top = (0,\,y_B - y_A)^\top$, dunque $C = (x_A,\,y_B)$. Geometricamente, $C$ è il punto di coordinate $(x_A,\,y_B)$, ovvero l'intersezione della verticale per $A$ con l'orizzontale per $B$. (c) Una combinazione naturale: \textit{carrello in $A$ con asse $\ab_y$ e cedimento $\bar{u}(A)$}, e \textit{carrello in $B$ con asse $\ab_x$ e cedimento $\bar{v}(B)$}. Il primo carrello consente lo spostamento di $A$ in direzione $\ab_x$ (consistente con $\ub(A) = \bar{u}(A)\,\ab_x$); il secondo, lo spostamento di $B$ in direzione $\ab_y$ (consistente con $\ub(B) = \bar{v}(B)\,\ab_y$). I due carrelli formano un sistema isostatico: due equazioni cinematiche per tre gradi di libertà del corpo, con un'unica configurazione cinematica ammessa.
	
\end{esercizio}


